% L’évolution de l’artisanat — Géographie Tle D — Cours signé OnBuch+
% Programme officiel MINESEC. Compiler avec : tectonic N06-evolution-artisanat.tex
\newcommand{\DOCMATIERE}{Géographie}
\newcommand{\DOCNIVEAU}{Tle D}
\newcommand{\DOCMODULE}{Module 1 : Le Cameroun}
\newcommand{\DOCLECON}{N06}
\newcommand{\DOCTITRE}{L’évolution de l’artisanat}
% OnBuch+ — préambule « cours signés » ENRICHI (compilé avec tectonic / XeLaTeX)
% Système visuel « Pop » d'OnBuch+ : fond crème (jamais blanc), bordures encre
% épaisses, ombres « dures » décalées, titres Archivo Black en capitales.
% Version MATHÉMATIQUES : graphes (pgfplots/tikz), boîtes pédagogiques
% (définition, propriété, méthode, attention, le savais-tu, exercice résolu,
% exercice, TP, bilan), page de garde et signature OnBuch+ ancrée en bas.
%
% Macros attendues dans main.tex AVANT \input{preamble} :
%   \DOCMATIERE  \DOCNIVEAU  \DOCMODULE  \DOCLECON (numéro)  \DOCTITRE
\documentclass[11pt]{article}

\usepackage{fontspec}
\usepackage[french]{babel}
\usepackage[a4paper,margin=2cm,top=3.4cm,bottom=2.8cm,headheight=1.4cm,footskip=1.3cm]{geometry}
\usepackage{amsmath,amssymb,amsfonts}
\usepackage{mathtools} % \xmapsto, \coloneqq... souvent utilisés par les modèles
\usepackage{cancel} % simplifications barrées (bilans de forces)
% \abs{z} (module, valeur absolue) et \norm{u} (norme), utilisés par les modèles
\providecommand{\abs}{}\let\abs\relax\DeclarePairedDelimiter{\abs}{\lvert}{\rvert}
\providecommand{\norm}{}\let\norm\relax\DeclarePairedDelimiter{\norm}{\lVert}{\rVert}
\usepackage[table]{xcolor} % [table] : fournit aussi \rowcolors (alternance de lignes)
\usepackage{pagecolor}
\usepackage{tikz}
\usetikzlibrary{arrows.meta,positioning,calc,shapes.geometric,shapes.misc,shapes.arrows,decorations.pathmorphing,decorations.pathreplacing,decorations.markings,patterns,shadows,mindmap,trees,fit,backgrounds,matrix,intersections,angles,quotes}
\usepackage{pgfplots}
\usepackage[version=4]{mhchem} % équations nucléaires \ce{^{235}_{92}U}
\usepackage[european,siunitx]{circuitikz} % schémas électriques
\pgfplotsset{compat=1.18}
\usepgfplotslibrary{fillbetween,groupplots} % aires entre courbes (calcul intégral)
\usepackage{enumitem}
\usepackage{fancyhdr}
% [most] charge toutes les bibliothèques tcolorbox (listings, minted, raster,
% documentation...) et gonfle inutilement la mémoire TeX sur un document long
% (~300 boîtes) : on ne charge que ce qui est réellement utilisé.
\usepackage[skins,breakable]{tcolorbox}
\usepackage{graphicx}
\graphicspath{{/home/user/t-t/geo-onbuch/images/}}
\usepackage{colortbl}
\usepackage{booktabs}
\usepackage{tabularx}
\usepackage{multirow}
\usepackage{diagbox} % case d'en-tête diagonale des tableaux à double entrée
% Colonnes extensibles centrée / gauche / droite (C, L, R), souvent utilisées par les modèles
\newcolumntype{C}{>{\centering\arraybackslash}X}
\newcolumntype{L}{>{\raggedright\arraybackslash}X}
\newcolumntype{R}{>{\raggedleft\arraybackslash}X}
\usepackage{array}
\usepackage{multicol}
\usepackage{siunitx}
% parse-numbers=false : accepte aussi \SI{2,5 \times 10^{-3}}{...}
\sisetup{locale=FR,output-decimal-marker={,},group-minimum-digits=4,parse-numbers=false}
\DeclareSIUnit\ohm{\ensuremath{\symup{\Omega}}} % U+2126 absent de la police
\DeclareSIUnit\division{div} % réglages d'oscilloscope (ms/div, V/div)
\DeclareSIUnit\tr{tr} % tours (vitesse de rotation, ex. \SI{1500}{\tr\per\minute})
\DeclareSIUnit\molar{M} % concentration molaire (ex. \SI{3}{\molar}, \SI{1}{\micro\molar})
\DeclareSIUnit\an{an} % année (le modèle confond parfois avec \year, un compteur TeX) — ex. \SI{5}{\kilo\meter\per\an}
\DeclareSIUnit\FCFA{FCFA} % franc CFA (ex. \SI{500}{\FCFA\per\kilo\gram})
\DeclareSIUnit\degreeBrix{\textdegree Brix} % teneur en sucre d'un sirop/jus (réfractométrie)
\DeclareSIUnit\month{mois} % le modèle utilise parfois \month, un compteur TeX, comme unité de durée
\DeclareSIUnit\microliter{\micro\liter} % le modèle écrit parfois \microliter en un seul token
\DeclareSIUnit\million{millions} % dénombrement (ex. \SI{15}{\million\per\mL} spermatozoïdes)
\usepackage{qrcode}
% \degree : commande fréquemment utilisée par les modèles pour un angle ou une
% température en dehors de \SI{}{} (non fournie par les paquets chargés).
\providecommand{\degree}{^{\circ}}
% \qed (fin de démonstration), utilisé par les modèles sans amsthm
\providecommand{\qed}{\hfill$\square$}
% \barre : double barre des valeurs interdites dans un tableau de variations (tkz-tab)
\providecommand{\barre}{\ensuremath{\|}}
% environnement proof (amsthm non chargé) utilisé par les modèles
\ifdefined\proof\else\newenvironment{proof}[1][Démonstration]{\par\noindent\textit{#1.}\ }{\qed\par}\fi
\usepackage{lastpage}
\usepackage{hyperref}
\hypersetup{hidelinks}

% ============================================================
% Palette « Pop » OnBuch+ — mêmes hex que la classe P (plus_ui.dart)
% ============================================================
\definecolor{poporange}{HTML}{F59321}
\definecolor{popdark}{HTML}{E07A0C}
\definecolor{popcream}{HTML}{FFF6E8}
\definecolor{popink}{HTML}{1C1714}
\definecolor{poppurple}{HTML}{7A5AE0}
\definecolor{popgreen}{HTML}{2E9E6A}
\definecolor{poppink}{HTML}{E0457B}
\definecolor{popblue}{HTML}{2F7DE1}
\definecolor{popgold}{HTML}{C98A12}
\definecolor{popmuted}{HTML}{8A7F76}
\definecolor{popline}{HTML}{EADFD2}
\definecolor{popgray}{HTML}{8A7F76}
\colorlet{popgrey}{popgray}
% Alias tolérants (le modèle nomme parfois les couleurs différemment)
\colorlet{popwhite}{white}
\colorlet{popblack}{popink}
\colorlet{popred}{poppink}
\colorlet{popyellow}{popgold}
% Teintes claires (fonds de tableaux/encadrés) — le modèle nomme parfois
% les couleurs avec un suffixe L (ex. popblueL) sans les avoir définies.
\colorlet{poporangeL}{poporange!20}
\colorlet{popdarkL}{popdark!20}
\colorlet{poppurpleL}{poppurple!20}
\colorlet{popgreenL}{popgreen!20}
\colorlet{poppinkL}{poppink!20}
\colorlet{popblueL}{popblue!20}
\colorlet{popgoldL}{popgold!20}
\colorlet{popredL}{popred!20}
\colorlet{popgrayL}{popgray!20}
% Teintes claires pour fonds de tableaux / zones
\colorlet{popblueL}{popblue!10!white}
\colorlet{poporangeL}{poporange!14!white}
\colorlet{popgreenL}{popgreen!12!white}
\colorlet{poppurpleL}{poppurple!12!white}
\colorlet{poppinkL}{poppink!12!white}
\colorlet{popgoldL}{popgold!14!white}

\pagecolor{popcream}

% ============================================================
% Typographie
% ============================================================
\defaultfontfeatures{Ligatures=TeX}
\setmainfont{PlusJakartaSans-Medium.ttf}[Path=../../../fonts/,BoldFont=PlusJakartaSans-Bold.ttf]
% Maths en sans-serif (Fira Math) avec chiffres et lettres droites de Plus
% Jakarta, pour que les formules se fondent dans le texte.
\usepackage[math-style=ISO,bold-style=ISO]{unicode-math}
\setmathfont{FiraMath-Regular.otf}[Path=../../../fonts/]
\setmathfont{PlusJakartaSans-Medium.ttf}[Path=../../../fonts/,range={up/num,up/{Latin,latin},\mathup}]
\usepackage{icomma} % virgule décimale sans espace en mode math
% Caractères mathématiques Unicode absents des polices de texte (Plus Jakarta,
% Archivo Black) : rendus par leur équivalent mathématique, en texte comme en
% formule — sinon ils s'affichent en carré vide (ex. « Arithmétique dans ℤ »).
\usepackage{newunicodechar}
\newunicodechar{ℝ}{\ensuremath{\mathbb{R}}}
\newunicodechar{ℤ}{\ensuremath{\mathbb{Z}}}
\newunicodechar{ℂ}{\ensuremath{\mathbb{C}}}
\newunicodechar{ℕ}{\ensuremath{\mathbb{N}}}
\newunicodechar{ℚ}{\ensuremath{\mathbb{Q}}}
\newunicodechar{𝔻}{\ensuremath{\mathbb{D}}}
\newunicodechar{α}{\ensuremath{\alpha}}
\newunicodechar{β}{\ensuremath{\beta}}
\newunicodechar{γ}{\ensuremath{\gamma}}
\newunicodechar{δ}{\ensuremath{\delta}}
\newunicodechar{ε}{\ensuremath{\varepsilon}}
\newunicodechar{θ}{\ensuremath{\theta}}
\newunicodechar{λ}{\ensuremath{\lambda}}
\newunicodechar{μ}{\ensuremath{\mu}}
\newunicodechar{σ}{\ensuremath{\sigma}}
\newunicodechar{τ}{\ensuremath{\tau}}
\newunicodechar{φ}{\ensuremath{\varphi}}
\newunicodechar{ψ}{\ensuremath{\psi}}
\newunicodechar{ω}{\ensuremath{\omega}}
\newunicodechar{Σ}{\ensuremath{\Sigma}}
% majuscules produites par \MakeUppercase dans les titres (α-aminés -> Α)
\newunicodechar{Α}{\ensuremath{\alpha}}
\newunicodechar{Β}{\ensuremath{\beta}}
\newunicodechar{Γ}{\ensuremath{\Gamma}}
\newunicodechar{Δ}{\ensuremath{\Delta}}
\newunicodechar{Ε}{\ensuremath{\varepsilon}}
\newunicodechar{Θ}{\ensuremath{\Theta}}
\newunicodechar{Λ}{\ensuremath{\Lambda}}
\newunicodechar{Μ}{\ensuremath{\mu}}
\newunicodechar{Τ}{\ensuremath{\tau}}
\newunicodechar{Φ}{\ensuremath{\Phi}}
\newunicodechar{Ψ}{\ensuremath{\Psi}}
\newunicodechar{Ω}{\ensuremath{\Omega}}
\newunicodechar{↦}{\ensuremath{\mapsto}}
\newunicodechar{∈}{\ensuremath{\in}}
\newunicodechar{∉}{\ensuremath{\notin}}
\newunicodechar{∧}{\ensuremath{\wedge}}
\newunicodechar{⇔}{\ensuremath{\Leftrightarrow}}
\newunicodechar{⟺}{\ensuremath{\Longleftrightarrow}}
\newunicodechar{⇒}{\ensuremath{\Rightarrow}}
\newunicodechar{⟹}{\ensuremath{\Longrightarrow}}
\newunicodechar{∘}{\ensuremath{\circ}}
\newunicodechar{∪}{\ensuremath{\cup}}
\newunicodechar{∩}{\ensuremath{\cap}}
\newunicodechar{≡}{\ensuremath{\equiv}}
\newunicodechar{⊂}{\ensuremath{\subset}}
\newunicodechar{⊥}{\ensuremath{\perp}}
\newunicodechar{∃}{\ensuremath{\exists}}
\newunicodechar{∀}{\ensuremath{\forall}}
\newunicodechar{∛}{\ensuremath{\sqrt[3]{\,}}}
\newunicodechar{∜}{\ensuremath{\sqrt[4]{\,}}}
\newunicodechar{⃗}{\ensuremath{{}^{\to}}}
\newunicodechar{ⁿ}{\ifmmode^{n}\else\textsuperscript{\ensuremath{n}}\fi}
\newunicodechar{ˣ}{\ifmmode^{x}\else\textsuperscript{\ensuremath{x}}\fi}
\newunicodechar{ᵇ}{\ifmmode^{b}\else\textsuperscript{\ensuremath{b}}\fi}
\newunicodechar{ʳ}{\ifmmode^{r}\else\textsuperscript{\ensuremath{r}}\fi}
\newunicodechar{ᵘ}{\ifmmode^{u}\else\textsuperscript{\ensuremath{u}}\fi}
\newunicodechar{ʸ}{\ifmmode^{y}\else\textsuperscript{\ensuremath{y}}\fi}
\newunicodechar{ᵃ}{\ifmmode^{a}\else\textsuperscript{\ensuremath{a}}\fi}
\newunicodechar{ᵉ}{\ifmmode^{e}\else\textsuperscript{\ensuremath{e}}\fi}
\newunicodechar{ᵏ}{\ifmmode^{k}\else\textsuperscript{\ensuremath{k}}\fi}
\newunicodechar{ᵖ}{\ifmmode^{p}\else\textsuperscript{\ensuremath{p}}\fi}
\newunicodechar{ᴺ}{\ifmmode^{N}\else\textsuperscript{\ensuremath{N}}\fi}
\newunicodechar{ᶜ}{\ifmmode^{c}\else\textsuperscript{\ensuremath{c}}\fi}
\newunicodechar{ʰ}{\ifmmode^{h}\else\textsuperscript{\ensuremath{h}}\fi}
\newunicodechar{ᵠ}{\ifmmode^{q}\else\textsuperscript{\ensuremath{q}}\fi}
\newunicodechar{ₙ}{\ifmmode_{n}\else\textsubscript{\ensuremath{n}}\fi}
\newunicodechar{ₐ}{\ifmmode_{a}\else\textsubscript{\ensuremath{a}}\fi}
\newunicodechar{ᵢ}{\ifmmode_{i}\else\textsubscript{\ensuremath{i}}\fi}
\newunicodechar{ᵣ}{\ifmmode_{r}\else\textsubscript{\ensuremath{r}}\fi}
\newunicodechar{ₖ}{\ifmmode_{k}\else\textsubscript{\ensuremath{k}}\fi}
\newunicodechar{ᵦ}{\ifmmode_{\beta}\else\textsubscript{\ensuremath{\beta}}\fi}
\newunicodechar{ₓ}{\ifmmode_{x}\else\textsubscript{\ensuremath{x}}\fi}
\newfontfamily\popdisplay{ArchivoBlack-Regular.ttf}[Path=../../../fonts/]
\newfontfamily\poplogo{SpaceGrotesk-Bold.ttf}[Path=../../../fonts/]
\newfontfamily\popsemi{PlusJakartaSans-SemiBold.ttf}[Path=../../../fonts/]
\newfontfamily\popxbold{PlusJakartaSans-ExtraBold.ttf}[Path=../../../fonts/]
\newfontfamily\popxboldls{PlusJakartaSans-ExtraBold.ttf}[Path=../../../fonts/,LetterSpace=3.0]

\setlist[enumerate]{leftmargin=1.7em,itemsep=5pt,topsep=4pt}
\setlist[itemize]{leftmargin=1.7em,itemsep=4pt,topsep=4pt,label={\color{poporange}\textbullet}}
\setlength{\parindent}{0pt}
\setlength{\parskip}{5pt}
\renewcommand{\arraystretch}{1.25}

% pgfplots : style Pop par défaut
\pgfplotsset{
  every axis/.append style={axis line style={popink,line width=1pt},tick style={popink},
    grid style={popline,line width=0.5pt},label style={font=\small},tick label style={font=\footnotesize}},
  popcurve/.style={poporange,line width=1.8pt,no markers,smooth},
}
% Même style disponible dans un \draw TikZ simple (hors pgfplots)
\tikzset{popcurve/.style={poporange,line width=1.8pt}}
\tikzset{popgrid/.style={popline,very thin}}
\colorlet{popcurve}{poporange} % le nom du style est parfois utilisé comme couleur

% ============================================================
% Pill / squiggle
% ============================================================
\newcommand{\poppill}[3]{%
  \tikz[baseline=-0.55ex]\node[draw=popink,line width=1.2pt,rounded corners=2.6pt,%
    fill=#2,inner xsep=6.5pt,inner ysep=3pt]{%
    \popxboldls\fontsize{7}{7}\selectfont\color{#3}\MakeUppercase{#1}};%
}
\newcommand{\popsquiggle}[1]{%
  \tikz[baseline=-0.4ex]{
    \draw[#1,line width=2.6pt,line cap=round]
      plot [smooth] coordinates {(0,0.09) (0.34,0.01) (0.68,0.21) (1.02,0.01) (1.36,0.10)};
  }%
}
% Mot-clé surligné (marqueur orange sous le texte)
\newcommand{\cle}[1]{{\popxbold\color{popdark}#1}}

% ============================================================
% En-tête / pied
% ============================================================
\pagestyle{fancy}
\fancyhf{}
\renewcommand{\headrulewidth}{0pt}
\renewcommand{\footrulewidth}{0pt}
\lhead{\raisebox{-0.15ex}{\poplogo\fontsize{13}{13}\selectfont\color{popink}OnBuch\color{poporange}+}}
\rhead{\poppill{\DOCMATIERE\ \textbullet\ \DOCNIVEAU\ \textbullet\ Leçon \DOCLECON}{white}{popink}}
\lfoot{\popxbold\fontsize{7.6}{7.6}\selectfont\color{popmuted}\MakeUppercase{Cours signé OnBuch+}\ \textbullet\ {\popsemi\DOCTITRE}}
\rfoot{\tikz[baseline=-0.6ex]\node[rounded corners=8.5pt,fill=popink,minimum height=17pt,minimum width=17pt,inner xsep=5pt,inner sep=0]{\popxbold\fontsize{8}{8}\selectfont\color{popcream}\thepage\,/\,\pageref*{LastPage}};}

% ============================================================
% Titres
% ============================================================
\newcommand{\chaptitre}[2]{%
  \begin{tcolorbox}[enhanced,colback=poporange,colframe=popink,boxrule=2.6pt,arc=11pt,
      drop shadow=popink,left=15pt,right=15pt,top=12pt,bottom=11pt,before skip=3pt,after skip=18pt]
    \poppill{#2}{popink}{popcream}\par\vspace{9pt}
    {\raggedright\hyphenpenalty=10000\exhyphenpenalty=10000\popdisplay\fontsize{22}{24}\selectfont\color{popcream}\MakeUppercase{#1}\par}
  \end{tcolorbox}
}
% Section numérotée : pastille orange avec le numéro + titre Archivo
\newcounter{popsec}
\newcommand{\coursec}[1]{%
  \stepcounter{popsec}\setcounter{popsub}{0}%
  \par\vspace{16pt}\noindent
  \begin{minipage}{\linewidth}
  \tikz[baseline=-0.7ex]\node[draw=popink,line width=1.6pt,fill=poporange,rounded corners=4pt,minimum size=22pt,inner sep=0]{\popdisplay\fontsize{12}{12}\selectfont\color{popcream}\thepopsec};%
  \hspace{8pt}{\popdisplay\fontsize{15}{17}\selectfont\color{popink}\MakeUppercase{#1}}\par
  \vspace{4pt}\noindent{\color{poporange}\rule{36pt}{3.6pt}}
  \end{minipage}\par\nopagebreak\vspace{8pt}
}
\newcounter{popsub}
\newcommand{\courssub}[1]{%
  \stepcounter{popsub}%
  \par\vspace{9pt}\noindent{\popxbold\fontsize{11.5}{13}\selectfont\color{popink}{\color{poporange}\thepopsec.\thepopsub}\hspace{6pt}#1}\par\nopagebreak\vspace{4pt}
}

% ============================================================
% Boîtes pédagogiques — même recette : fond blanc, bordure épaisse colorée,
% ombre dure encre, badge plein. Titre optionnel [#1].
% ============================================================
\tcbset{popbase/.style={breakable,enhanced,colback=white,boxrule=2.3pt,arc=9pt,
  shadow={3pt}{-3pt}{0pt}{popink},left=14pt,right=14pt,top=11pt,bottom=11pt,before skip=11pt,after skip=15pt,
  fontupper=\popsemi\fontsize{10.6}{14.5}\selectfont\color{popink},
  fontlower=\popsemi\fontsize{10.6}{14.5}\selectfont\color{popink}}}
% Coche dessinée (le glyphe ✓ n'existe pas dans les polices du document)
\newcommand{\popcheck}[1]{\tikz[baseline=-0.5ex]\draw[#1,line width=1.6pt,line cap=round,line join=round](-0.13,0.0)--(-0.04,-0.09)--(0.13,0.1);}
\providecommand{\coche}{\popcheck{popgreen}}
\providecommand{\resultat}[1]{\par\smallskip\noindent\fcolorbox{popgreen}{popgreenL}{\parbox{\dimexpr\linewidth-2\fboxsep-2\fboxrule\relax}{\textbf{Résultat.} #1}}\par\smallskip}
\newcommand{\popboxhead}[3]{\poppill{#1}{#2}{white}\ifx\relax#3\relax\else\hspace{6pt}{\popxbold\fontsize{10.6}{12}\selectfont\color{popink}#3}\fi\par\vspace{7pt}}

\newtcolorbox{aretenir}[1][]{popbase,colframe=popgreen,before upper={\popboxhead{À retenir}{popgreen}{#1}}}
\newtcolorbox{exemplebox}[1][]{popbase,colframe=poporange,before upper={\popboxhead{Exemple}{poporange}{#1}}}
\newtcolorbox{definition}[1][]{popbase,colframe=popblue,before upper={\popboxhead{Définition}{popblue}{#1}}}
\newtcolorbox{propriete}[1][]{popbase,colframe=poppurple,before upper={\popboxhead{Propriété}{poppurple}{#1}}}
\newtcolorbox{methode}[1][]{popbase,colframe=popgold,before upper={\popboxhead{Méthode}{popgold}{#1}}}
\newtcolorbox{attention}[1][]{popbase,colframe=poppink,before upper={\popboxhead{Attention}{poppink}{#1}}}
\newtcolorbox{experience}[1][]{popbase,colframe=popblue,colback=popblueL,before upper={\popboxhead{Expérience}{popblue}{#1}}}
% « Le savais-tu ? » : carte violette pleine (contexte, culture, Cameroun)
\newtcolorbox{savaistu}[1][]{popbase,colback=poppurple,colframe=popink,
  fontupper=\popsemi\fontsize{10.4}{14.2}\selectfont\color{white},
  before upper={\poppill{Le savais-tu\,?}{white}{poppurple}\ifx\relax#1\relax\else\hspace{6pt}{\popxbold\color{white}#1}\fi\par\vspace{7pt}}}
% Exercice résolu : énoncé en haut, \tcblower puis la solution sur fond crème
\newcounter{popexo}\newcounter{popexr}
\newtcolorbox{exoresolu}[1][]{popbase,colframe=poporange,colbacklower=popcream,
  segmentation style={popink,line width=1.2pt,dashed},
  before upper={\stepcounter{popexr}\popboxhead{Exercice résolu \thepopexr}{poporange}{#1}},
  before lower={\poppill{Solution}{popink}{popcream}\par\vspace{6pt}}}
% Exercice (non résolu en place) — la correction est dans la partie Corrigés
\newtcolorbox{exercice}[1][]{popbase,colframe=popink,
  before upper={\stepcounter{popexo}\popboxhead{Exercice \thepopexo}{popink}{#1}}}
% Corrigé
\newtcolorbox{corrige}[1][]{popbase,colframe=popgreen,colback=popgreenL,
  before upper={\popboxhead{Corrigé}{popgreen}{#1}}}
% Objectifs / prérequis / bilan
\newtcolorbox{objectifs}{popbase,colframe=popink,colback=poporangeL,
  before upper={\popboxhead{Objectifs de la leçon}{popink}{}}}
\newtcolorbox{prerequis}{popbase,colframe=popink,colback=popblueL,
  before upper={\popboxhead{Prérequis}{popblue}{}}}
\newtcolorbox{bilan}{popbase,colframe=popink,colback=white,boxrule=2.8pt,
  before upper={\popboxhead{Fiche bilan}{poporange}{L'essentiel à connaître pour l'examen}}}
% Figure encadrée avec légende
\newtcolorbox{popfigure}[1][]{enhanced,breakable=false,colback=white,colframe=popline,boxrule=1.4pt,arc=8pt,
  left=10pt,right=10pt,top=10pt,bottom=8pt,before skip=10pt,after skip=12pt,halign=center,
  lower separated=false,halign lower=center,
  fontlower=\popsemi\fontsize{9}{11.5}\selectfont\color{popmuted},
  #1}
\newcommand{\legende}[1]{\par\vspace{6pt}{\popsemi\fontsize{9}{11.5}\selectfont\color{popmuted}#1}}

% ============================================================
% Signature OnBuch+
% ============================================================
\newcommand{\poplogoinline}[1]{{\poplogo\fontsize{#1}{#1}\selectfont\color{popink}OnBuch\color{poporange}+}}
\newcommand{\signatureonbuch}{%
  \begin{tcolorbox}[enhanced,colback=popink,colframe=popink,boxrule=2.2pt,arc=10pt,
      drop shadow=poporange,left=14pt,right=14pt,top=10pt,bottom=10pt,before skip=0pt,after skip=0pt]
    \tikz[baseline=-0.7ex]\node[circle,fill=popgreen,draw=popcream,line width=1.3pt,minimum size=22pt,inner sep=0]{\popcheck{white}};%
    \hspace{9pt}\begin{minipage}[c]{0.62\linewidth}
      {\popxbold\fontsize{12}{13}\selectfont\color{popcream}Signé\ \textbullet\ {\poplogo OnBuch\color{poporange}+}}\par\vspace{2pt}
      {\raggedright\popsemi\fontsize{8}{10.5}\selectfont\color{popline}\MakeUppercase{Équipe pédagogique OnBuch+}\\\MakeUppercase{Conforme au programme officiel MINESEC}\par}
    \end{minipage}\hfill
    {\poplogo\fontsize{22}{22}\selectfont\color{popcream}OnBuch\color{poporange}+}
  \end{tcolorbox}%
}

% ============================================================
% Page de garde
% #1 titre · #2 matière · #3 niveau · #4 module · #5 n° leçon · #6 durée ·
% #7 description / objectifs (texte ou itemize)
% ============================================================
\newcommand{\pagegarde}[7]{%
  \thispagestyle{empty}
  \newgeometry{a4paper,margin=1.7cm,top=1.6cm,bottom=1.4cm}%
  \begin{tikzpicture}[remember picture,overlay]
    \fill[poporange!18] ([xshift=-3.2cm,yshift=-3.6cm]current page.north east) circle (4.6cm);
    \fill[poppurple!14] ([xshift=2.4cm,yshift=7.6cm]current page.south west) circle (3.8cm);
  \end{tikzpicture}%
  {\poplogo\fontsize{30}{30}\selectfont\color{popink}OnBuch\color{poporange}+}\hfill
  \raisebox{0.6ex}{\poppill{Cours signé \textbullet\ Premium}{popink}{popcream}}\par
  \vspace{1pt}\popsquiggle{poppurple}\par
  \vspace{8pt}
  \begin{tcolorbox}[enhanced,colback=poporange,colframe=popink,boxrule=2.8pt,arc=13pt,
      drop shadow=popink,left=18pt,right=18pt,top=12pt,bottom=13pt,after skip=10pt]
    \poppill{#2 \textbullet\ #3}{popink}{popcream}\hspace{5pt}\poppill{#4}{white}{popink}\par\vspace{8pt}
    {\popdisplay\fontsize{14}{14}\selectfont\color{popink}LEÇON #5}\par\vspace{4pt}
    {\raggedright\hyphenpenalty=10000\exhyphenpenalty=10000\popdisplay\fontsize{25}{27}\selectfont\color{popcream}\MakeUppercase{#1}\par}
    \vspace{8pt}
    {\popxbold\fontsize{9.5}{9.5}\selectfont\color{popink}\MakeUppercase{Durée conseillée : #6}}
  \end{tcolorbox}
  \begin{tcolorbox}[enhanced,colback=white,colframe=popink,boxrule=2.2pt,arc=10pt,
      drop shadow=popink,left=16pt,right=16pt,top=10pt,bottom=10pt,after skip=8pt]
    \poppill{Dans cette leçon}{poppurple}{white}\par\vspace{5pt}
    \popsemi\fontsize{9.3}{12.6}\selectfont\color{popink}#7
  \end{tcolorbox}
  \noindent\begin{minipage}[t]{0.487\linewidth}
    \begin{tcolorbox}[enhanced,colback=popgreen,colframe=popink,boxrule=2.2pt,arc=10pt,
        drop shadow=popink,left=11pt,right=11pt,top=10pt,bottom=10pt,height=3.1cm,valign=center]
      \tcbox[colback=white,colframe=popink,boxrule=1.3pt,arc=5pt,boxsep=0pt,nobeforeafter,
        left=3pt,right=3pt,top=3pt,bottom=3pt,valign=center]{\qrcode[height=1.9cm]{https://play.google.com/store/apps/details?id=com.luvvix.onbuch&hl=fr}}%
      \hspace{8pt}\begin{minipage}[c]{0.5\linewidth}
        {\popdisplay\fontsize{11}{12.5}\selectfont\color{white}\MakeUppercase{Télécharge OnBuch}}\par\vspace{4pt}
        {\raggedright\popsemi\fontsize{8.4}{11}\selectfont\color{white}Gratuit \textbullet\ Google Play\\Tuteur IA, annales, concours, résultats\par}
      \end{minipage}
    \end{tcolorbox}
  \end{minipage}\hfill
  \begin{minipage}[t]{0.487\linewidth}
    \begin{tcolorbox}[enhanced,colback=poppurple,colframe=popink,boxrule=2.2pt,arc=10pt,
        drop shadow=popink,left=11pt,right=11pt,top=10pt,bottom=10pt,height=3.1cm,valign=center]
      \tikz[baseline=-0.7ex]\node[circle,fill=white,draw=popink,line width=1.3pt,minimum size=1.6cm,inner sep=0]{\popdisplay\fontsize{24}{24}\selectfont\color{poppurple}+};%
      \hspace{8pt}\begin{minipage}[c]{0.6\linewidth}
        {\popdisplay\fontsize{11}{12.5}\selectfont\color{white}\MakeUppercase{Passe à OnBuch+}}\par\vspace{4pt}
        {\raggedright\popsemi\fontsize{8.4}{11}\selectfont\color{white}Tous les cours signés, Tuteur IA illimité, fiches PDF — onglet OnBuch+ de l'app\par}
      \end{minipage}
    \end{tcolorbox}
  \end{minipage}
  \vspace{8pt}
  \popcontenu
  \vfill
  \signatureonbuch
  \restoregeometry
  \newpage
}

% Rangée « Ce cours contient » (page de garde)
\newcommand{\poptile}[3]{% #1 couleur #2 gros texte #3 légende
  \begin{tcolorbox}[enhanced,colback=#1,colframe=popink,boxrule=2pt,arc=9pt,shadow={2.5pt}{-2.5pt}{0pt}{popink},
      left=6pt,right=6pt,top=9pt,bottom=9pt,halign=center,valign=center,height=1.75cm,width=\linewidth,nobeforeafter]
    {\popdisplay\fontsize{12.5}{13}\selectfont\color{white}\MakeUppercase{#2}}\par\vspace{3pt}
    {\centering\popsemi\fontsize{7.8}{9.5}\selectfont\color{white}#3\par}
  \end{tcolorbox}}
\newcommand{\popcontenu}{%
  \par\noindent{\popxboldls\fontsize{7.5}{7.5}\selectfont\color{popmuted}CE COURS SIGNÉ CONTIENT}\par\vspace{7pt}
  \noindent\begin{minipage}[t]{0.235\linewidth}\poptile{popblue}{Cours}{complet, illustré, pas à pas}\end{minipage}\hfill
  \begin{minipage}[t]{0.235\linewidth}\poptile{poporange}{Méthodes}{et exercices résolus}\end{minipage}\hfill
  \begin{minipage}[t]{0.235\linewidth}\poptile{poppink}{Exercices}{type examen + corrigés}\end{minipage}\hfill
  \begin{minipage}[t]{0.235\linewidth}\poptile{popgreen}{Fiche bilan}{l'essentiel à retenir}\end{minipage}\par}

% Sommaire
\newtcolorbox{sommaire}{enhanced,colback=white,colframe=popink,boxrule=2.2pt,arc=10pt,
  drop shadow=popink,left=18pt,right=18pt,top=13pt,bottom=13pt,before skip=6pt,after skip=20pt,
  before upper={\poppill{Sommaire}{poppurple}{white}\par\vspace{9pt}\popsemi\fontsize{10.6}{14.5}\selectfont\color{popink}}}

% Signature de fin de cours — poussée tout en bas de la dernière page
\newcommand{\findecours}{%
  \par\vfill\nopagebreak
  \begin{center}\popsquiggle{poporange}\end{center}\vspace{4pt}
  \signatureonbuch
}
\AtBeginDocument{\def\llbracket{\mathopen{[\mkern-3.5mu[}}\def\rrbracket{\mathclose{]\mkern-3.5mu]}}} % intervalles d'entiers (glyphe ⟦ absent de la police)
% Glyphes absents de Fira Math : redéfinis à partir de glyphes disponibles
\AtBeginDocument{%
  \def\setminus{\mathbin{\backslash}}\let\smallsetminus\setminus
  \def\vdots{\mathord{\vbox{\baselineskip3.2pt\lineskiplimit0pt\kern4pt\hbox{.}\hbox{.}\hbox{.}}}}%
  \def\ddots{\mathinner{\mkern1mu\raise6.5pt\hbox{.}\mkern2mu\raise3.6pt\hbox{.}\mkern2mu\raise0.7pt\hbox{.}\mkern1mu}}%
  \def\triangle{\mathord{\scalebox{0.9}{$\symup{\Delta}$}}}%
  \def\nabla{\mathord{\raisebox{\depth}{\scalebox{1}[-1]{$\symup{\Delta}$}}}}%
  \def\checkmark{\popcheck{popdark}}%
  \def\star{\mathbin{\ast}}\def\bigstar{\mathord{\ast}}%
  \def\diamond{\mathbin{\scalebox{0.6}{$\Diamond$}}}%
  \def\bigcup{\mathop{\mathchoice{\raisebox{-0.3em}{\scalebox{1.7}{$\cup$}}}{\scalebox{1.25}{$\cup$}}{\cup}{\cup}}\displaylimits}%
  \def\bigcap{\mathop{\mathchoice{\raisebox{-0.3em}{\scalebox{1.7}{$\cap$}}}{\scalebox{1.25}{$\cap$}}{\cap}{\cap}}\displaylimits}%
  \def\lesseqqgtr{\mathrel{\lessgtr}}\def\bigoplus{\mathop{\oplus}\displaylimits}%
}
\providecommand{\ssi}{si et seulement si} % abréviation fréquente
\AtBeginDocument{\providecommand{\wideparen}{\overparen}} % arc de cercle

%% FIGFIX-BEGIN v5 — correctifs globaux des figures (docs/onbuch-generation/tools/figfix)
\usepackage{adjustbox}\usepackage{etoolbox}\tcbuselibrary{raster}
% toute figure TikZ/pgfplots plus large que la ligne est réduite à la largeur disponible
\BeforeBeginEnvironment{tikzpicture}{\begin{adjustbox}{max width=\linewidth}}
\AfterEndEnvironment{tikzpicture}{\end{adjustbox}}
% légendes pgfplots lisibles par-dessus les courbes
\pgfplotsset{every axis/.append style={legend style={fill=white,fill opacity=0.92,text opacity=1,draw=black!15,inner sep=3pt}}}
% étiquettes d'arêtes (syntaxe edge["..."]) avec fond blanc
\tikzset{every edge quotes/.append style={fill=white,inner sep=1.2pt}}
%% FIGFIX-END
\begin{document}

\pagegarde{L’évolution de l’artisanat}{Géographie}{Tle D}{Module 1 : Le Cameroun}{N06}{2 h}{Cette leçon étudie les transformations profondes de l'artisanat camerounais : modernisation des équipements, apparition de produits nouveaux, amélioration de la qualité et respect des normes, organisation des filières et des circuits de vente. Elle montre comment l'artisanat, secteur traditionnel, devient un levier d'intégration économique nationale.
\vspace{4pt}
\begin{multicols}{2}\small
\begin{itemize}
\item À la fin de cette leçon, je sais définir l'artisanat, la qualité et la norme
\item À la fin de cette leçon, je sais observer et décrire, à partir de photographies, les équipements modernes utilisés par les artisans camerounais
\item À la fin de cette leçon, je sais inventorier et classer les produits artisanaux nouveaux par domaine (décoration, habillement, agroalimentaire)
\item À la fin de cette leçon, je sais expliquer comment le conditionnement et le contrôle qualité améliorent la compétitivité des produits artisanaux
\item À la fin de cette leçon, je sais décrire l'organisation d'une filière artisanale et ses circuits de vente
\item À la fin de cette leçon, je sais localiser sur une carte du Cameroun les principaux pôles artisanaux et légender cette carte
\end{itemize}\end{multicols}}

\begin{sommaire}
\begin{multicols}{2}
\begin{enumerate}
\item L'artisanat camerounais : définitions et héritage traditionnel
\item La modernisation des équipements de travail
\item Des produits nouveaux pour des marchés nouveaux
\item L'amélioration de la qualité et le respect des normes
\item L'organisation des filières et des circuits de vente
\item Travaux pratiques : cartographier et schématiser l'artisanat camerounais
\item Activité d'intégration
\item Méthodes et exercices résolus
\item Exercices
\item Corrigés des exercices
\item Fiche bilan
\end{enumerate}
\end{multicols}
\end{sommaire}

\begin{objectifs}
\begin{itemize}
\item À la fin de cette leçon, je sais définir l'artisanat, la qualité et la norme
\item À la fin de cette leçon, je sais observer et décrire, à partir de photographies, les équipements modernes utilisés par les artisans camerounais
\item À la fin de cette leçon, je sais inventorier et classer les produits artisanaux nouveaux par domaine (décoration, habillement, agroalimentaire)
\item À la fin de cette leçon, je sais expliquer comment le conditionnement et le contrôle qualité améliorent la compétitivité des produits artisanaux
\item À la fin de cette leçon, je sais décrire l'organisation d'une filière artisanale et ses circuits de vente
\item À la fin de cette leçon, je sais localiser sur une carte du Cameroun les principaux pôles artisanaux et légender cette carte
\item À la fin de cette leçon, je sais construire un croquis ou un diagramme présentant une filière artisanale
\item À la fin de cette leçon, je sais interpréter des tableaux statistiques et des documents sur la production artisanale camerounaise
\end{itemize}
\end{objectifs}

\begin{prerequis}
\begin{itemize}
\item Notion d'artisanat et place du secteur informel dans l'économie camerounaise (classe de Première)
\item Les activités économiques du Cameroun : agriculture, industrie, services (début d'année de Terminale)
\item Lecture et légendage d'une carte du Cameroun (régions, villes principales)
\item Notion de filière et de circuit de commercialisation
\item Construction de diagrammes simples (histogrammes, diagrammes circulaires)
\end{itemize}
\end{prerequis}

\newpage

\coursec{L'artisanat camerounais : définitions et héritage traditionnel}

\courssub{Situation d'accroche et problématique}

À Foumban, dans la région de l'Ouest, un sculpteur qui taillait autrefois ses masques à la seule machette utilise désormais une scie électrique et une ponceuse : sa production a fortement augmenté et ses œuvres partent jusqu'à Douala, voire jusqu'à Paris. À Bafoussam, une coopérative de femmes conditionne les feuilles de ndolé et le piment dans des emballages étiquetés conformes aux normes, vendus jusque dans les rayons des supermarchés. Pourtant, à quelques kilomètres de là, d'autres artisans peinent à vendre leurs produits, faute d'équipements modernes et de circuits de commercialisation organisés.

Ces trois scènes, observables chaque jour au Cameroun, posent la question centrale de cette leçon : \textbf{comment l'artisanat camerounais s'est-il transformé ? Quels facteurs expliquent la réussite de certains artisans et la marginalisation d'autres ?} Pour y répondre, il faut d'abord comprendre ce qu'est l'artisanat, d'où il vient, et quel poids il occupe dans la vie économique et sociale du pays. C'est l'objet de cette première partie : poser les définitions et mesurer l'héritage traditionnel sur lequel s'est construite l'évolution actuelle.

\courssub{Qu'est-ce que l'artisanat ?}

\textbf{L'intuition d'abord.} Imagine le cordonnier installé sous un hangar au marché de Mokolo à Yaoundé : il répare et fabrique des chaussures avec ses mains, quelques outils (marteau, alêne, couteau, parfois une machine à coudre à pédale), seul ou aidé d'un ou deux apprentis. Il achète le cuir, découpe, assemble, colle, polit et vend lui-même. Il maîtrise \emph{toutes} les étapes du produit, de la matière première à la vente. Voilà l'artisanat dans sa forme la plus pure.

\begin{definition}[Artisanat]
L'\cle{artisanat} est une activité de transformation et de production réalisée \textbf{manuellement} ou avec des \textbf{outils simples}, par un \textbf{artisan} seul ou avec peu de main-d'œuvre (familiale ou apprentis). Il se caractérise par :
\begin{itemize}
\item une production à \textbf{petite échelle}, en petites séries ou à l'unité ;
\item une \textbf{faible division du travail} : l'artisan maîtrise le produit de bout en bout, de la matière première à la vente ;
\item un atelier souvent \textbf{familial}, installé au domicile, au marché ou en bordure de route ;
\item un savoir-faire transmis par \textbf{apprentissage}, le plus souvent de maître à élève, parfois de parents à enfants.
\end{itemize}
\end{definition}

\begin{attention}[Ne pas confondre artisanat et industrie]
L'erreur la plus fréquente au Baccalauréat consiste à utiliser « artisanat » et « industrie » comme des synonymes. Or l'\textbf{industrie} repose sur la mécanisation poussée, la production de masse et une forte division du travail (chaque ouvrier n'effectue qu'une tâche parcellaire sur une chaîne). L'\textbf{artisan}, lui, conçoit, fabrique et vend souvent lui-même son produit. Un sculpteur de Foumban n'est pas un ouvrier d'usine : c'est un créateur-producteur-vendeur. Dans une copie, précise toujours ce critère de \cle{maîtrise du produit de bout en bout}.
\end{attention}

\courssub{Les domaines traditionnels de l'artisanat camerounais}

L'artisanat camerounais est d'une richesse exceptionnelle, héritée de siècles de savoir-faire propres à chaque aire culturelle. On peut classer les domaines traditionnels en grandes familles :

\begin{itemize}
\item la \textbf{sculpture sur bois} : masques, statuettes, trônes, tam-tams, portes sculptées (Ouest et Nord-Ouest surtout) ;
\item la \textbf{poterie} : jarres, marmites, canaris, objets décoratifs (Nord, Adamaoua, Ouest) ;
\item la \textbf{vannerie} : paniers, nattes, chapeaux, meubles en rotin et en raphia (Littoral, Sud, Ouest) ;
\item le \textbf{tissage et la broderie} : pagnes, togues, ndop, gandouras brodées (Nord-Ouest, Ouest, Nord) ;
\item la \textbf{forge} : houes, machettes, couteaux, pointes de flèches et de lances (toutes les régions, jadis vitale pour l'agriculture et la chasse) ;
\item la \textbf{bijouterie et le travail du bronze} : colliers, bracelets, pipes et figurines en bronze coulé à la cire perdue (Foumban notamment) ;
\item la \textbf{transformation alimentaire domestique} : séchage du poisson et du manioc, fabrication de l'huile de palme, du beurre de karité, des feuilles de ndolé conditionnées, des épices moulues.
\end{itemize}

Chaque région s'est ainsi spécialisée en fonction de ses matières premières (bois, argile, raphia, cuir, métal) et de ses traditions culturelles. Cette spécialisation géographique constitue la base des \cle{pôles artisanaux} que nous allons localiser.

\courssub{Les grands pôles artisanaux traditionnels du Cameroun}

L'observation d'une carte des activités artisanales révèle que l'artisanat n'est pas réparti au hasard : il se concentre dans des \textbf{pôles} où se combinent matières premières, traditions culturelles fortes et présence d'une clientèle (cours royales, marchés, axes commerciaux).

\begin{popfigure}
\begin{center}
\includegraphics[width=\linewidth,height=9.5cm,keepaspectratio]{N01/cmr_map.jpg}
\end{center}
\legende{Fond de carte : carte générale du Cameroun. Les pôles artisanaux traditionnels s'y repèrent à partir des villes : Foumban et Bafoussam (Ouest), Bamenda (Nord-Ouest), Maroua (Extrême-Nord), Douala et Yaoundé (grandes villes du Sud). Leurs spécialités (sculpture, bronze, poterie, tissage, cuir) sont décrites dans le texte ci-dessous. \\ Source : Wikimedia Commons, CIA, domaine public.}
\end{popfigure}

Décrivons brièvement ces pôles :

\begin{itemize}
\item \textbf{Foumban} (région de l'Ouest) : capitale historique du royaume bamoun, c'est le pôle artisanal le plus célèbre du pays. Le « quartier des artisans » y regroupe sculpteurs sur bois, bronziers (technique de la fonte à la cire perdue), brodeurs et tisserands. Le musée du Palais royal conserve des siècles de création artistique.
\item \textbf{Foumbot}, voisine de Foumban, est réputée pour sa poterie et ses bronzes.
\item \textbf{Bafoussam} (Ouest) : capitale de la broderie, de la vannerie en raphia et du tissage.
\item \textbf{Bamenda} (Nord-Ouest) : réputée pour le tissage du célèbre pagne \emph{ndop} et la vannerie.
\item \textbf{Maroua} (Extrême-Nord) : pôle du travail du \textbf{cuir} (sacs, sandales, poufs) et de la \textbf{poterie}, héritier des traditions peules et des échanges transsahéliens.
\item \textbf{Kribi et les zones côtières} (Sud et Littoral) : vannerie, nattes et objets en raphia et en rotin, matières abondantes en zone forestière humide.
\item \textbf{Batouri} (Est) : sculpture et vannerie, dans une région forestière riche en bois précieux.
\end{itemize}

\begin{savaistu}[Foumban, la « cité des arts »]
Foumban, capitale du royaume bamoun, est surnommée la \textbf{« cité des arts »}. Son quartier des artisans regroupe sculpteurs, bronziers et brodeurs depuis des générations, sous le patronage historique des sultans. Le roi Njoya (fin du \textsc{xix}\textsuperscript{e} -- début du \textsc{xx}\textsuperscript{e} siècle), inventeur d'une écriture (le \emph{shu-mom}), y encouragea toutes les formes de création. Aujourd'hui encore, la Route des arts et le marché artisanal de Foumban attirent collectionneurs et touristes venus de Douala, d'Europe et d'Amérique.
\end{savaistu}

\courssub{Les limites de l'artisanat traditionnel}

Si l'artisanat traditionnel est riche culturellement, il souffre de faiblesses structurelles qui expliquent la marginalisation de nombreux artisans évoquée dans la situation d'accroche. Observons un atelier traditionnel de poterie ou de forge : on y constate presque toujours les mêmes contraintes.

\begin{itemize}
\item \textbf{Un outillage rudimentaire} : machette, herminette, marteau, tour à poterie manuel, four à bois. Le travail est lent et physiquement épuisant.
\item \textbf{Une faible productivité} : un sculpteur à la machette peut passer plusieurs jours sur un seul masque ; une potière ne cuit que quelques jarres par fournée. La production ne permet pas de répondre à une commande importante.
\item \textbf{Une qualité irrégulière} : sans normes ni contrôle, deux pièces d'un même artisan peuvent différer fortement ; les dimensions, la finition et la solidité varient, ce qui décourage les acheteurs exigeants (hôtels, exportateurs, supermarchés).
\item \textbf{Une vente locale et informelle} : les produits sont vendus au bord de la route, au marché ou aux touristes de passage, sans étiquette, sans emballage, sans facture. Les circuits de commercialisation sont courts et les marges faibles, souvent captées par les intermédiaires.
\end{itemize}

Ces limites forment un cercle vicieux : faible productivité $\Rightarrow$ faibles revenus $\Rightarrow$ impossibilité d'investir dans de meilleurs outils $\Rightarrow$ faible productivité. C'est précisément ce cercle que la \emph{modernisation} — objet des sections suivantes — cherche à briser.

\courssub{Le poids économique et social du secteur}

Malgré ces limites, l'artisanat joue un rôle considérable dans l'économie camerounaise, car il est l'une des principales composantes du \cle{secteur informel}.

\begin{exemplebox}[Un secteur qui fait vivre des millions de Camerounais]
Selon les enquêtes emploi menées au Cameroun (enquêtes EESI de l'INS), le \textbf{secteur informel} — dont l'artisanat est une composante majeure, aux côtés du petit commerce et de l'agriculture familiale — emploie la grande majorité de la population active : environ \textbf{9 actifs sur 10} y exercent leur activité principale ou secondaire (ordre de grandeur à retenir). Concrètement : tailleurs, coiffeurs, menuisiers, cordonniers, forgerons, potières, vanniers, transformatrices de produits alimentaires... L'artisanat fournit ainsi des \textbf{emplois} et des \textbf{revenus} aussi bien en milieu rural (où il complète l'agriculture pendant les saisons creuses) qu'en milieu urbain (où il absorbe une partie des migrants et des jeunes sans emploi salarié).
\end{exemplebox}

Au-delà de l'économie, l'artisanat a une valeur \textbf{sociale et culturelle} :

\begin{itemize}
\item il \textbf{transmet les savoir-faire} entre générations (apprentissage chez un maître artisan) et structure la vie sociale des quartiers et des villages ;
\item il constitue un pilier de l'\textbf{identité culturelle nationale} : masques, trônes bamiléké, bronzes bamoun, pagnes ndop et gandouras brodées sont des symboles du patrimoine camerounais, exposés dans les musées et recherchés à l'étranger ;
\item il alimente le \textbf{tourisme} et les exportations d'objets d'art, sources de devises modestes mais réelles.
\end{itemize}

\begin{aretenir}[L'essentiel]
L'artisanat est une production manuelle ou semi-manuelle, à petite échelle, maîtrisée de bout en bout par l'artisan. Le Cameroun possède des pôles artisanaux traditionnels puissants (Foumban, Foumbot, Bafoussam, Bamenda, Maroua, Kribi, Batouri), spécialisés selon les matières premières et les cultures. Mais l'outillage rudimentaire, la faible productivité, la qualité irrégulière et la vente informelle limitent ce secteur, qui emploie pourtant, au sein du secteur informel, l'essentiel de la population active. Moderniser l'artisanat, c'est donc agir à la fois sur l'économie, l'emploi et le rayonnement culturel du pays.
\end{aretenir}

\courssub{De l'artisanat traditionnel à l'artisanat modernisé : une comparaison}

Le schéma suivant synthétise, en trois critères, le passage de l'artisanat hérité du passé à l'artisanat en voie de modernisation que nous étudierons dans les sections suivantes.

\begin{popfigure}
\begin{center}
\begin{tikzpicture}[every node/.style={font=\small}]
% Boîte gauche : traditionnel
\node[draw=popmuted, fill=popmuted!15, rounded corners, text width=5.6cm, align=left, minimum height=5.2cm, anchor=north west] (trad) at (0,0) {};
\node[anchor=north, font=\small\bfseries, text=popdark] at (trad.north) {Artisanat traditionnel};
\node[anchor=north west, text width=5.2cm, align=left] at ([yshift=-0.9cm]trad.north west)
{\textbf{Outillage :} machette, herminette, tour manuel, four à bois\\[0.15cm]
\textbf{Production :} lente, pièces uniques, qualité irrégulière\\[0.15cm]
\textbf{Débouchés :} marché local, vente informelle au bord de la route};
% Boîte droite : modernisé
\node[draw=popgreen, fill=popgreenL, rounded corners, text width=5.6cm, align=left, minimum height=5.2cm, anchor=north west] (mod) at (7.6,0) {};
\node[anchor=north, font=\small\bfseries, text=popdark] at (mod.north) {Artisanat modernisé};
\node[anchor=north west, text width=5.2cm, align=left] at ([yshift=-0.9cm]mod.north west)
{\textbf{Outillage :} scies électriques, ponceuses, perceuses, machines à écraser\\[0.15cm]
\textbf{Production :} plus rapide, séries régulières, qualité contrôlée et normée\\[0.15cm]
\textbf{Débouchés :} supermarchés, coopératives, exportation, filières organisées};
% Flèche centrale
\draw[-{Stealth[length=4mm]}, line width=2.5pt, poporange] (5.9,-2.6) -- (7.4,-2.6);
\node[font=\footnotesize\itshape, text=poporange, align=center] at (6.65,-3.1) {modernisation\\(équipements, normes,\\filières)};
\end{tikzpicture}
\end{center}
\legende{Schéma comparatif de l'artisanat traditionnel et de l'artisanat modernisé selon trois critères : outillage, production, débouchés. La modernisation (flèche orange) sera détaillée dans les sections suivantes de la leçon.}
\end{popfigure}

\courssub{Savoir-faire : décrire une photographie d'atelier artisanal}

Au Baccalauréat, tu peux recevoir la photographie d'un atelier (sculpteur de Foumban, potière de Maroua, coopérative de Bafoussam) avec la consigne : « Décrivez ce document ». Voici la démarche attendue.

\begin{methode}[Décrire une photographie d'atelier artisanal : du général au particulier]
\begin{enumerate}
\item \textbf{Identifier le document} : nature (photographie), lieu (atelier, région si précisée), thème (artisanat).
\item \textbf{Décrire le cadre général} : type d'espace (hangar, boutique au marché, atelier en plein air), taille, organisation de l'espace (zone de travail, stockage des matières premières, présentation des produits finis).
\item \textbf{Identifier les acteurs} : nombre de personnes, artisan seul ou avec apprentis, hommes/femmes, répartition des tâches.
\item \textbf{Repérer les outils et équipements} : outils manuels ou machines électriques ; c'est l'indice clé pour juger du degré de modernisation.
\item \textbf{Décrire les produits} : nature, stade de fabrication, conditionnement (en vrac ou emballés, étiquetés).
\item \textbf{Conclure par une interprétation} : atelier traditionnel ou modernisé ? Quels indices le prouvent ? Quelles difficultés ou quels atouts peut-on en déduire ?
\end{enumerate}
\end{methode}

\begin{exoresolu}[Décrire une photographie d'atelier de sculpture à Foumban]
\textbf{Énoncé.} Sur une photographie prise dans le quartier des artisans de Foumban, on observe un homme assis devant un hangar en tôle, entouré de billes de bois et de masques accrochés au mur. Il tient une herminette ; à ses pieds, un jeune garçon dégrossit une statuette à la machette. Aucun emballage n'est visible ; les œuvres sont empilées au sol. Décrivez cette photographie et dégagez le niveau de modernisation de l'atelier.
\tcblower
\textbf{Solution détaillée.}
\emph{1. Identification} : il s'agit d'une photographie montrant un atelier de sculpture sur bois à Foumban (Ouest-Cameroun), pôle artisanal traditionnel.
\emph{2. Cadre général} : l'atelier est un hangar sommaire en tôle, ouvert sur l'extérieur ; l'espace sert à la fois de lieu de production (billes de bois au sol) et de lieu d'exposition-vente (masques accrochés au mur).
\emph{3. Acteurs} : deux personnes — un maître artisan (herminette à la main, travail de finition) et un apprenti (dégrossissage à la machette) — ce qui illustre la transmission du savoir-faire par apprentissage et la faible division du travail.
\emph{4. Outils} : uniquement des outils manuels (herminette, machette) ; aucune machine électrique n'est visible.
\emph{5. Produits} : masques et statuettes, vendus en l'état, sans emballage ni étiquette, empilés au sol.
\emph{6. Interprétation} : il s'agit d'un atelier \textbf{traditionnel} : outillage rudimentaire, production lente et irrégulière, vente informelle probable (touristes, marché local). Pour se moderniser, cet artisan devrait s'équiper (scie électrique, ponceuse), améliorer la finition et trouver des circuits de vente organisés — exactement le chemin parcouru par le sculpteur de la situation d'accroche.
\end{exoresolu}

\begin{exercice}[Pour t'entraîner]
\begin{enumerate}
\item Définis l'artisanat en citant trois de ses caractéristiques, puis explique en quatre lignes ce qui le distingue de l'industrie.
\item Sur un croquis simplifié du Cameroun, place et nomme quatre pôles artisanaux traditionnels en précisant leur spécialité dominante (n'oublie ni le titre, ni la légende, ni l'orientation).
\item Explique, à partir d'un exemple choisi, comment les limites de l'artisanat traditionnel forment un « cercle vicieux » de la pauvreté artisanale.
\end{enumerate}
\end{exercice}

\begin{corrige}[Exercice]
\begin{enumerate}
\item L'artisanat est une activité de transformation et de production réalisée manuellement ou avec des outils simples, par un artisan seul ou avec peu de main-d'œuvre. Caractéristiques (au choix) : production à petite échelle ; faible division du travail ; atelier familial ; transmission par apprentissage. Contrairement à l'industrie (mécanisation, production de masse, tâches parcellaires), l'artisan maîtrise son produit de la matière première à la vente.
\item Croquis attendu : contour simplifié du Cameroun, flèche du nord, titre et légende. Exemples de pôles : Foumban (sculpture, bronze), Bafoussam (broderie, vannerie), Bamenda (tissage ndop, vannerie), Maroua (cuir, poterie), Kribi (vannerie, nattes), Foumbot (poterie), Batouri (sculpture).
\item Exemple du potier traditionnel : son four à bois et son tour manuel limitent sa production à quelques jarres par semaine (faible productivité) ; ses revenus restent donc faibles ; il ne peut pas épargner pour acheter un four amélioré ou un tour électrique ; sa production reste donc faible : le cercle vicieux se referme. La modernisation des équipements et l'organisation des filières visent à rompre ce cercle.
\end{enumerate}
\end{corrige}

\coursec{La modernisation des équipements de travail}

Dans la section précédente, nous avons vu que l'artisanat camerounais s'enracine dans un riche héritage traditionnel : la poterie de Foumban, les tisserands de l'Extrême-Nord, les sculpteurs de l'Ouest, les vanneries du Sud travaillaient depuis des générations avec des outils purement manuels — la hache, le couteau, la râpe en bois, le tour à poterie actionné à la main. Cet héritage n'a pas disparu : il s'est transformé. Imagine un menuisier de la rue des Briqueteries à Yaoundé : son père passait une journée entière à débiter une planche avec une scie à main ; lui, en une heure, grâce à une scie électrique, il débite la même quantité de bois, avec des coupes plus régulières. C'est ce passage de l'outil manuel à la machine que nous allons étudier : la \cle{modernisation des équipements de travail}.

\courssub{Du geste manuel à la machine : la mécanisation de l'artisanat}

Commençons par une intuition simple. Pourquoi un artisan adopte-t-il une machine ? Parce que la machine fait \emph{plus vite}, \emph{plus régulièrement} et \emph{avec moins d'effort} ce que la main faisait lentement et péniblement. C'est exactement la définition de la mécanisation.

\begin{definition}[Mécanisation]
La \cle{mécanisation} est le remplacement du travail purement manuel par des machines, actionnées le plus souvent par une énergie extérieure (électricité, moteur thermique). Elle augmente la \cle{productivité} de l'artisan, c'est-à-dire la quantité de produits fabriqués par unité de temps, tout en réduisant la pénibilité du travail.
\end{definition}

Au Cameroun, cette mécanisation touche toutes les branches de l'artisanat. Inventorions les principaux équipements modernes :

\begin{itemize}
\item \textbf{Les machines à écraser et à presser} : dans la transformation agroalimentaire, les moulins et machines à écraser traitent le manioc (pour la farine, les bâtons de manioc, le gari), l'arachide (pâte et huile) et les épices (piments, poivre de Penja). Les presses à huile mécaniques remplacent le pilage au mortier.
\item \textbf{Les scies électriques} : scies circulaires, scies à ruban et tronçonneuses ont transformé le travail des menuisiers et ébénistes, notamment dans les quartiers artisanaux de Yaoundé (Nkolndongo, Briqueterie) et de Douala (New-Bell, Bonabéri).
\item \textbf{Les ponceuses et perceuses} : elles assurent des surfaces lisses et des assemblages précis, impossibles à obtenir régulièrement avec le papier de verre et la vrille manuels.
\item \textbf{Les tours à poterie électriques} : ils permettent un tournage régulier et rapide des poteries, là où le tour manuel exigeait un assistant et beaucoup d'adresse.
\item \textbf{Les machines à coudre industrielles} : dans les ateliers de couture et de maroquinerie, elles remplacent la machine à pédale et permettent des points solides, rapides et uniformes.
\end{itemize}

\begin{popfigure}
\begin{tikzpicture}[scale=0.95, every node/.style={font=\small}]
% Sol
\draw[thick, popdark] (0,0) -- (13,0);
\fill[popcream] (0,0) rectangle (13,0.15);
% Établi
\fill[popgoldL] (0.5,0.15) rectangle (3.5,1.3);
\draw[popdark, thick] (0.5,0.15) rectangle (3.5,1.3);
\draw[popdark] (0.7,1.3) -- (0.7,0.15); \draw[popdark] (3.3,1.3) -- (3.3,0.15);
\node[popdark] at (2,0.75) {\textbf{Établi}};
% Planche sur établi
\fill[poporangeL] (0.7,1.3) rectangle (3.3,1.55);
\draw[popdark] (0.7,1.3) rectangle (3.3,1.55);
% Scie circulaire
\draw[popblue, very thick] (5.2,0.15) rectangle (7.4,1.1);
\draw[popblue, very thick, fill=popblueL] (6.3,1.6) circle (0.55);
\foreach \a in {0,45,...,315} \draw[popblue] (6.3,1.6) -- ++(\a:0.55);
\draw[popblue, thick] (6.3,1.6) -- (6.3,1.1);
\node[popblue, font=\small\bfseries] at (6.3,0.55) {Scie électrique};
% Planche en cours de coupe
\fill[poporangeL] (4.9,1.1) rectangle (7.7,1.35);
\draw[popdark] (4.9,1.1) rectangle (7.7,1.35);
% Ponceuse
\draw[popgreen, very thick, fill=popgreenL, rounded corners] (8.6,0.15) rectangle (10.1,1.0);
\draw[popgreen, thick] (9.35,1.0) -- (9.35,1.5);
\draw[popgreen, thick, fill=popgreenL] (9.0,1.5) rectangle (9.7,1.8);
\node[popgreen, font=\small\bfseries] at (9.35,0.5) {Ponceuse};
% Perceuse à colonne
\draw[poppurple, very thick] (11.2,0.15) -- (11.2,2.6);
\draw[poppurple, very thick] (11.2,2.6) -- (12.4,2.6);
\draw[poppurple, very thick, fill=poppurpleL] (11.9,2.1) rectangle (12.6,2.6);
\draw[poppurple, thick] (12.25,2.1) -- (12.25,1.5);
\draw[poppurple, thick] (11.5,1.0) -- (12.9,1.0);
\draw[poppurple, thick] (11.2,0.15) -- (12.9,0.15);
\node[poppurple, font=\small\bfseries] at (12.1,0.5) {Perceuse};
% Prise électrique
\draw[poporange, thick, fill=poporangeL] (5.6,3.4) rectangle (6.4,3.9);
\node[poporange, font=\footnotesize] at (6,3.65) {220 V};
\draw[poporange, dashed, thick] (6,3.4) -- (6.3,2.2);
\draw[poporange, dashed, thick] (6,3.4) -- (9.35,1.9);
\draw[poporange, dashed, thick] (6,3.4) -- (12.25,2.7);
\node[poporange, font=\footnotesize, align=center] at (3.2,3.65) {Arrivée du\\courant électrique};
% Artisan
\draw[popink, thick, fill=poppinkL] (4.2,2.5) circle (0.3);
\draw[popink, thick] (4.2,2.2) -- (4.2,1.3);
\draw[popink, thick] (4.2,1.9) -- (4.9,1.5);
\draw[popink, thick] (4.2,1.9) -- (3.6,1.5);
\draw[popink, thick] (4.2,1.3) -- (3.9,0.4);
\draw[popink, thick] (4.2,1.3) -- (4.5,0.4);
\node[popink, font=\footnotesize] at (4.2,3.0) {L'artisan};
\end{tikzpicture}
\legende{Schéma d'un atelier de menuiserie modernisé : 1. établi de travail ; 2. scie électrique circulaire pour le débit des planches ; 3. ponceuse pour le lissage des surfaces ; 4. perceuse à colonne pour les assemblages ; 5. alimentation électrique commune (220 V) reliant toutes les machines. L'artisan reste au centre du dispositif : c'est lui qui guide chaque machine.}
\end{popfigure}

\courssub{Les effets de la mécanisation sur la production}

Pour mesurer rigoureusement les effets de la mécanisation, menons une petite démarche d'investigation : comparons la production journalière d'un même atelier de menuiserie avant et après l'acquisition de machines (scie électrique, ponceuse, perceuse). Les données ci-dessous sont fictives mais réalistes, construites à partir d'ordres de grandeur observés dans les petits ateliers urbains camerounais.

\begin{popfigure}
\begin{center}
\begin{tikzpicture}
\begin{axis}[
    ybar, bar width=22pt,
    width=0.85\linewidth, height=6.5cm,
    ymin=0, ymax=14,
    ylabel={Production journalière (unités)},
    symbolic x coords={Planches débitées, Chaises poncées, Trous percés},
    xtick=data,
    x tick label style={font=\small},
    ytick={0,2,4,6,8,10,12,14},
    legend style={at={(0.5,1.02)}, anchor=south, legend columns=2, font=\small},
    nodes near coords, nodes near coords style={font=\small},
    every axis plot/.append style={draw=popdark}
]
\addplot[fill=popmuted!60] coordinates {(Planches débitées,3) (Chaises poncées,2) (Trous percés,6)};
\addplot[fill=popblue] coordinates {(Planches débitées,12) (Chaises poncées,9) (Trous percés,13)};
\legend{Avant (outils manuels), Après (machines électriques)}
\end{axis}
\end{tikzpicture}
\end{center}
\legende{Comparaison de la production journalière d'un atelier de menuiserie avant et après mécanisation (données fictives réalistes, en unités par journée de travail). La production est multipliée par 3 à 4 pour le débitage et le ponçage.}
\end{popfigure}

Lecture du graphique : le débitage des planches passe de 3 à 12 unités par jour, soit une multiplication par 4 ; le ponçage passe de 2 à 9 chaises (multiplication par 4,5) ; le perçage est plus que doublé. On en tire quatre effets généraux de la mécanisation :

\begin{enumerate}
\item \textbf{Le gain de temps} : les opérations les plus longues (sciage, écrasement, ponçage) sont raccourcies de manière spectaculaire.
\item \textbf{L'augmentation de la production} : à journée de travail égale, l'artisan produit davantage, ce qui accroît son revenu potentiel.
\item \textbf{La réduction de la pénibilité} : l'effort musculaire (scier, piler, pédaler) est transféré à la machine ; l'artisan se fatigue moins et préserve sa santé.
\item \textbf{La régularité des produits} : les coupes sont droites, les surfaces uniformes, les points de couture identiques. Cette régularité prépare l'artisan aux exigences de \cle{qualité} et de \cle{norme} que nous étudierons plus loin.
\end{enumerate}

\begin{exemplebox}[Un exemple chiffré : la scie électrique]
Un menuisier de New-Bell (Douala) qui débitait à la scie à main une dizaine de planches par journée complète de travail débite désormais, avec une scie circulaire électrique, la même quantité en environ une heure. Autrement dit, \textbf{une heure de sciage mécanisé équivaut à une journée de sciage manuel}. Le temps libéré est réinvesti dans le façonnage, l'assemblage et la finition, étapes où la valeur ajoutée est la plus forte.
\end{exemplebox}

\begin{exoresolu}[Calculer un gain de productivité]
Énoncé : une transformatrice de manioc de la région du Centre pilait manuellement environ 20 kg de manioc par jour. Après l'achat d'une machine à écraser, elle traite 100 kg par jour. Calcule le taux d'augmentation de sa production journalière.
\tcblower
Solution détaillée : le taux de variation se calcule par la formule
\[ t = \frac{V_{\text{après}} - V_{\text{avant}}}{V_{\text{avant}}} \times 100 = \frac{100 - 20}{20} \times 100 = \frac{80}{20} \times 100 = 400\%. \]
La production journalière a donc augmenté de 400 \% : elle a été multipliée par 5. Ce gain permet à la transformatrice d'approvisionner régulièrement les marchés de sa commune et d'augmenter son chiffre d'affaires, à condition de maîtriser les coûts (électricité ou carburant, entretien de la machine).
\end{exoresolu}

\courssub{L'approvisionnement en équipements : d'où viennent les machines ?}

Les artisans camerounais s'équipent par trois voies principales :

\begin{itemize}
\item \textbf{L'importation} : la plupart des machines (scies, ponceuses, machines à coudre industrielles, moulins) sont importées d'Asie (Chine notamment), d'Europe ou du Nigeria voisin, puis vendues dans les quincailleries et les marchés d'équipements de Douala et de Yaoundé. Les machines d'occasion, moins chères, occupent une place importante.
\item \textbf{La fabrication locale adaptée} : des forgerons et mécaniciens camerounais construisent eux-mêmes des machines simples et robustes — presses à huile, décortiqueuses d'arachide, râpes à manioc motorisées — adaptées aux matières premières locales et réparables sur place.
\item \textbf{Les appuis de projets et de coopératives} : l'État (à travers le ministère en charge de l'artisanat, le MINPMEESA, et les programmes d'appui), les ONG, les coopératives et les GIC (Groupements d'Initiative Commune) financent ou subventionnent l'achat d'équipements collectifs, par exemple une machine à écraser partagée par un groupement de femmes transformatrices.
\end{itemize}

\begin{savaistu}[L'innovation frugale, une force camerounaise]
Certains artisans camerounais fabriquent eux-mêmes leurs machines à partir de matériaux de récupération : une presse à huile montée sur une vieille transmission de voiture, une décortiqueuse construite avec des pièces de moto hors d'usage, un tour à poterie entraîné par un moteur de machine à laver. Cette ingéniosité porte un nom : l'\cle{innovation frugale}, c'est-à-dire la capacité à produire des solutions techniques efficaces avec très peu de moyens. Elle réduit le coût d'accès à la mécanisation et garantit que les pièces de rechange restent disponibles localement.
\end{savaistu}

\courssub{Localisation des pôles artisanaux modernisés}

La modernisation n'est pas uniforme sur le territoire. Elle se concentre dans les grands pôles urbains et les bassins de production agricole.

\begin{popfigure}
\begin{center}
\includegraphics[width=\linewidth,height=9.5cm,keepaspectratio]{N01/cmr_map.jpg}
\end{center}
\legende{Fond de carte : carte générale du Cameroun. Les pôles où la modernisation de l'artisanat est la plus avancée sont les grandes villes et leurs régions : Douala et Yaoundé, puis Bafoussam et Bamenda dans l'Ouest et le Nord-Ouest. On les repère sur la carte à leur position le long des routes et du chemin de fer. \\ Source : Wikimedia Commons, CIA, domaine public.}
\end{popfigure}

\courssub{Les obstacles à la modernisation}

La mécanisation n'est pas automatique : de nombreux artisans restent bloqués par quatre obstacles majeurs.

\begin{enumerate}
\item \textbf{Le coût des machines} : une scie circulaire, un moulin ou une machine à coudre industrielle représentent un investissement de plusieurs centaines de milliers de francs CFA, souvent inaccessible sans crédit bancaire — or les artisans ont peu accès au financement.
\item \textbf{L'accès à l'électricité} : dans de nombreuses zones rurales et même dans certains quartiers urbains, l'électricité est absente, chère ou instable (délestages fréquents). Une machine électrique sans courant fiable est inutilisable ; certains artisans recourent alors à des groupes électrogènes, ce qui alourdit les coûts.
\item \textbf{L'entretien et les pièces détachées} : une machine importée dont la pièce de rechange est introuvable au Cameroun peut rester immobilisée des mois. La maintenance exige aussi des techniciens qualifiés, rares en milieu rural.
\item \textbf{La formation des artisans} : utiliser, régler et entretenir une machine demande un apprentissage. Sans formation, la machine est mal utilisée, s'use vite et peut provoquer des accidents (les scies électriques sont dangereuses sans consignes de sécurité).
\end{enumerate}

\begin{attention}[La machine ne remplace pas le savoir-faire]
Erreur fréquente : croire que la modernisation supprime l'artisanat. C'est faux. La mécanisation \textbf{amplifie} le savoir-faire artisanal, elle ne l'abolit pas : la scie électrique ne choisit ni l'essence de bois, ni le sens de la coupe, ni le dessin du meuble — c'est toujours l'œil et la main de l'artisan qui guident la machine. Un artisan sans compétence produira des défauts... plus rapidement. À l'examen, évite donc d'opposer « artisanat traditionnel » et « artisanat moderne » comme deux mondes séparés : il s'agit d'une \emph{continuité transformée}.
\end{attention}

\courssub{Des exemples camerounais de modernisation réussie}

Deux cas concrets illustrent cette dynamique. D'une part, les \textbf{menuisiers de Yaoundé et de Douala} : regroupés dans des quartiers spécialisés (Briqueterie, Nkolndongo à Yaoundé ; New-Bell, Bonabéri à Douala), ils se sont massivement équipés de scies électriques, de ponceuses et de perceuses, ce qui leur permet de fournir meubles, portes et cadres de fenêtres aux chantiers urbains en forte croissance. D'autre part, les \textbf{transformatrices de manioc du Centre et du Littoral} : organisées en GIC et en coopératives, souvent appuyées par des projets de développement (ex. PADMIR, PEA-Jeunes), elles utilisent des presses et des machines à écraser pour produire bâtons de manioc, gari et farine en quantités régulières, destinées aux marchés urbains. Dans les deux cas, la mécanisation a transformé une activité de subsistance en véritable petite entreprise.

\begin{methode}[Analyser un document photographique d'équipement artisanal]
Au Baccalauréat, on peut te présenter la photographie d'un atelier ou d'une machine. Pour l'analyser rigoureusement, procède en quatre étapes :
\begin{enumerate}
\item \textbf{Identifier l'outil ou la machine} : nomme-le précisément (scie circulaire, presse à huile, machine à coudre industrielle...).
\item \textbf{Dégager sa fonction} : quelle opération effectue-t-il (couper, écraser, lisser, assembler) et dans quelle filière (bois, agroalimentaire, textile...) ?
\item \textbf{Préciser sa source d'énergie} : manuelle, électrique ou thermique ? C'est le critère qui distingue l'outil traditionnel de la machine moderne.
\item \textbf{Évaluer son effet sur la production} : gain de temps, quantité produite, régularité, réduction de la pénibilité — et, si possible, les contraintes associées (coût, électricité, entretien).
\end{enumerate}
Conclus toujours en reliant le document à la notion du cours : ici, la modernisation des équipements de travail.
\end{methode}

\begin{exercice}[À toi de jouer]
Une coopérative de femmes de la région du Littoral hésite à acheter une machine à écraser le manioc coûtant 450 000 FCFA. Actuellement, ses membres pilent 150 kg de manioc par jour à la main ; la machine permettrait d'en traiter 600 kg par jour. 1) Calcule le taux d'augmentation de la production. 2) Présente deux avantages et deux obstacles que la coopérative doit examiner avant l'achat.
\end{exercice}

\begin{corrige}[Exercice : la machine à écraser du Littoral]
1) Taux d'augmentation : $t = \dfrac{600 - 150}{150} \times 100 = \dfrac{450}{150} \times 100 = 300\%$. La production serait multipliée par 4.
2) Avantages possibles : gain de temps et réduction de la pénibilité pour les femmes ; augmentation du chiffre d'affaires grâce à une production régulière permettant de fournir les marchés urbains. Obstacles à examiner : le coût élevé de la machine (450 000 FCFA) qui suppose un crédit ou une subvention ; la disponibilité de l'électricité (ou le coût d'un groupe électrogène) ainsi que l'entretien et les pièces détachées ; on peut ajouter la nécessité de former les utilisatrices.
\end{corrige}

\begin{aretenir}[L'essentiel de la section]
La modernisation des équipements de travail désigne le passage de l'outil manuel à la machine : machines à écraser (manioc, arachide, épices), scies électriques, ponceuses, perceuses, tours à poterie électriques, machines à coudre industrielles. Cette \cle{mécanisation} apporte un gain de temps, une augmentation de la production, une réduction de la pénibilité et une meilleure régularité des produits. Les équipements proviennent de l'importation, de la fabrication locale adaptée (innovation frugale) et des appuis de projets et coopératives. Mais quatre obstacles freinent la modernisation : le coût des machines, l'accès à l'électricité, l'entretien et les pièces détachées, la formation des artisans. Enfin, la machine ne supprime pas le savoir-faire : elle l'amplifie, car elle reste guidée par la compétence de l'artisan.
\end{aretenir}

Cette modernisation des outils a une conséquence directe : produisant plus vite et plus régulièrement, l'artisan peut désormais viser de nouveaux clients et imaginer de nouveaux produits. C'est précisément l'objet de la section suivante : des produits nouveaux pour des marchés nouveaux.

\coursec{Des produits nouveaux pour des marchés nouveaux}

Dans la section précédente, nous avons vu comment l'artisan camerounais a modernisé ses \cle{équipements de travail} : la scie électrique a remplacé la scie à main dans les ateliers d'ébénisterie de Foumban, la machine à écraser a transformé le travail des productrices d'huile de palme, la ponceuse a rendu les surfaces plus lisses. Mais un outil moderne ne suffit pas : encore faut-il fabriquer quelque chose que le client veut acheter. C'est précisément l'objet de cette section : comprendre comment l'artisanat camerounais a renouvelé son \cle{offre} en créant des \cle{produits nouveaux} adaptés à des marchés nouveaux — la ville, le tourisme, la diaspora, l'exportation.

\begin{definition}[Produit nouveau]
Un \cle{produit nouveau} est un bien artisanal \textbf{inédit} (qui n'existait pas dans l'artisanat traditionnel) ou \textbf{relooké} (produit traditionnel dont la forme, la matière, la finition ou la présentation ont été transformées), fabriqué pour répondre à une \textbf{demande nouvelle} : clientèle urbaine, touristique, diaspora ou marché d'exportation. Le produit nouveau n'est donc pas défini seulement par sa technique de fabrication, mais par le \emph{marché} auquel il s'adresse.
\end{definition}

\courssub{Des produits nouveaux dans la décoration}

\textbf{Intuition.} Entre dans un salon moderne à Yaoundé ou à Douala : tu y verras souvent un salon en cuir ou en tissu aux lignes contemporaines, mais dont les accoudoirs portent des motifs sculptés bamiléké ; un luminaire en rotin tressé ; des cadres en bois précieux (ébène, iroko) ; parfois un tabouret royal transformé en table basse. Rien de tout cela n'existait dans la case traditionnelle : l'artisan a \emph{réinterprété} son héritage.

\textbf{Analyse.} La décoration est le domaine où la créativité artisanale camerounaise est la plus visible. On distingue plusieurs familles de produits nouveaux :

\begin{itemize}
\item les \textbf{meubles modernes inspirés des motifs traditionnels} : salons complets, tables, lits et bibliothèques produits dans les ateliers d'ébénisterie de Foumban, de Bafoussam, de Yaoundé (quartier Mvog-Ada) et de Douala (New-Bell), qui combinent bois local, cuir, tissu et sculpture ;
\item les \textbf{objets d'art pour le tourisme} : masques, statuettes, tambours, pipes et plateaux sculptés, vendus dans les villages artisanaux (marché artisanal de Yaoundé, centre artisanal de Foumban près du Palais royal des Bamoun) ;
\item les \textbf{luminaires et cadres} : lampes en rotin, en calebasse ou en bois tourné, cadres de tableaux en bois précieux ;
\item les \textbf{articles en matériaux recyclés} : sacs et porte-monnaie tissés à partir d'emballages plastiques, sculptures en métal récupéré (bidons, ferraille), objets en pneus usagés — une filière en pleine expansion à Douala et Yaoundé, qui répond à la fois au souci écologique et à la recherche de prix bas.
\end{itemize}

\begin{exemplebox}[Le tabouret bamiléké devient table basse]
Le tabouret royal sculpté dans un seul bloc de bois, autrefois réservé aux notables des chefferies de l'Ouest, est aujourd'hui reproduit en série et vendu comme table basse ou élément de décoration intérieure aux citadins camerounais, aux touristes et même à des acheteurs européens via Internet. L'objet n'a pas changé de technique de fabrication ; il a changé de \emph{fonction} et de \emph{clientèle} : c'est un produit relooké.
\end{exemplebox}

\courssub{Des produits nouveaux dans l'habillement}

\textbf{Intuition.} Lors d'une fête scolaire ou d'un mariage, observe les tenues : le pagne \emph{ndop} des Grassfields ou le pagne tissé \emph{kaba ngondo} de la côte ne sont plus seulement portés en grand boubou traditionnel ; ils sont coupés en vestes cintrées, en pantalons, en robes de soirée, en sacs à main et même en baskets.

\textbf{Analyse.} L'habillement illustre parfaitement le passage du produit traditionnel au produit nouveau :

\begin{itemize}
\item les \textbf{pagnes et tissus teints modernisés} : le \emph{ndop} (tissu de coton teint à la réserve, originaire de la région du Nord-Ouest et de l'Ouest) et les bazins teints sont déclinés en nouvelles couleurs et nouveaux motifs pour séduire une clientèle jeune et urbaine ;
\item la \textbf{broderie sur vêtements contemporains} : la broderie du Grand Nord (héritée de la tradition des boubous) orne désormais chemises, robes et vestes de coupe européenne ;
\item les \textbf{accessoires de mode} : sacs en cuir ou en raphia, bijoux en perles, en os, en graines ou en bronze (fondeurs de Foumban), chaussures en cuir fabriquées par les cordonniers de Douala et de Maroua, chapeaux et ceintures tressés.
\end{itemize}

Ces produits visent une clientèle qui veut afficher une identité culturelle camerounaise tout en suivant la mode contemporaine : c'est le phénomène du \emph{« porter local »}, encouragé par les stylistes et créateurs de mode camerounais qui collaborent directement avec les artisans.

\courssub{Des produits nouveaux dans la transformation des produits alimentaires}

\textbf{Intuition.} Sur les étals du marché Mokolo à Yaoundé ou du marché Mboppi à Douala, à côté des tas de manioc et de plantain, on trouve désormais des bouteilles de jus de gingembre et de foléré (bissap), des sachets d'épices moulues et étiquetées, des chips de banane plantain emballées, des pots de miel. Ces produits sont le résultat d'un artisanat de \cle{transformation} qui ajoute de la valeur aux produits agricoles bruts.

\textbf{Analyse.} L'artisanat agroalimentaire est le secteur où les produits nouveaux se multiplient le plus vite :

\begin{itemize}
\item les \textbf{jus de fruits artisanaux} : jus d'ananas, de mangue, de goyave, de gingembre, de foléré, produits en petites unités et vendus frais en ville ;
\item les \textbf{épices conditionnées} : piment de Penja (renommé pour sa qualité), poivre de Penja, njansang, pebé, pèbè et mélanges moulus, vendus en sachets ou en pots étiquetés ;
\item les \textbf{farines} : farine de manioc, de maïs, d'igname, conditionnées en sachets pour les ménages urbains pressés ;
\item les \textbf{huiles} : huile de palme rouge, huile d'arachide et huile de coco artisanales ;
\item le \textbf{miel} : miel de l'Adamaoua et de l'Ouest, filtré et mis en pot ;
\item les \textbf{chips de banane plantain et de manioc} : frites, séchées et emballées, elles concurrencent les chips industrielles importées ;
\item les \textbf{produits laitiers locaux} : lait caillé et fromage peul de l'Adamaoua et du Nord.
\end{itemize}

\begin{exemplebox}[Jus artisanaux et épices conditionnées : une part croissante des étals]
Sur les marchés de Yaoundé et de Douala, les jus de fruits artisanaux et les épices conditionnées occupent une part croissante des étals. Ordre de grandeur indicatif : dans certains marchés urbains, les produits alimentaires transformés artisanalement représenteraient aujourd'hui environ un tiers des produits artisanaux vendus, contre une part très faible il y a une vingtaine d'années (donnée d'observation de terrain, à manier avec prudence, les statistiques officielles de l'artisanat restant fragmentaires). Ce succès s'explique par des prix inférieurs à ceux des produits industriels et par le goût des consommateurs pour les saveurs locales.
\end{exemplebox}

\courssub{L'adaptation de l'offre à la demande : des marchés nouveaux}

Le produit nouveau n'existe que parce qu'il existe une \cle{demande nouvelle}. L'artisan camerounais adapte son offre à quatre grandes clientèles :

\begin{enumerate}
\item \textbf{les goûts urbains} : la croissance rapide des villes (Douala et Yaoundé comptent chacune plusieurs millions d'habitants, ordre de grandeur : plus de 3 millions chacune au début des années 2020) crée une demande de meubles modernes, de vêtements pratiques, d'aliments transformés rapides à cuisiner ;
\item \textbf{la clientèle touristique} : les touristes recherchent des objets d'art authentiques, portables et présentés (masques, statuettes, bijoux), ce qui pousse les artisans à standardiser les formats et à soigner l'emballage ;
\item \textbf{la diaspora} : les Camerounais vivant à l'étranger commandent pagnes, épices, œuvres d'art et produits alimentaires, entretenant un marché de la nostalgie et de l'identité ;
\item \textbf{l'exportation} : certains produits (piment et poivre de Penja, café et cacao artisanaux transformés, objets d'art) atteignent les marchés africains, européens et américains, parfois via les plateformes de vente en ligne.
\end{enumerate}

\begin{savaistu}[Le « made in Cameroon » gagne les supermarchés]
Longtemps confinés aux marchés et aux bordures de routes, les produits artisanaux camerounais entrent désormais dans les rayons des supermarchés de Douala et de Yaoundé : miel de l'Adamaoua en pot étiqueté, café artisanal torréfié et moulu de l'Ouest (région du café arabica des hautes terres), piment de Penja conditionné en sachets. Le poivre de Penja bénéficie même d'une Indication Géographique Protégée (IGP), premier label de ce type en Afrique subsaharienne, obtenu dans le cadre de l'Organisation Africaine de la Propriété Intellectuelle (OAPI) : preuve qu'un produit local de qualité peut conquérir un marché mondial.
\end{savaistu}

\courssub{L'innovation comme réponse à la concurrence des produits industriels importés}

\textbf{Problème.} Les marchés camerounais sont inondés de produits industriels importés : meubles en mélaminé, chaussures et sacs asiatiques à bas prix, jus en briques, chips industrielles. Comment l'artisan peut-il survivre face à cette concurrence ?

\textbf{Réponse.} Par l'\cle{innovation}, c'est-à-dire par la création de produits nouveaux qui exploitent les avantages propres de l'artisanat :

\begin{itemize}
\item l'\textbf{authenticité} : un masque sculpté à Foumban ou un pagne ndop véritable ne peuvent pas être copiés par une usine étrangère sans perdre leur valeur symbolique ;
\item la \textbf{personnalisation} : l'ébéniste fabrique un meuble sur mesure, adapté au salon du client, ce qu'aucun meuble importé standardisé ne peut offrir ;
\item la \textbf{fraîcheur et le goût local} : un jus de gingembre artisanal, des épices fraîchement moulues ou des chips de plantain répondent à des préférences alimentaires que les produits importés ne satisfont pas ;
\item le \textbf{prix et la proximité} : les matériaux recyclés et les circuits courts (vente directe à l'atelier ou au marché) permettent de rester compétitif.
\end{itemize}

L'innovation artisanale est donc une \emph{stratégie économique} : elle permet à l'artisan de se différencier, de fidéliser une clientèle et de contribuer à l'\cle{intégration nationale} en valorisant les ressources et les identités culturelles de toutes les régions du Cameroun.

\begin{attention}[Un produit nouveau n'est pas forcément un produit de qualité]
Attention au piège de raisonnement : innover ne suffit pas. Un jus de fruit mal conditionné peut être contaminé ; un meuble mal assemblé se casse ; un colorant textile de mauvaise qualité déteint au premier lavage. Pour durer sur le marché — surtout face aux produits industriels importés — l'innovation doit s'accompagner du \cle{respect des normes} (hygiène, solidité, étiquetage, conditionnement). C'est l'objet de la section suivante. Dans une copie de Bac, ne confonds jamais « produit nouveau » (critère de marché) et « produit de qualité » (critère de conformité aux normes).
\end{attention}

\begin{popfigure}
\begin{center}
\begin{tikzpicture}[font=\small]
% Titre des colonnes
\fill[popblueL] (0,0) rectangle (4.6,0.8);
\fill[poporangeL] (4.8,0) rectangle (9.4,0.8);
\fill[popgreenL] (9.6,0) rectangle (14.2,0.8);
\node[text=popdark,font=\bfseries] at (2.3,0.4) {D\'ecoration};
\node[text=popdark,font=\bfseries] at (7.1,0.4) {Habillement};
\node[text=popdark,font=\bfseries] at (11.9,0.4) {Agroalimentaire};
% Corps des colonnes
\fill[popblue!8] (0,-4.6) rectangle (4.6,0);
\fill[poporange!8] (4.8,-4.6) rectangle (9.4,0);
\fill[popgreen!8] (9.6,-4.6) rectangle (14.2,0);
\draw[popline] (0,0.8) rectangle (4.6,-4.6);
\draw[popline] (4.8,0.8) rectangle (9.4,-4.6);
\draw[popline] (9.6,0.8) rectangle (14.2,-4.6);
% Contenu colonne 1
\node[text=popdark,align=left,anchor=north west] at (0.15,-0.15)
{\textbullet\ Meubles modernes \`a motifs\\ \hspace{0.3cm}traditionnels (Foumban)\\[2pt]
 \textbullet\ Objets d'art pour le tourisme\\ \hspace{0.3cm}(masques, statuettes)\\[2pt]
 \textbullet\ Luminaires en rotin, cadres\\ \hspace{0.3cm}en bois pr\'ecieux\\[2pt]
 \textbullet\ Articles en mat\'eriaux\\ \hspace{0.3cm}recycl\'es (plastique, m\'etal)};
% Contenu colonne 2
\node[text=popdark,align=left,anchor=north west] at (4.95,-0.15)
{\textbullet\ Pagnes modernis\'es (ndop,\\ \hspace{0.3cm}kaba ngondo)\\[2pt]
 \textbullet\ Broderie du Grand Nord sur\\ \hspace{0.3cm}v\^etements contemporains\\[2pt]
 \textbullet\ Sacs en cuir et raphia\\[2pt]
 \textbullet\ Bijoux en perles, bronze\\[2pt]
 \textbullet\ Chaussures en cuir (Douala,\\ \hspace{0.3cm}Maroua)};
% Contenu colonne 3
\node[text=popdark,align=left,anchor=north west] at (9.75,-0.15)
{\textbullet\ Jus de fruits artisanaux\\ \hspace{0.3cm}(fol\'er\'e, gingembre, mangue)\\[2pt]
 \textbullet\ \'Epices conditionn\'ees (piment\\ \hspace{0.3cm}et poivre de Penja, njansang)\\[2pt]
 \textbullet\ Farines, huiles, miel de\\ \hspace{0.3cm}l'Adamaoua\\[2pt]
 \textbullet\ Chips de banane et de manioc\\[2pt]
 \textbullet\ Produits laitiers locaux};
\end{tikzpicture}
\legende{Classification des produits artisanaux nouveaux camerounais en trois domaines : (1) d\'ecoration, (2) habillement, (3) transformation agroalimentaire. Chaque domaine combine h\'eritage traditionnel et adaptation \`a une demande nouvelle (urbaine, touristique, diaspora, export).}
\end{center}
\end{popfigure}

\begin{popfigure}
\begin{center}
\begin{tikzpicture}[font=\small]
% Diagramme circulaire : parts fictives mais r\'ealistes
% Agroalimentaire 35\% (0 -> 126 deg)
\fill[popgreen] (0,0) -- (126:2.6) arc (126:0:2.6) -- cycle;
% D\'ecoration 25\% (126 -> 216)
\fill[popblue] (0,0) -- (126:2.6) arc (126:216:2.6) -- cycle;
% Habillement 20\% (216 -> 288)
\fill[poporange] (0,0) -- (216:2.6) arc (216:288:2.6) -- cycle;
% Objets recycl\'es 12\% (288 -> 331.2)
\fill[poppurple] (0,0) -- (288:2.6) arc (288:331.2:2.6) -- cycle;
% Autres 8\% (331.2 -> 360)
\fill[popgold] (0,0) -- (331.2:2.6) arc (331.2:360:2.6) -- cycle;
\draw[popdark] (0,0) circle (2.6);
% Etiquettes
\node[text=white,font=\bfseries] at (63:1.6) {35\,\%};
\node[text=white,font=\bfseries] at (171:1.6) {25\,\%};
\node[text=white,font=\bfseries] at (252:1.6) {20\,\%};
\node[text=white,font=\bfseries] at (309:1.7) {12\,\%};
\node[text=white,font=\bfseries] at (345:1.7) {8\,\%};
% L\'egende
\fill[popgreen] (3.3,2.2) rectangle (3.7,2.5); \node[anchor=west,text=popdark] at (3.8,2.35) {Produits agroalimentaires transform\'es};
\fill[popblue] (3.3,1.5) rectangle (3.7,1.8); \node[anchor=west,text=popdark] at (3.8,1.65) {Articles de d\'ecoration};
\fill[poporange] (3.3,0.8) rectangle (3.7,1.1); \node[anchor=west,text=popdark] at (3.8,0.95) {Habillement et accessoires};
\fill[poppurple] (3.3,0.1) rectangle (3.7,0.4); \node[anchor=west,text=popdark] at (3.8,0.25) {Articles en mat\'eriaux recycl\'es};
\fill[popgold] (3.3,-0.6) rectangle (3.7,-0.3); \node[anchor=west,text=popdark] at (3.8,-0.45) {Autres (poterie, vannerie\dots)};
\end{tikzpicture}
\legende{R\'epartition fictive mais r\'ealiste des produits artisanaux vendus sur un march\'e urbain camerounais (type march\'e Mokolo \`a Yaound\'e ou march\'e Mboppi \`a Douala), en pourcentage des \'etals observ\'es. Source : simulation p\'edagogique d'apr\`es des observations de terrain ; les statistiques officielles d\'etaill\'ees font d\'efaut, d'o\`u l'importance des enqu\^etes de terrain en TP.}
\end{center}
\end{popfigure}

\begin{methode}[Inventorier et classer des documents (photos, catalogues, \'etals de march\'e)]
Pour classer des produits artisanaux observ\'es sur des photographies, dans un catalogue ou lors d'une sortie au march\'e, proc\`ede en quatre \'etapes :
\begin{enumerate}
\item \textbf{D\'efinis d'abord des crit\`eres de classement} avant de classer : le domaine (d\'ecoration, habillement, agroalimentaire), la mati\`ere premi\`ere (bois, cuir, textile, produits agricoles, mat\'eriaux recycl\'es) ou la client\`ele vis\'ee (urbaine, touristique, diaspora, export).
\item \textbf{Observe et inventorie} chaque produit : note son nom, sa mati\`ere, son lieu de fabrication probable, son prix \'eventuel.
\item \textbf{Classe} chaque produit dans une rubrique en appliquant strictement tes crit\`eres ; un m\^eme produit peut relever de deux crit\`eres diff\'erents (un sac en raphia : habillement par le domaine, touristique par la client\`ele) — pr\'cise alors ton crit\`ere principal.
\item \textbf{Interpr\`ete} ton tableau : quelle rubrique domine ? Quelle client\`ele est la plus vis\'ee ? Qu'est-ce que cela r\'ev\`ele de l'\'evolution de l'artisanat ?
\end{enumerate}
Erreur \`a \'eviter : commencer \`a classer sans avoir fix\'e de crit\`eres — le classement devient alors incoh\'erent et impossible \`a interpr\'eter.
\end{methode}

\begin{exoresolu}[Classer et interpr\'eter des produits artisanaux]
\textbf{\'Enonc\'e.} Lors d'une visite au march\'e artisanal de Yaound\'e, un \'el\`eve rel\`eve les produits suivants : (a) chips de plantain emball\'ees ; (b) masque sculpt\'e bamoun ; (c) sac \`a main en pagne ndop ; (d) miel de l'Adamaoua en pot \'etiquet\'e ; (e) lampe en calebasse ; (f) chaussures en cuir. 1) Classe ces produits par domaine. 2) Indique, pour chacun, la client\`ele la plus probablement vis\'ee. 3) Que r\'ev\`ele cette diversit\'e de l'\'evolution de l'artisanat camerounais ?
\tcblower
\textbf{Solution d\'etaill\'ee.}
1) \emph{Classement par domaine} : agroalimentaire : (a) et (d) ; d\'ecoration : (b) et (e) ; habillement et accessoires : (c) et (f).

2) \emph{Client\`ele vis\'ee} : (a) chips de plantain : client\`ele urbaine (consommation quotidienne, prix accessible) ; (b) masque bamoun : client\`ele touristique et diaspora (objet d'art identitaire) ; (c) sac en pagne ndop : client\`ele urbaine jeune et diaspora (mode et identit\'e culturelle) ; (d) miel \'etiquet\'e : client\`ele urbaine et supermarch\'es, potentiellement l'export ; (e) lampe en calebasse : client\`ele urbaine et touristique (d\'ecoration) ; (f) chaussures en cuir : client\`ele urbaine.

3) \emph{Interpr\'etation} : cette diversit\'e r\'ev\`ele que l'artisanat camerounais ne se contente plus de reproduire les objets traditionnels : il cr\'ee des \cle{produits nouveaux} adapt\'es \`a des \cle{march\'es nouveaux} (ville, tourisme, diaspora, export). L'artisan est devenu un entrepreneur qui analyse la demande, transforme les produits agricoles locaux (plantain, miel) et valorise les identit\'es culturelles (ndop, art bamoun). Cette adaptation est une r\'eponse \`a la concurrence des produits industriels import\'es et contribue \`a l'int\'egration nationale en faisant circuler les produits de toutes les r\'egions (Adamaoua, Ouest, Nord-Ouest) vers les grandes villes du Centre et du Littoral.
\end{exoresolu}

\begin{exercice}[\`A toi de jouer]
Observe les \'etals d'un march\'e de ta localit\'e (ou un catalogue de produits artisanaux en ligne). 1) Inventorie dix produits artisanaux. 2) Classe-les selon deux crit\`eres : le domaine et la client\`ele vis\'ee. 3) Construis un tableau \`a double entr\'ee croisant ces deux crit\`eres. 4) R\'edige un court commentaire (une dizaine de lignes) : quelle client\`ele domine ? L'innovation artisanale est-elle visible ? Un corrig\'e type te sera propos\'e en classe.
\end{exercice}

\begin{aretenir}[L'essentiel de la section]
\begin{itemize}
\item Un \cle{produit nouveau} est un bien artisanal in\'edit ou relook\'e, destin\'e \`a une demande nouvelle (ville, tourisme, diaspora, export).
\item Trois grands domaines d'innovation : la \textbf{d\'ecoration} (meubles \`a motifs traditionnels, objets d'art, mat\'eriaux recycl\'es), l'\textbf{habillement} (pagnes modernis\'es, broderie, accessoires en cuir), l'\textbf{agroalimentaire} (jus, \'epices conditionn\'ees, farines, miel, chips).
\item L'innovation est la r\'eponse de l'artisan \`a la \textbf{concurrence des produits industriels import\'es} : il mise sur l'authenticit\'e, la personnalisation, la fra\^icheur et la proximit\'e.
\item Mais un produit nouveau n'est pas automatiquement un produit de qualit\'e : il doit respecter des \cle{normes} pour durer — c'est l'objet de la prochaine section.
\end{itemize}
\end{aretenir}

\coursec{L'amélioration de la qualité et le respect des normes}

Dans la section précédente, nous avons vu que l'artisanat camerounais produit désormais des \cle{produits nouveaux} : jus de fruits locaux embouteillés, savons parfumés, tissus teints modernes, meubles design. Mais fabriquer un produit nouveau ne suffit pas. Pour entrer dans les rayons de Santa Lucia à Douala, de Mahima à Yaoundé, ou pour être exporté vers l'Europe, ce produit doit encore franchir une étape décisive : prouver sa \cle{qualité} et sa conformité aux \cle{normes}. C'est toute la différence entre le jus de gingembre vendu dans une bouteille recyclée non étiquetée au bord de la route et le même jus, embouteillé, étiqueté et contrôlé, vendu jusqu'à trois fois plus cher en supermarché. Cette section explique cette révolution silencieuse de l'artisanat camerounais.

\courssub{Deux notions à ne pas confondre : la qualité et la norme}

\textbf{Situation concrète.} Au marché Mokolo de Yaoundé, deux vendeuses proposent du piment moulu. Le premier est vendu en vrac, dans un sac ouvert : il est excellent au goût, mais on ne sait ni quand il a été moulu, ni s'il contient des colorants ajoutés. Le second est conditionné en sachet scellé portant une étiquette complète, avec une date d'expiration et une mention de conformité. Lequel choisirait un supermarché ? Lequel choisirait un importateur européen ? Répondre à cette question suppose de distinguer deux notions que les élèves confondent souvent.

\begin{definition}[La qualité]
La \cle{qualité} est l'ensemble des caractéristiques d'un produit qui lui confèrent l'\textbf{aptitude à satisfaire les besoins exprimés ou implicites du consommateur}. Ces caractéristiques sont de plusieurs ordres :
\begin{itemize}
\item la \textbf{durabilité} (un meuble en bois d'iroko bien séché ne se fend pas au bout de six mois) ;
\item la \textbf{finition} (un pagne sans bavure de teinture, un bronze de Foumban sans défaut de moulage) ;
\item la \textbf{sécurité} (un jus sans germes pathogènes, un jouet en bois sans écharde ni peinture toxique) ;
\item le \textbf{goût} et l'aspect pour les produits alimentaires (miel limpide, café bien torréfié).
\end{itemize}
La qualité se juge donc du point de vue du \emph{client} : un produit de qualité est un produit qui donne satisfaction.
\end{definition}

\begin{definition}[La norme]
Une \cle{norme} est un document établi par \textbf{consensus} et approuvé par un organisme reconnu, qui fournit des \textbf{règles, des lignes directrices ou des caractéristiques} pour des activités ou leurs résultats. Au Cameroun, c'est l'\textbf{ANOR} (Agence des Normes et de la Qualité), créée par la loi n°2009/006 du 14 janvier 2009 et basée à Yaoundé, qui élabore les normes nationales (préfixe NC, « Norme Camerounaise ») et délivre la marque de conformité. Au niveau international, on cite l'\textbf{ISO} (Organisation internationale de normalisation, ex. ISO 9001 pour le management de la qualité, ISO 22000 pour la sécurité des denrées alimentaires) et le \textbf{Codex Alimentarius} pour les produits alimentaires.
\end{definition}

\begin{attention}[Erreur fréquente : confondre qualité et norme]
La \textbf{qualité} renvoie à la \emph{satisfaction du client} ; la \textbf{norme} renvoie à la \emph{conformité à une règle technique}. Ce ne sont pas des synonymes :
\begin{itemize}
\item un produit peut être \textbf{conforme à une norme mais de qualité médiocre} (un sachet de piment parfaitement étiqueté et hygiénique, mais fade car fabriqué avec un piment de mauvaise variété) ;
\item inversement, un produit peut être d'\textbf{excellente qualité gustative mais non conforme} (un fromage artisanal délicieux fabriqué dans des conditions d'hygiène non déclarées, donc interdit à la vente en supermarché).
\end{itemize}
Au Baccalauréat, rédiger « la qualité, c'est-à-dire la norme » est une erreur de définition sanctionnée. Présente toujours les deux notions séparément, puis montre leur complémentarité : la norme est un \emph{plancher} obligatoire, la qualité est un \emph{objectif} de satisfaction qui peut aller bien au-delà.
\end{attention}

\courssub{Les nouvelles méthodes de conditionnement}

Le \cle{conditionnement} désigne l'ensemble des opérations d'emballage qui protègent le produit entre l'atelier et le consommateur. C'est souvent la première modernisation visible de l'artisanat camerounais.

\textbf{Observation guidée.} Compare deux photographies mentales : celle d'un débit d'huile de palme vendue en bouteilles de boisson usagées, bouchées avec un bouchon de fortune, et celle d'une huile de palme artisanale de la région du Sud-Ouest conditionnée en bidons neufs scellés, avec capsule inviolable. Le contenu peut être identique ; la valeur marchande, la durée de conservation et l'accès aux circuits modernes sont radicalement différents.

Les nouvelles méthodes de conditionnement comprennent :
\begin{itemize}
\item les \textbf{emballages plastiques et en verre scellés} : bouteilles et pots neufs à capsule inviolable pour les jus, miels, huiles et liqueurs artisanales ; le verre est privilégié pour les produits haut de gamme car il préserve mieux les arômes ;
\item les \textbf{sachets étiquetés} : sachets plastiques thermosoudés pour les épices moulues (piment, njansang, pebè), le gari, la farine de manioc, les tisanes ; la soudure garantit que le produit n'a pas été ouvert ;
\item les \textbf{cartons imprimés} : emballages en carton portant le logo et les couleurs du producteur, utilisés pour les savons, les œuvres d'art, les chaussures en cuir de la Benoué ;
\item l'\textbf{emballage sous vide} : l'air est retiré du sachet avant soudure, ce qui ralentit l'oxydation et le développement des moisissures ; technique de plus en plus utilisée pour le poisson fumé artisanal de Youpwe et de Limbé, ainsi que pour les viandes séchées (kilichi).
\end{itemize}

\begin{propriete}[Fonctions du conditionnement moderne]
Un bon conditionnement remplit quatre fonctions : \textbf{protéger} le produit (contre l'humidité, la lumière, les chocs, les microbes), \textbf{conserver} (allonger la durée de vie), \textbf{informer} (support de l'étiquette) et \textbf{valoriser} (rendre le produit attractif et transportable vers des marchés lointains).
\end{propriete}

\courssub{Les nouvelles méthodes de présentation : étiquetage, branding, vitrines}

Le conditionnement protège ; la \cle{présentation} vend. Trois outils dominent.

\textbf{L'étiquetage.} L'étiquette est la carte d'identité du produit. Elle transforme un objet anonyme en produit identifiable, traçable et responsable.

\begin{exemplebox}[Les mentions obligatoires d'une étiquette alimentaire]
Un produit alimentaire artisanal étiqueté doit mentionner \textbf{au minimum} :
\begin{enumerate}
\item la \textbf{dénomination} du produit (ex. « Jus de gingembre ») ;
\item la \textbf{liste des ingrédients} par ordre décroissant de poids (gingembre, eau, sucre, citron) ;
\item le \textbf{poids net} ou le volume net (ex. 50 cl) ;
\item les \textbf{dates de fabrication et d'expiration} (ou la date limite de consommation, DLC) ;
\item le \textbf{nom et l'adresse du fabricant} (ex. « GIC Saveurs du Mbam, Bafia »).
\end{enumerate}
S'y ajoutent souvent le numéro de lot, les conditions de conservation (« à conserver au frais ») et, pour les produits certifiés, la marque de conformité à la norme.
\end{exemplebox}

\textbf{Le branding local.} De plus en plus d'artisans camerounais créent de véritables \textbf{marques} : nom, logo, slogan, charte de couleurs. Le poivre de Penja, le miel de l'Adamaoua (marque Oku pour le miel blanc des monts Oku), le café de l'Ouest sont devenus des noms qui évoquent une origine et une promesse de qualité. Ce branding local permet de sortir de la concurrence par les prix : on ne vend plus « du miel », on vend « le miel blanc d'Oku ».

\textbf{Les stands et vitrines.} La présentation physique évolue aussi : stands aménagés lors des foires (Foire internationale de l'artisanat, salons PROMOTE à Yaoundé), boutiques-vitrines dans les villes (quartier des artisans de Foumban, villages artisanaux de Yaoundé et Douala), et désormais vitrines numériques (pages Facebook, Instagram, plateformes de commerce électronique) où la photographie soignée du produit joue le rôle d'étalage.

\begin{popfigure}
\begin{center}
\begin{tikzpicture}[scale=1, transform shape]
% Fond de l'étiquette
\draw[fill=popcream, draw=popdark, line width=1.2pt, rounded corners=6pt] (0,0) rectangle (10.5,6.4);
\draw[draw=popgold, line width=0.8pt, rounded corners=4pt] (0.25,0.25) rectangle (10.25,6.15);
% Bandeau marque
\fill[popgreen] (0.25,4.9) rectangle (10.25,6.15);
\node[text=white, font=\bfseries\large] at (5.25,5.75) {SAVEURS DU MBAM};
\node[text=white, font=\small] at (5.25,5.2) {Produit artisanal du Cameroun --- GIC Saveurs du Mbam, Bafia (Région du Centre)};
% Dénomination
\node[font=\bfseries\large, text=popink] at (5.25,4.35) {JUS DE GINGEMBRE NATUREL};
% Mentions
\node[anchor=west, font=\small, text=popdark] at (0.7,3.7) {\textbf{Ingrédients :} gingembre frais, eau, sucre, jus de citron.};
\node[anchor=west, font=\small, text=popdark] at (0.7,3.15) {\textbf{Volume net :} 50 cl \quad \textbf{Lot :} L2024-117};
\node[anchor=west, font=\small, text=popdark] at (0.7,2.6) {\textbf{Fabriqué le :} 05/01/2025 \quad \textbf{À consommer avant le :} 05/04/2025};
\node[anchor=west, font=\small, text=popdark] at (0.7,2.05) {\textbf{Conservation :} à conserver au frais, à l'abri de la lumière.};
\node[anchor=west, font=\small, text=popdark] at (0.7,1.5) {\textbf{Fabricant :} GIC Saveurs du Mbam, B.P. 12 Bafia --- Tél. : 6XX XX XX XX};
% Logo conformité
\draw[fill=popblueL, draw=popblue, line width=1pt] (8.0,0.55) circle (0.62);
\node[font=\scriptsize\bfseries, text=popblue, align=center] at (8.0,0.55) {Conforme\\NC};
% Code-barres schématique
\foreach \x/\w in {0.9/0.06,1.05/0.03,1.18/0.06,1.32/0.03,1.45/0.09,1.62/0.03,1.75/0.06,1.9/0.03,2.03/0.09,2.2/0.03,2.33/0.06,2.48/0.03,2.6/0.06,2.75/0.03,2.88/0.09,3.05/0.03,3.18/0.06}
  \fill[popdark] (\x,0.45) rectangle (\x+\w,1.15);
\node[font=\scriptsize, text=popdark] at (1.95,0.32) {6 1 9 0 0 0 1 2 3 4 5 6};
% Flèches de légende
\draw[-{Stealth}, poporange, thick] (5.25,4.15) -- (5.25,3.95);
\node[anchor=west, font=\scriptsize, text=poporange] at (5.5,4.05) {1. Dénomination};
\draw[-{Stealth}, poporange, thick] (0.3,3.7) -- (0.65,3.7);
\node[anchor=east, font=\scriptsize, text=poporange] at (0.25,3.7) {2. Ingrédients + Volume};
\draw[-{Stealth}, poporange, thick] (0.3,2.6) -- (0.65,2.6);
\node[anchor=east, font=\scriptsize, text=poporange] at (0.25,2.6) {3. Dates};
\draw[-{Stealth}, poporange, thick] (8.65,0.9) -- (9.3,1.6);
\node[anchor=west, font=\scriptsize, text=poporange] at (8.6,1.75) {4. Marque de conformité NC};
\end{tikzpicture}
\end{center}
\legende{Maquette d'une étiquette normalisée d'un produit artisanal camerounais (jus de gingembre). Mentions obligatoires : 1. dénomination du produit ; 2. liste des ingrédients et poids/volume net ; 3. dates de fabrication et d'expiration ; 4. marque de conformité à la norme camerounaise (NC), délivrée par l'ANOR. Le code-barres facilite la vente en supermarché.}
\end{popfigure}

\courssub{Le contrôle de la qualité : de la matière première au produit fini}

La qualité ne se décrète pas sur l'étiquette : elle se \textbf{construit} à chaque étape de la production. C'est le rôle du \cle{contrôle de la qualité}, qui intervient à quatre niveaux.

\begin{enumerate}
\item \textbf{La vérification des matières premières.} Un bon jus exige des fruits sains et mûrs ; un bon fromage exige un lait frais non frelaté ; un bon meuble exige un bois correctement séché (un bois humide se fend et se déforme). L'artisan moderne sélectionne ses fournisseurs et refuse les lots défectueux.
\item \textbf{L'hygiène de production.} Locaux propres et aérés, eau potable, lavage des mains et des équipements, port de tenues (blouse, charlotte, gants), protection contre les insectes et les poussières. Pour les produits alimentaires, c'est la condition première de la sécurité sanitaire.
\item \textbf{Le contrôle en fin de chaîne.} Avant conditionnement, chaque lot est vérifié : goût, aspect, odeur, poids, solidité des soudures, exactitude des étiquettes. Les produits défectueux sont écartés.
\item \textbf{Les analyses en laboratoire.} Pour les produits destinés aux supermarchés ou à l'exportation, des échantillons sont envoyés dans des laboratoires accrédités (laboratoire de l'ANOR, laboratoires universitaires comme ceux de la FSMB ou de l'ENSP à Yaoundé, laboratoires privés spécialisés en agroalimentaire) qui mesurent par exemple la teneur en eau, la présence de germes, de pesticides ou de colorants interdits.
\end{enumerate}

\begin{popfigure}
\begin{center}
\begin{tikzpicture}[
  etape/.style={rectangle, rounded corners=4pt, draw=popdark, line width=0.9pt, fill=popblueL, text width=2.5cm, minimum height=1.5cm, align=center, font=\small},
  ctrl/.style={rectangle, rounded corners=3pt, draw=poporange, line width=0.9pt, fill=poporangeL, text width=2.5cm, align=center, font=\scriptsize},
  fleche/.style={-{Stealth[length=3mm]}, line width=1.4pt, popgreen}
]
\node[etape, align=center] (a) at (0,0){\textbf{Matière première}\\fruits, bois, laine, cuir};
\node[etape, right=1.1cm of a, align=center] (b){\textbf{Production hygiénique}\\locaux propres, tenues, eau potable};
\node[etape, right=1.1cm of b, align=center] (c){\textbf{Contrôle qualité}\\goût, aspect, poids, analyses};
\node[etape, right=1.1cm of c, align=center] (d){\textbf{Conditionnement}\\sachets scellés, verre, sous vide};
\node[etape, below=1.4cm of d, align=center] (e){\textbf{Étiquetage}\\dénomination, dates, fabricant};
\node[etape, left=1.1cm of e, align=center] (f){\textbf{Mise en marché}\\supermarchés, foires, exportation};
\draw[fleche] (a) -- (b);
\draw[fleche] (b) -- (c);
\draw[fleche] (c) -- (d);
\draw[fleche] (d) -- (e);
\draw[fleche] (e) -- (f);
% Contrôles intermédiaires
\node[ctrl, above=0.5cm of a, xshift=1.9cm] (c1) {Contrôle 1 : sélection des matières premières};
\draw[-{Stealth}, poporange, dashed] (c1) -- (a.north east);
\node[ctrl, above=0.5cm of c, xshift=1.9cm] (c2) {Contrôle 2 : hygiène de production};
\draw[-{Stealth}, poporange, dashed] (c2) -- (b.north east);
\node[ctrl, below=0.5cm of f, xshift=-1.9cm] (c3) {Contrôle 3 : fin de chaîne + laboratoire};
\draw[-{Stealth}, poporange, dashed] (c3) -- (c.south west);
\end{tikzpicture}
\end{center}
\legende{Schéma de la chaîne qualité dans un atelier artisanal modernisé. Les étapes de production (rectangles bleus) sont doublées de contrôles (encadrés orange) : vérification des matières premières, hygiène de production, contrôle en fin de chaîne complété par des analyses de laboratoire. Chaque maillon défaillant compromet la qualité finale.}
\end{popfigure}

\begin{methode}[Analyser la qualité d'un produit artisanal (démarche d'observation)]
Pour évaluer concrètement la qualité d'un produit artisanal (par exemple lors d'une visite d'atelier ou d'une foire), procède en quatre étapes :
\begin{enumerate}
\item \textbf{Observer} le produit : aspect, finition, odeur, régularité de la fabrication.
\item \textbf{Lire} l'étiquette : les mentions obligatoires sont-elles présentes et lisibles ?
\item \textbf{Questionner} le producteur : origine des matières premières, méthodes d'hygiène, existence de contrôles.
\item \textbf{Conclure} en distinguant les deux plans : le produit est-il \emph{conforme} (normes respectées) ? est-il de \emph{bonne qualité} (satisfaction probable du consommateur) ?
\end{enumerate}
\end{methode}

\courssub{Les enjeux : supermarchés, exportation, protection du consommateur, image du pays}

Pourquoi cet effort de qualité et de normalisation est-il si important ? Quatre enjeux majeurs peuvent être dégagés.

\textbf{L'accès aux supermarchés.} Les grandes et moyennes surfaces (Santa Lucia, Mahima, Carrefour Market à Douala et Yaoundé) n'acceptent que des produits conditionnés, étiquetés et traçables, avec code-barres et dates de conservation. La normalisation est donc le \emph{sésame} qui fait passer l'artisanat du marché de rue au circuit de distribution moderne, beaucoup plus rémunérateur.

\textbf{L'accès à l'exportation.} Les marchés européen, américain et même régionaux (CEMAC) exigent le respect de normes sanitaires et techniques strictes (normes ISO, règlements européens, normes du Codex Alimentarius). Un produit non conforme est refoulé à la frontière. À l'inverse, un produit certifié peut être exporté à forte valeur ajoutée.

\begin{savaistu}[Le poivre de Penja, première IGP d'Afrique subsaharienne]
Le \textbf{poivre de Penja}, cultivé dans le département du Moungo (région du Littoral) sur des sols volcaniques d'exception, a obtenu en 2014 une \textbf{Indication Géographique Protégée (IGP)} délivrée par l'Organisation Africaine de la Propriété Intellectuelle (OAPI) : c'est la \textbf{première IGP d'Afrique subsaharienne}. Ce label garantit l'origine et la qualité du produit, interdit l'usage du nom « poivre de Penja » pour un poivre produit ailleurs, et permet d'exporter vers l'Europe à des prix très supérieurs (le kilogramme peut se vendre plusieurs dizaines de milliers de FCFA chez les grands distributeurs européens, contre quelques milliers pour un poivre ordinaire). Il montre qu'un label de qualité peut transformer un produit artisanal en produit de luxe.
\end{savaistu}

\textbf{La protection du consommateur.} Le respect des normes sanitaires protège la santé des Camerounais : jus sans contamination, huiles sans adultération, cosmétiques artisanaux sans substances dangereuses (certains savons éclaircissants artisanaux contiennent de l'hydroquinone ou des corticoïdes interdits, d'où l'importance du contrôle). La norme est ici un bouclier sanitaire.

\textbf{L'image du « made in Cameroon ».} Des produits artisanaux bien conditionnés, étiquetés et fiables construisent une réputation nationale. À l'inverse, quelques produits défectueux ou frauduleux discréditent l'ensemble de la production locale et renforcent le préjugé en faveur des produits importés. La qualité est donc aussi un enjeu d'\textbf{intégration nationale} et de fierté économique : consommer local devient possible lorsque le produit local inspire confiance.

\begin{exoresolu}[Commentaire de document : deux conditionnements, deux destins commerciaux]
\textbf{Énoncé.} Document : « À Bafoussam, deux coopératives transforment le même lait de vache en fromage. La coopérative A vend ses fromages emballés dans des feuilles de bananier, sans étiquette, au marché local : prix 1 000 FCFA le fromage. La coopérative B a investi dans une salle de production hygiénique, conditionne sous vide, étiquette ses produits (dénomination, ingrédients, dates, fabricant) et a obtenu une attestation de conformité : ses fromages sont vendus 3 500 FCFA dans les supermarchés de Yaoundé et de Douala. »
Question 1 : Identifie dans le document les éléments qui relèvent de la qualité et ceux qui relèvent de la norme.
Question 2 : Explique l'écart de prix observé.
\tcblower
\textbf{Solution détaillée.}
\textbf{Question 1.} Relèvent de la \textbf{qualité} (satisfaction du consommateur) : le goût et la fraîcheur du fromage (identiques à la base puisqu'il s'agit du même lait), la conservation prolongée permise par le vide, la présentation soignée. Relèvent de la \textbf{norme} (conformité à une règle) : l'hygiène de la salle de production, les mentions obligatoires de l'étiquette (dénomination, ingrédients, dates de fabrication et d'expiration, nom du fabricant), l'attestation de conformité délivrée par l'organisme compétent (ANOR).
\textbf{Question 2.} L'écart de prix (de 1 000 à 3 500 FCFA, soit 3,5 fois plus) s'explique par trois mécanismes : (a) le conditionnement sous vide et l'étiquetage ont un coût, mais ils allongent la durée de conservation et permettent le transport vers les grandes villes ; (b) la conformité aux normes ouvre l'accès aux supermarchés, circuit de distribution où la clientèle solvable accepte des prix élevés ; (c) la confiance créée par l'étiquette et l'attestation agit comme une garantie : le consommateur paie non seulement le fromage, mais aussi la sécurité et la traçabilité. La modernisation de la qualité transforme donc un produit artisanal local en produit commercialisé sur un marché national rémunérateur.
\end{exoresolu}

\begin{attention}[Pièges de rédaction au Baccalauréat]
\begin{itemize}
\item Ne réduis pas la qualité à l'emballage : un bel emballage peut cacher un mauvais produit. La qualité porte d'abord sur le produit lui-même (durabilité, finition, sécurité, goût).
\item Cite des \textbf{exemples camerounais précis} (ANOR, poivre de Penja IGP, miel d'Oku) plutôt que des généralités : c'est ce qui distingue une bonne copie.
\item N'affirme pas que « tous les produits artisanaux sont désormais normalisés » : la normalisation reste inégale, coûteuse pour les petits artisans, et une grande partie de la production échappe encore au contrôle. Nuancer est une compétence.
\end{itemize}
\end{attention}

\begin{aretenir}[L'essentiel de la section]
\begin{itemize}
\item La \cle{qualité} est l'aptitude d'un produit à satisfaire le consommateur (durabilité, finition, sécurité, goût) ; la \cle{norme} est une règle technique ou sanitaire établie par un organisme (ANOR au Cameroun, ISO au niveau international).
\item Les nouvelles méthodes de \textbf{conditionnement} (sachets scellés, verre, cartons imprimés, emballage sous vide) et de \textbf{présentation} (étiquetage complet, branding local, stands et vitrines) valorisent les produits artisanaux.
\item Le \textbf{contrôle de la qualité} s'exerce à tous les stades : matières premières, hygiène de production, fin de chaîne, analyses de laboratoire.
\item Les enjeux sont majeurs : accès aux supermarchés et à l'exportation, protection du consommateur, rayonnement du « made in Cameroon » — illustré par le succès de l'IGP du poivre de Penja.
\end{itemize}
\end{aretenir}

\begin{exercice}[Pour t'entraîner]
1. Définis les termes \cle{qualité} et \cle{norme} en une phrase chacun, puis construis un exemple montrant qu'un produit peut être conforme à une norme sans être de bonne qualité.
2. Énumère les cinq mentions minimales d'une étiquette de produit alimentaire artisanal.
3. Construis un tableau à deux colonnes présentant, pour chaque étape de la chaîne qualité, le contrôle correspondant.
4. Rédige un paragraphe argumenté (10 à 15 lignes) montrant comment le respect des normes peut contribuer à l'intégration nationale du Cameroun.
\end{exercice}

\coursec{L'organisation des filières et des circuits de vente}

Dans la section précédente, nous avons vu comment l'artisan camerounais améliore la \cle{qualité} de ses produits et respecte les \cle{normes} de fabrication : un meuble bien poncé, un jus de fruit correctement conditionné et étiqueté, un tissu teint avec des colorants fixés. Mais un beau produit, même conforme aux normes, ne rapporte rien s'il reste empilé au fond de l'atelier. Encore faut-il qu'il trouve son acheteur. Prenons une situation concrète : à Foumban, dans la région de l'Ouest, un sculpteur a passé trois semaines à tailler un masque en bois d'ébène. Qui lui a fourni le bois ? Qui va polir et teinter le masque ? Où sera-t-il vendu : au bord de la route, dans une boutique de la rue des Artisans, à la foire de Yaoundé, ou à un touriste européen ? Répondre à ces questions, c'est décrire une \cle{filière} et ses \cle{circuits de vente}. C'est l'objet de cette section.

\courssub{La notion de filière artisanale}

\textbf{De l'intuition à la définition.} Observe le parcours d'un objet artisanal, par exemple un panier tressé vendu sur le marché de Bafoussam. Avant d'arriver entre les mains de la cliente, ce panier a connu plusieurs étapes : la récolte du raphia ou de la paille dans la brousse, le séchage et la teinture des fibres, le tressage proprement dit, le transport jusqu'au marché, puis la vente. Chaque étape est assurée par un acteur différent, et chaque acteur ajoute de la valeur au produit. L'ensemble de ces étapes enchaînées forme une filière.

\begin{definition}[Filière]
Une \cle{filière} est la succession des activités économiques qui vont de la matière première au produit fini vendu au consommateur. Dans une filière artisanale, on distingue classiquement quatre maillons : l'\cle{approvisionnement} en matière première (bois, cuir, argile, fibres, métal), la \cle{fabrication} (transformation de la matière en produit), le \cle{conditionnement} (emballage, étiquetage, présentation) et la \cle{distribution} (vente au consommateur final). Chaque maillon est tenu par des acteurs qui se transmettent le produit en y ajoutant de la valeur.
\end{definition}

\begin{definition}[Circuit de vente]
Un \cle{circuit de vente} (ou circuit de distribution) est le cheminement d'un produit du producteur au consommateur. Il est dit \cle{court} lorsque la vente se fait directement ou avec un seul intermédiaire (artisan $\rightarrow$ consommateur), et \cle{long} lorsque plusieurs intermédiaires se succèdent (artisan $\rightarrow$ grossiste $\rightarrow$ détaillant $\rightarrow$ consommateur).
\end{definition}

\begin{attention}[Ne pas réduire la filière à la seule vente]
L'erreur la plus fréquente au Baccalauréat consiste à décrire uniquement les circuits de vente en oubliant l'amont de la filière. Or l'approvisionnement en matière première de qualité est un maillon tout aussi décisif : un sculpteur de Foumban qui ne trouve plus de bois d'ébène de bonne qualité, ou un potier de Bafmeng qui reçoit une argile impure, ne pourra pas produire, quels que soient ses débouchés commerciaux. Dans une copie, décris toujours la filière \emph{en entier}, de la matière première au consommateur.
\end{attention}

\begin{methode}[Construire un croquis de filière]
Pour représenter une filière artisanale sur un croquis :
\begin{enumerate}
\item Identifier tous les maillons, de la matière première au consommateur final, sans en oublier (approvisionnement, fabrication, finition, conditionnement, vente).
\item Représenter chaque maillon par un cadre rectangulaire contenant un nom d'acteur et un lieu précis (forêt de l'Ouest, atelier de Foumban, foire de Yaoundé).
\item Relier les cadres par des flèches orientées dans le sens du produit, de l'amont vers l'aval.
\item Annoter le croquis : nature du flux (bois brut, produit fini, argent), acteurs, lieux, et éventuellement les bifurcations (client local, touriste, exportateur).
\item Ajouter un titre, une légende et, si le croquis a une dimension spatiale, une orientation.
\end{enumerate}
\end{methode}

\begin{popfigure}
\begin{center}
\begin{tikzpicture}[
  maillon/.style={draw=popblue, fill=popblueL, rounded corners=2pt, minimum width=2.6cm, minimum height=1.1cm, align=center, font=\small, text=popdark},
  flux/.style={-{Stealth[length=3mm]}, very thick, poporange},
  note/.style={font=\scriptsize\itshape, text=popmuted}
]
\node[maillon, align=center] (foret) at (0,0){\textbf{Forêt (Ouest)}\\ coupe du bois\\ (ébène, iroko)};
\node[maillon, align=center] (sculpt) at (4.2,0){\textbf{Sculpteur}\\ Foumban\\ taille du masque};
\node[maillon, align=center] (finition) at (8.4,0){\textbf{Atelier de finition}\\ ponçage, teinture,\\ cirage};
\node[maillon, align=center] (vente) at (12.6,0){\textbf{Boutique / foire}\\ rue des Artisans,\\ Yaoundé, Douala};
\node[maillon, fill=popgreenL, draw=popgreen] (local) at (9.4,-2.6) {Client local};
\node[maillon, fill=popgreenL, draw=popgreen] (touriste) at (12.6,-2.6) {Touriste};
\node[maillon, fill=popgreenL, draw=popgreen, align=center] (export) at (15.8,-2.6){Exportateur\\ (diaspora, Europe)};
\draw[flux] (foret) -- (sculpt);
\draw[flux] (sculpt) -- (finition);
\draw[flux] (finition) -- (vente);
\draw[flux] (vente.south) -- (local.north);
\draw[flux] (vente.south) -- (touriste.north);
\draw[flux] (vente.east) .. controls (15.5,0.3) .. (export.north);
\node[note] at (2.1,0.55) {bois brut};
\node[note] at (6.3,0.55) {objet taillé};
\node[note] at (10.5,0.55) {produit fini};
\end{tikzpicture}
\end{center}
\legende{Croquis de la filière « sculpture de Foumban » (région de l'Ouest). Les cadres bleus représentent les maillons de la filière (approvisionnement, fabrication, finition, vente) ; les cadres verts représentent les débouchés finaux. Les flèches orange indiquent le sens du flux du produit, de la matière première vers le consommateur.}
\end{popfigure}

\courssub{L'organisation des artisans : se regrouper pour exister}

L'artisan camerounais travaille rarement seul face au marché. Pour acheter la matière première moins cher, partager des machines coûteuses, se former et négocier avec les pouvoirs publics, les artisans se regroupent en organisations professionnelles.

\begin{itemize}
\item Les \cle{coopératives} : les artisans mettent en commun leurs moyens (achat groupé de bois, de cuir ou de tissu ; atelier partagé ; magasin de vente collectif) et se partagent les bénéfices. Les coopératives de potières de l'Extrême-Nord ou de tisserands de l'Adamaoua en sont des exemples.
\item Les \cle{groupements d'initiative commune} (GIC) : forme d'association souple et très répandue au Cameroun, particulièrement en milieu rural. Un GIC de transformatrices de manioc ou de couturières peut ouvrir un compte bancaire, recevoir des dons et contracter des commandes groupées.
\item Les \cle{associations de quartier} : à l'échelle d'un village ou d'un quartier urbain, elles organisent la solidarité entre artisans (tontines, entraide, défense des prix).
\item La \cle{Chambre des Métiers du Cameroun (CMC)} et ses démembrements départementaux : structure de représentation professionnelle qui défend les intérêts des artisans, encadre l'apprentissage et le compagnonnage, et sert d'interlocuteur privilégié avec l'État. Elle joue un rôle clé dans la formation des apprentis et la délivrance des qualifications professionnelles.
\end{itemize}

Ces organisations sont le ciment de la filière : elles transforment des milliers d'artisans isolés en acteurs capables de peser sur les prix, d'accéder au crédit et de répondre à de grosses commandes (ameublement d'un hôtel, fourniture de tenues pour une cérémonie officielle).

\courssub{Les circuits de vente traditionnels}

Pendant des générations, l'artisanat camerounais s'est écoulé par des circuits simples et de proximité, qui restent aujourd'hui majoritaires.

\begin{itemize}
\item Les \cle{marchés locaux} : lieux de vente par excellence, comme les marchés de Foumban, de Bafoussam, de Maroua ou le marché des Fleurs de Douala. L'artisan ou un revendeur y tient un étal, souvent les jours de marché hebdomadaires.
\item La \cle{vente à domicile} : l'artisan vend directement depuis sa cour ou sa concession, à des clients du quartier qui connaissent sa réputation de bouche à oreille.
\item La \cle{boutique d'atelier} : l'artisan expose et vend ses produits sur son lieu de fabrication. C'est le modèle de la rue des Artisans de Foumban, où le touriste voit le sculpteur travailler et achète l'objet à la sortie de l'atelier.
\item Les \cle{foires artisanales} : événements périodiques où les artisans de plusieurs régions exposent leurs produits. Elles permettent de toucher une clientèle urbaine et fortunée et de prendre des commandes.
\end{itemize}

Ces circuits ont un point commun : la proximité entre le producteur et l'acheteur, qui permet la négociation directe du prix et limite le nombre d'intermédiaires.

\courssub{Les circuits modernes : conquérir de nouveaux marchés}

Depuis quelques décennies, et surtout depuis les années 2010, de nouveaux circuits de vente se sont ouverts aux artisans camerounais.

\begin{itemize}
\item Les \cle{supermarchés} : certaines enseignes de Douala et de Yaoundé (ex. : Casino, Mahima, Spar) consacrent des rayons aux produits artisanaux transformés (jus locaux, épices conditionnées, miel, savons artisanaux). L'accès à ces rayons exige un conditionnement soigné et des volumes réguliers, d'où l'importance des normes étudiées dans la section précédente.
\item Les \cle{boutiques spécialisées} : magasins de décoration, galeries d'art, boutiques de prêt-à-porter en pagne, qui valorisent le produit artisanal auprès d'une clientèle urbaine aisée.
\item Les \cle{salons et expositions} : la Foire internationale de Yaoundé (FIYA), le Salon de l'artisanat de Douala, ainsi que les salons internationaux de l'artisanat (SIAO à Ouagadougou), offrent une vitrine exceptionnelle et des contacts avec des acheteurs étrangers.
\item Le \cle{commerce électronique et les réseaux sociaux} : l'artisan photographie ses produits, les publie, reçoit les commandes par messagerie et expédie par les agences de transport.
\item L'\cle{exportation vers la diaspora} : les Camerounais vivant en Europe, en Amérique du Nord ou ailleurs en Afrique constituent un marché fidèle pour les tissus, les sculptures, les épices et les produits alimentaires transformés, expédiés par fret aérien ou maritime.
\end{itemize}

\begin{savaistu}[L'artisanat à l'ère de WhatsApp]
De nombreux artisans camerounais vendent désormais via WhatsApp et Facebook : ils photographient leurs produits avec un simple téléphone, publient dans des groupes ou sur leur statut, négocient par message, se font payer par mobile money (\textsc{Orange Money}, \textsc{MTN Mobile Money}) et livrent par les agences de transport voyageurs qui relient toutes les villes du pays. Ce circuit de vente numérique, en pleine croissance, supprime les intermédiaires et permet à une couturière de Bafoussam de vendre à une cliente de Douala sans quitter son atelier.
\end{savaistu}

\begin{popfigure}
\begin{center}
\begin{tikzpicture}[
  acteur/.style={draw=poppurple, fill=poppurpleL, rounded corners=2pt, minimum width=2.5cm, minimum height=0.9cm, align=center, font=\small, text=popdark},
  flux/.style={-{Stealth[length=3mm]}, very thick, poppink},
  titre/.style={font=\small\bfseries, text=popdark}
]
\node[titre] at (2.5,1.6) {Circuit court};
\node[acteur] (a1) at (0,0.6) {Artisan};
\node[acteur] (c1) at (5,0.6) {Consommateur};
\draw[flux] (a1) -- (c1) node[midway, above, font=\scriptsize\itshape, text=popmuted] {vente directe (atelier, marché, WhatsApp)};
\node[titre] at (2.5,-1.2) {Circuit long};
\node[acteur] (a2) at (0,-2.2) {Artisan};
\node[acteur] (g2) at (3.4,-2.2) {Grossiste};
\node[acteur] (d2) at (6.8,-2.2) {Détaillant};
\node[acteur] (c2) at (10.2,-2.2) {Consommateur};
\draw[flux] (a2) -- (g2);
\draw[flux] (g2) -- (d2);
\draw[flux] (d2) -- (c2);
\end{tikzpicture}
\end{center}
\legende{Comparaison des circuits de vente. En circuit court, l'artisan vend directement au consommateur et conserve l'essentiel de la marge. En circuit long, grossiste et détaillant prélèvent chacun une marge, ce qui réduit la part revenant à l'artisan mais élargit la diffusion du produit.}
\end{popfigure}

\begin{exemplebox}[Circuit court, circuit long : qui gagne quoi ?]
Un sac en cuir fabriqué à Maroua est vendu \SI{15000}{FCFA} en vente directe à l'atelier : l'artisan conserve la quasi-totalité de la marge. Le même sac, vendu \SI{30000}{FCFA} dans une boutique de Douala après passage par un grossiste, ne rapporte parfois que \SI{12000}{FCFA} à \SI{14000}{FCFA} à l'artisan : les intermédiaires captent plus de la moitié du prix final. En contrepartie, le circuit long donne accès à une clientèle que l'artisan seul n'aurait jamais atteinte. Le choix du circuit est donc un arbitrage entre marge unitaire et volume des ventes.
\end{exemplebox}

\begin{exoresolu}[Calculer la part des intermédiaires]
Énoncé. Une sculpture achetée \SI{20000}{FCFA} au sculpteur de Foumban est revendue \SI{45000}{FCFA} à un touriste par une boutique de Yaoundé, qui l'avait elle-même acquise \SI{30000}{FCFA} auprès d'un grossiste. 1) Calcule la part du prix final revenant à l'artisan. 2) Calcule la part totale captée par les intermédiaires. 3) Commente.
\tcblower
Solution détaillée.
\begin{enumerate}
\item Part de l'artisan : $\dfrac{20000}{45000} \times 100 \approx \num{44,4}\%$ du prix final.
\item Marge du grossiste : $30000 - 20000 = \SI{10000}{FCFA}$ ; marge de la boutique : $45000 - 30000 = \SI{15000}{FCFA}$. Total des intermédiaires : $10000 + 15000 = \SI{25000}{FCFA}$, soit $\dfrac{25000}{45000} \times 100 \approx \num{55,6}\%$ du prix final.
\item Commentaire : les intermédiaires captent plus de la moitié du prix final (\num{55,6}\%), alors que l'artisan, qui a fourni l'essentiel du travail, n'en reçoit que \num{44,4}\%. Cela illustre l'intérêt des circuits courts (vente directe à l'atelier, foires, vente en ligne) pour améliorer les revenus des artisans, sans nier le rôle de diffusion joué par grossistes et détaillants.
\end{enumerate}
\end{exoresolu}

\courssub{Le rôle de l'État et des partenaires}

L'État camerounais accompagne l'organisation des filières artisanales à plusieurs niveaux.

\begin{itemize}
\item Le \cle{Ministère des Petites et Moyennes Entreprises, de l'Économie Sociale et de l'Artisanat (MINPMEESA)} élabore la politique nationale de l'artisanat, recense les artisans, délivre les cartes professionnelles et organise les grandes manifestations de promotion.
\item Les \cle{centres de formation} (centres d'artisanat, lycées techniques, centres de formation professionnelle) transmettent les techniques modernes : utilisation des machines, gestion d'un atelier, techniques de conditionnement et de marketing.
\item Les \cle{appuis financiers} : microcrédits (via le \textsc{Fonds National de l'Emploi}, les programmes \textsc{PADMIR}, \textsc{PIAASI}), subventions et dons d'équipements accordés par l'État, les collectivités territoriales décentralisées (CTD), les ONG et les partenaires internationaux (coopérations, programmes de développement) pour moderniser les ateliers et structurer les GIC et coopératives.
\item La promotion du \cle{Made in Cameroon} : campagnes encourageant les Camerounais à consommer local (pagne, meubles, produits alimentaires transformés), labellisation des produits et participation aux foires internationales pour valoriser l'artisanat national.
\end{itemize}

\courssub{Les difficultés persistantes}

Malgré ces progrès, les filières artisanales camerounaises restent fragiles.

\begin{itemize}
\item Le \cle{financement} : les banques commerciales prêtent difficilement aux artisans sans garanties réelles ; les tontines et institutions de microfinance ne suffisent pas toujours à acheter des machines coûteuses (scies à ruban, tours à bois, fours à céramique).
\item La \cle{concurrence des importations} : les produits manufacturés importés (meubles en kit, chaussures en plastique, tissus imprimés industriels, objets de décoration à bas prix) concurrencent durement la production locale, souvent plus chère du fait du coût de la main-d'œuvre et de l'énergie.
\item La \cle{faiblesse de l'emballage} : un conditionnement fragile ou peu attractif ferme les portes des supermarchés et de l'exportation, même pour un produit de bonne qualité intrinsèque.
\item L'\cle{accès limité aux marchés internationaux} : coûts élevés du transport aérien et maritime, exigences strictes des normes sanitaires et phytosanitaires étrangères (ex. : normes \textsc{UE} pour le bois, l'alimentaire), méconnaissance des démarches douanières et d'exportation, et faible organisation collective pour mutualiser les coûts logistiques pénalisent les artisans qui visent la clientèle extérieure.
\end{itemize}

\begin{aretenir}[L'essentiel de la section]
\begin{itemize}
\item Une filière artisanale enchaîne quatre maillons : approvisionnement, fabrication, conditionnement, distribution.
\item Les artisans s'organisent en coopératives, GIC, associations de quartier et au sein de la Chambre des Métiers du Cameroun (CMC) pour peser sur le marché.
\item Les circuits traditionnels (marchés, vente à domicile, boutique d'atelier, foires) coexistent avec les circuits modernes (supermarchés, salons, commerce électronique, exportation vers la diaspora).
\item Le circuit court maximise la marge de l'artisan ; le circuit long élargit la diffusion mais les intermédiaires peuvent capter plus de la moitié du prix final.
\item L'État soutient les filières via le MINPMEESA (formation, financement, promotion du Made in Cameroon), mais le financement, la concurrence importée, l'emballage et l'accès aux marchés internationaux restent des obstacles majeurs.
\end{itemize}
\end{aretenir}

\begin{exercice}[Pour t'entraîner]
Une coopérative de potières de l'Extrême-Nord (région de Maroua) vend ses jarres de trois façons : au marché local (vente directe), à un grossiste de Maroua qui alimente les boutiques de Garoua, et par commandes WhatsApp livrées à Yaoundé via une agence de transport.
\begin{enumerate}
\item Identifie, pour chaque mode de vente, le type de circuit (court ou long).
\item Construis le croquis de cette filière, de l'extraction de l'argile jusqu'aux clients, en appliquant la méthode étudiée (acteurs, lieux, flèches, annotations).
\item Explique deux avantages que la coopérative retire de la vente par WhatsApp et deux difficultés qu'elle peut rencontrer pour exporter ses jarres vers l'Europe.
\end{enumerate}
\end{exercice}

\begin{corrige}[Exercice 1]
\begin{enumerate}
\item \textbf{Analyse des circuits :}
      \begin{itemize}
      \item Marché local (vente directe) : \textbf{circuit court} (Artisan $\rightarrow$ Consommateur).
      \item Grossiste de Maroua $\rightarrow$ Boutiques de Garoua : \textbf{circuit long} (Artisan $\rightarrow$ Grossiste $\rightarrow$ Détaillant $\rightarrow$ Consommateur).
      \item Commandes WhatsApp $\rightarrow$ Livraison Yaoundé : \textbf{circuit court} (Artisan $\rightarrow$ Consommateur), le transporteur n'étant qu'un prestataire logistique et non un intermédiaire commercial qui achète/revend.
      \end{itemize}
\item \textbf{Croquis de la filière « Poterie de l'Extrême-Nord » (schéma descriptif) :}
      \begin{center}
      \begin{tikzpicture}[
        maillon/.style={draw=popblue, fill=popblueL, rounded corners=2pt, minimum width=2.8cm, minimum height=1cm, align=center, font=\small, text=popdark},
        debouche/.style={draw=popgreen, fill=popgreenL, rounded corners=2pt, minimum width=2.5cm, minimum height=0.9cm, align=center, font=\small, text=popdark},
        flux/.style={-{Stealth[length=3mm]}, very thick, poporange},
        note/.style={font=\scriptsize\itshape, text=popmuted}
      ]
      \node[maillon, align=center] (argile) at (0,0){\textbf{Carrière d'argile}\\ Maroua / Mokolo\\ extraction};
      \node[maillon, align=center] (prepa) at (4,0){\textbf{Préparation}\\ Coopérative\\ malaxage, dégraissage};
      \node[maillon, align=center] (fab) at (8,0){\textbf{Fabrication}\\ Atelier coopératif\\ modelage, séchage, cuisson};
      \node[maillon, align=center] (cond) at (12,0){\textbf{Conditionnement}\\ emballage, étiquetage\\ « Made in Cameroon »};
      \node[maillon, align=center] (vente) at (16,0){\textbf{Point de vente}\\ Local coopératif / Maroua};
      \node[debouche, align=center] (marche) at (14,-2.5){Client local\\ (Marché Maroua)};
      \node[debouche, align=center] (garoua) at (17,-2.5){Client Garoua\\ (Boutique détaillant)};
      \node[debouche, align=center] (yaounde) at (20,-2.5){Client Yaoundé\\ (WhatsApp)};
      \draw[flux] (argile) -- (prepa) node[midway, above] {argile brute};
      \draw[flux] (prepa) -- (fab) node[midway, above] {argile prête};
      \draw[flux] (fab) -- (cond) node[midway, above] {jarres cuites};
      \draw[flux] (cond) -- (vente) node[midway, above] {jarres finies};
      \draw[flux] (vente.south west) -- (marche.north) node[midway, left] {circuit court};
      \draw[flux] (vente.south) -- (garoua.north) node[midway, right] {via grossiste};
      \draw[flux] (vente.south east) -- (yaounde.north) node[midway, right] {transport direct};
      \end{tikzpicture}
      \end{center}
      \textit{(Le croquis doit comporter un titre, une légende pour les couleurs bleu/vert, et l'orientation Nord si spatialisé).}
\item \textbf{Avantages du WhatsApp :}
      \begin{itemize}
      \item Suppression des intermédiaires commerciaux : la coopérative fixe son prix et conserve \num{100}\% de la marge commerciale (hors frais de port).
      \item Élargissement du bassin de clientèle sans déplacement : accès direct aux clients des grandes métropoles (Yaoundé, Douala) et de la diaspora, commande à l'unité possible.
      \end{itemize}
      \textbf{Difficultés pour l'export vers l'Europe :}
      \begin{itemize}
      \item \textbf{Logistique et coûts :} Fragilité des jarres (casse), volume/poids élevé $\rightarrow$ fret aérien prohibitif, fret maritime long (groupage difficile pour petits volumes). Coût d'emballage spécifique (caisses bois, calage) renchérissant le prix final.
      \item \textbf{Normes et visibilité :} Exigences sanitaires (plomb dans les émaux ?), étiquetage \textsc{UE} (traçabilité, composition), absence de label reconnu, méconnaissance des démarches douanières (code NC, certificat d'origine EUR.1), concurrence de la céramique industrielle européenne bon marché.
      \end{itemize}
\end{enumerate}
\end{corrige}

\coursec{Travaux pratiques : cartographier et schématiser l'artisanat camerounais}

Dans la section précédente, nous avons vu comment les artisans camerounais s'organisent en \cle{filières} et en \cle{circuits de vente} pour écouler leurs produits, du marché de Foumban aux boutiques de Douala. Ces réalités géographiques — des lieux, des flux, des réseaux — ne se racontent pas seulement avec des mots : elles se \emph{montrent}. Le géographe dispose pour cela de trois outils fondamentaux : la \cle{carte}, le \cle{croquis} et le \cle{diagramme}. Cette dernière section est entièrement consacrée aux travaux pratiques : tu vas apprendre à cartographier les pôles artisanaux du Cameroun, à schématiser une filière, à construire un diagramme statistique et à analyser un dossier documentaire. Ce sont exactement les compétences évaluées au Baccalauréat.

\courssub{TP 1 — Cartographier les pôles artisanaux du Cameroun}

\textbf{Objectif.}\par Localiser et légender, sur un fond de carte du Cameroun, les grands pôles artisanaux par spécialité : la sculpture et le bronze à Foumban (Ouest), le cuir à Maroua (Extrême-Nord), la broderie à Bafoussam (Ouest), la vannerie dans la région du Sud.

\textbf{Intuition.}\par Quand tu expliques à un ami où se trouve le meilleur atelier de bronze, tu lui dis « à Foumban, dans l'Ouest ». Une carte fait la même chose, mais pour tout le pays et pour tous les métiers à la fois : elle transforme une liste de noms en une \emph{vision spatiale}. C'est le passage du récit à la géographie.

\begin{methode}[Légender une carte : les règles d'or]
\begin{enumerate}
\item \textbf{Choisir les figurés.} Un figuré est le symbole graphique qui représente un phénomène : un point pour une ville, un pictogramme pour une spécialité artisanale, une couleur pour une région, une flèche pour un flux.
\item \textbf{Un même phénomène = un même figuré.} Toutes les villes de sculpture doivent avoir exactement le même symbole. On ne mélange jamais les codes à l'intérieur d'une même légende.
\item \textbf{Hiérarchiser.} Les phénomènes les plus importants reçoivent les figurés les plus visibles (points plus gros, couleurs plus vives). Les éléments secondaires (villes repères, frontières) restent discrets.
\item \textbf{Rédiger un titre précis.} Le titre annonce le sujet, le lieu et éventuellement la date : « Les pôles artisanaux du Cameroun par spécialité ».
\item \textbf{Compléter.} Une bonne carte comporte toujours : un titre, une légende complète, une orientation (flèche du nord) et, si possible, une échelle indicative.
\end{enumerate}
\end{methode}

\begin{popfigure}
\begin{center}
\includegraphics[width=\linewidth,height=11cm,keepaspectratio]{N06/cmr_muette.png}
\end{center}
\legende{Fond de carte vierge du Cameroun (limites des régions et des départements, fleuves, lacs, barre d'échelle), à compléter par l'élève : le croquis attendu situe, par un figuré distinct pour chaque spécialité, les pôles artisanaux (Foumban, Bafoussam, Bamenda, Maroua, Douala, Yaoundé) ; chaque figuré est repris à l'identique dans la légende, avec un titre et la flèche du nord. \\ Source : Wikimedia Commons, Flappiefh, CC BY-SA 4.0.}
\end{popfigure}

\begin{attention}[Erreur fréquente au Bac]
Beaucoup de candidats produisent une carte \textbf{sans titre}, ou une carte dont la légende utilise des figurés qui \textbf{ne correspondent pas} à ceux dessinés sur la carte (un rond dans la légende, un carré sur la carte). C'est une faute de méthode lourdement sanctionnée. Avant de rendre ta copie, vérifie point par point : titre, légende complète, nord, cohérence figurés-carte.
\end{attention}

\courssub{TP 2 — Construire le croquis d'une filière artisanale}

\textbf{Objectif.}\par Représenter sous forme de croquis les \cle{maillons} d'une filière artisanale étudiée, par exemple la filière miel de l'Adamaoua ou la filière bronze de Foumban.

\textbf{Intuition.}\par Une filière est une chaîne : le miel ne passe pas directement de la ruche au pot vendu à Douala. Il y a des étapes (production, collecte, transformation, conditionnement, distribution) et des acteurs à chaque étape. Le croquis de filière rend cette chaîne visible d'un coup d'œil, là où un texte demanderait un paragraphe entier.

\begin{methode}[Construire un croquis de filière en 4 étapes]
\begin{enumerate}
\item \textbf{Lister les maillons} dans l'ordre chronologique : matière première $\rightarrow$ production $\rightarrow$ collecte $\rightarrow$ transformation $\rightarrow$ conditionnement $\rightarrow$ commercialisation.
\item \textbf{Identifier les acteurs} de chaque maillon (apiculteurs, coopérative, unité de conditionnement, commerçants, consommateurs).
\item \textbf{Choisir un support graphique} : boîtes rectangulaires pour les maillons, flèches pour les flux de produits, mentions des lieux sous chaque boîte.
\item \textbf{Légender et titrer} : titre explicite, flèches orientées dans le sens du flux, lieux géographiques indiqués.
\end{enumerate}
\end{methode}

\begin{popfigure}
\begin{center}
\begin{tikzpicture}[scale=1.0,
  maillon/.style={rectangle, draw=popdark, fill=popblueL, rounded corners, minimum width=2.6cm, minimum height=0.9cm, align=center, font=\scriptsize},
  trou/.style={rectangle, draw=poporange, dashed, fill=poporangeL, rounded corners, minimum width=2.6cm, minimum height=0.9cm, align=center, font=\scriptsize},
  flux/.style={->, very thick, poporange}]
\node[maillon, align=center] (m1) at (0,0){Production\\ \emph{(apiculteurs, Ngaoundéré)}};
\node[trou, align=center] (m2) at (3.6,0){Maillon 2\\ \emph{(à compléter)}};
\node[trou, align=center] (m3) at (7.2,0){Maillon 3\\ \emph{(à compléter)}};
\node[maillon, align=center] (m4) at (10.8,0){Vente\\ \emph{(marchés, boutiques)}};
\draw[flux] (m1) -- (m2);
\draw[flux] (m2) -- (m3);
\draw[flux] (m3) -- (m4);
\node[below=0.15cm of m2, font=\scriptsize\itshape, text=popmuted] {Qui ? Où ?};
\node[below=0.15cm of m3, font=\scriptsize\itshape, text=popmuted] {Qui ? Où ?};
\end{tikzpicture}
\end{center}
\legende{Modèle de croquis de la filière miel de l'Adamaoua « à trous ». Consigne : complète les maillons 2 et 3 en indiquant l'étape (collecte par la coopérative, puis transformation et conditionnement en pots étiquetés), les acteurs et les lieux. Attendu : maillon 2 = collecte et contrôle qualité par la coopérative ; maillon 3 = conditionnement normé (pots propres, étiquettes) avant expédition vers Yaoundé et Douala.}
\end{popfigure}

\begin{exercice}[Croquis de la filière bronze de Foumban]
À partir de tes connaissances sur l'artisanat de Foumban, construis le croquis complet de la filière bronze : extraction et achat des métaux, fonte à la cire perdue dans les ateliers du quartier des artisans, finition et polissage, vente au marché artisanal de Foumban et exportation vers les galeries de Douala et de l'étranger. N'oublie ni le titre, ni l'orientation des flèches, ni les lieux.
\end{exercice}

\begin{corrige}[Exercice : croquis de la filière bronze de Foumban]
Le croquis attendu comporte cinq maillons reliés par des flèches orientées :
\begin{enumerate}
\item \textbf{Approvisionnement} : achat de cuivre, de zinc et de vieux métaux recyclés (marchés de Foumban et de Bafoussam) ;
\item \textbf{Fonte} : technique de la cire perdue dans les ateliers familiaux du quartier artisanal de Foumban (région de l'Ouest) ;
\item \textbf{Finition} : ciselure, polissage, patine — étape où la qualité et les finitions font la valeur de l'objet ;
\item \textbf{Commercialisation locale} : marché artisanal de Foumban, boutiques touristiques, commandes des notables ;
\item \textbf{Commercialisation externe} : revendeurs de Douala et Yaoundé, exportation vers l'Europe et l'Amérique du Nord.
\end{enumerate}
Le titre (« La filière bronze de Foumban »), les flèches de flux et les noms de lieux sous chaque maillon sont obligatoires pour la note maximale.
\end{corrige}

\courssub{TP 3 — Construire un diagramme statistique}

\textbf{Objectif.}\par À partir d'un tableau de données, construire un histogramme (diagramme en barres) ou un diagramme circulaire, puis dégager une tendance.

\textbf{Intuition.}\par Un tableau de chiffres parle peu ; un graphique parle immédiatement. Si une coopérative vend trois fois plus au quatrième trimestre, la barre trois fois plus haute se voit en une seconde. Le diagramme est un \emph{outil de démonstration} : il rend l'argument visible.

\begin{exemplebox}[Les ventes trimestrielles d'une coopérative artisanale]
Une coopérative d'artisans de l'Ouest vend ses produits (sculptures, tissus brodés, paniers) sur une année. Les ventes, en millions de FCFA (données fictives réalistes construites pour l'exercice), sont les suivantes :
\begin{center}
\begin{tabularx}{\linewidth}{|l|X|X|X|X|}\hline
\rowcolor{popblueL} \textbf{Trimestre} & T1 (janv.\textendash mars) & T2 (avr.\textendash juin) & T3 (juil.\textendash sept.) & T4 (oct.\textendash déc.) \\ \hline
\textbf{Ventes (millions FCFA)} & 8 & 10 & 12 & 24 \\ \hline
\end{tabularx}
\end{center}
Total annuel : $8 + 10 + 12 + 24 = 54$ millions de FCFA. Le quatrième trimestre représente $\dfrac{24}{54} \times 100 \approx 44\,\%$ des ventes annuelles : c'est presque autant que les trois autres trimestres réunis.
\end{exemplebox}

\begin{popfigure}
\begin{center}
\begin{tikzpicture}
\begin{axis}[
  ybar, bar width=0.9cm,
  width=0.85\linewidth, height=6cm,
  xlabel={Trimestres}, ylabel={Ventes (millions de FCFA)},
  symbolic x coords={T1, T2, T3, T4},
  xtick=data,
  ymin=0, ymax=28,
  ytick={0,5,10,15,20,25},
  nodes near coords,
  every node near coord/.append style={font=\scriptsize},
  axis lines=left,
  enlarge x limits=0.2]
\addplot[fill=popblue, draw=popdark] coordinates {(T1,8) (T2,10) (T3,12) (T4,24)};
\end{axis}
\end{tikzpicture}
\end{center}
\legende{Histogramme des ventes trimestrielles d'une coopérative artisanale de l'Ouest (données fictives réalistes, en millions de FCFA). La barre du T4, deux fois plus haute que celle du T3, traduit la saisonnalité des fêtes de fin d'année.}
\end{popfigure}

\begin{exoresolu}[Construire un histogramme et dégager une tendance]
\textbf{Énoncé.} À partir du tableau des ventes trimestrielles ci-dessus : 1) construis un histogramme en respectant les règles de construction ; 2) dégage la tendance principale et propose une explication géographique.
\tcblower
\textbf{Solution détaillée.}
\begin{enumerate}
\item \textbf{Construction.} On trace deux axes perpendiculaires : l'axe horizontal porte les trimestres (variable en catégories), l'axe vertical porte les ventes en millions de FCFA, gradué de 0 à 28 avec un pas régulier de 5. Chaque barre a la même largeur, est séparée de ses voisines, et sa hauteur est proportionnelle à la valeur : 8, 10, 12, 24. On ajoute un titre (« Ventes trimestrielles de la coopérative ») et les étiquettes d'axes avec leurs unités.
\item \textbf{Tendance.} Les ventes progressent régulièrement de T1 à T3 ($+2$ millions par trimestre), puis explosent au T4 : elles doublent ($12 \rightarrow 24$) et atteignent environ 44\,\% du total annuel.
\item \textbf{Explication.} Cette \cle{saisonnalité} correspond aux fêtes de fin d'année (Noël, Nouvel An, mariages et cérémonies concentrés en décembre), période où la demande d'objets de décoration, de cadeaux et de tenues brodées est maximale. L'artisanat camerounais est donc fortement rythmé par le calendrier culturel : les coopératives doivent constituer leurs stocks dès le troisième trimestre.
\end{enumerate}
\end{exoresolu}

\begin{attention}[Pièges de construction des diagrammes]
Trois erreurs reviennent sans cesse : \textbf{(1)} des barres de largeurs inégales ou collées dans un histogramme de catégories ; \textbf{(2)} un axe vertical qui ne commence pas à zéro, ce qui exagère visuellement les écarts ; \textbf{(3)} l'absence d'unité ou de titre. Dans un diagramme circulaire, vérifie toujours que la somme des angles fait bien $360^{\circ}$ : pour le tableau ci-dessus, le T4 occuperait $\dfrac{24}{54} \times 360 = 160^{\circ}$.
\end{attention}

\courssub{TP 4 — Analyser un dossier documentaire}

\textbf{Objectif.}\par Analyser un dossier composé de trois documents de natures différentes : des photographies d'ateliers (décrites), un tableau de données et un extrait de texte sur les normes. C'est l'exercice le plus proche de l'épreuve du Bac.

\begin{experience}[Dossier documentaire : l'atelier de transformation agroalimentaire de Bafoussam]
\textbf{Document 1 — Photographies décrites.} Photo A : un atelier traditionnel de transformation du manioc en bâtons (bobolo) ; les femmes travaillent à la main, sans emballage, les produits sont exposés à l'air libre. Photo B : le même type d'atelier après modernisation ; on aperçoit une machine à écraser électrique, des blouses de travail, des sachets plastiques étiquetés avec une date de fabrication.
\medskip

\textbf{Document 2 — Tableau.} Évolution de la production de l'atelier (données fictives réalistes) :
\begin{center}
\begin{tabularx}{\linewidth}{|l|X|X|X|}\hline
\rowcolor{popblueL} \textbf{Année} & Avant modernisation & 1 an après & 3 ans après \\ \hline
Production (tonnes/an) & 12 & 20 & 35 \\ \hline
Rebut au contrôle qualité (\,\%) & 15 & 8 & 3 \\ \hline
\end{tabularx}
\end{center}

\textbf{Document 3 — Texte.} « Pour vendre dans les supermarchés de Douala, la coopérative a dû respecter des normes d'hygiène : locaux carrelés, emballages alimentaires, étiquetage mentionnant la composition et la date de péremption. Un contrôle qualité est effectué sur chaque lot avant expédition. »
\medskip

\textbf{Consignes d'analyse.} 1) Décris les transformations visibles entre les photos A et B. 2) Calcule l'évolution de la production entre « avant » et « 3 ans après ». 3) Montre, à partir des trois documents croisés, le lien entre modernisation, respect des normes et accès à de nouveaux marchés.
\end{experience}

\begin{methode}[Analyser un dossier documentaire en 4 temps]
\begin{enumerate}
\item \textbf{Identifier} chaque document (nature, auteur ou source, sujet, date).
\item \textbf{Extraire} les informations utiles de chacun, sans les mélanger.
\item \textbf{Croiser} les documents : que confirme le tableau par rapport aux photos ? Qu'explique le texte ?
\item \textbf{Conclure} en répondant précisément à la question posée, en citant les documents (« comme le montre le document 2... »).
\end{enumerate}
\end{methode}

\begin{exoresolu}[Exploitation du dossier de Bafoussam]
\textbf{Énoncé.} Calcule le taux de variation de la production entre la période « avant modernisation » et « 3 ans après », puis formule la conclusion géographique du dossier.
\tcblower
\textbf{Solution.}
\begin{align*}
\text{Taux de variation} &= \frac{35 - 12}{12} \times 100 = \frac{23}{12} \times 100 \approx 192\,\%.
\end{align*}
La production a presque triplé ($\times 2,9$) en trois ans, tandis que le taux de rebut au contrôle qualité est passé de 15\,\% à 3\,\% : l'atelier produit \emph{plus} et \emph{mieux}. Conclusion : la modernisation des équipements (machine à écraser, doc. 1) a permis d'augmenter les quantités (doc. 2), et le respect des normes d'hygiène et d'étiquetage (doc. 3) a ouvert les circuits de vente modernes (supermarchés), ce qui illustre parfaitement le passage d'un artisanat de subsistance à un artisanat de marché.
\end{exoresolu}

\courssub{Restitution et critique collective}

La dernière étape du travail pratique est la \cle{restitution} : chaque groupe présente sa carte, son croquis ou son diagramme à la classe, qui le critique selon une grille commune. Cette démarche, loin d'être une formalité, est le cœur de la formation du géographe : on apprend à \emph{lire} en apprenant à \emph{juger}.

\begin{methode}[Grille de relecture collective d'une production cartographique]
\begin{enumerate}
\item \textbf{Justesse des localisations} : Foumban et Bafoussam sont-ils bien dans l'Ouest ? Maroua bien dans l'Extrême-Nord ? La vannerie bien rattachée au Sud forestier ?
\item \textbf{Clarté de la légende} : chaque figuré de la carte se retrouve-t-il dans la légende, et réciproquement ? Un lecteur extérieur comprend-il la carte sans explication orale ?
\item \textbf{Cohérence des schémas} : les flèches de la filière vont-elles toutes dans le sens du flux ? Les maillons sont-ils dans le bon ordre chronologique ?
\item \textbf{Présence des éléments obligatoires} : titre, nord, échelle indicative, unités sur les axes des diagrammes.
\end{enumerate}
\end{methode}

\begin{savaistu}[Le croquis, arme secrète de ta copie de Bac]
Aux épreuves de géographie du Baccalauréat camerounais, les correcteurs valorisent explicitement les croquis et schémas \textbf{propres, légendés et orientés}. Un croquis clair, même simple, peut rehausser l'ensemble d'une copie : il prouve que le candidat sait organiser spatialement ses connaissances. À l'inverse, un croquis brouillon, sans titre ni légende, dessert même un bon raisonnement. Entraîne-toi à tracer en cinq minutes un fond de carte du Cameroun : c'est un investissement rentable.
\end{savaistu}

\begin{aretenir}[L'essentiel des travaux pratiques]
\begin{itemize}
\item Une \textbf{carte} se construit : fond de carte, figurés choisis et hiérarchisés, titre, légende complète, nord, échelle indicative. Un même phénomène = un même figuré.
\item Un \textbf{croquis de filière} enchaîne les maillons (production $\rightarrow$ collecte $\rightarrow$ transformation $\rightarrow$ conditionnement $\rightarrow$ vente) reliés par des flèches orientées, avec acteurs et lieux.
\item Un \textbf{diagramme} (histogramme, circulaire) traduit un tableau statistique en image ; il doit avoir titre, axes légendés et unités, et permettre de dégager une tendance (comme la saisonnalité des ventes de fin d'année).
\item L'\textbf{analyse de dossier} croise photographies, tableaux et textes pour construire un raisonnement géographique argumenté.
\end{itemize}
\end{aretenir}

\begin{exercice}[Pour t'entraîner : synthèse finale]
Sur un fond de carte du Cameroun tracé à main levée : 1) place et légende les quatre pôles artisanaux étudiés (Foumban, Maroua, Bafoussam, le Sud) ; 2) ajoute par des flèches les flux commerciaux entre ces pôles et les deux métropoles de consommation (Douala, Yaoundé) ; 3) rédige en cinq lignes un commentaire qui dégage l'idée principale de ta carte. Barème indicatif : localisations 4 pts, légende et figurés 4 pts, flux 4 pts, titre et nord 2 pts, commentaire 6 pts.
\end{exercice}

\newpage

\coursec{Activité d'intégration}

\coursec{Leçon 06 : L'évolution de l'artisanat}

\courssub{Objectifs de l'activité}

À l'issue de cette activité, tu seras capable de :

\begin{itemize}
\item \cle{observer} et \cle{décrire} deux photographies d'ateliers artisanaux, puis les \cle{comparer} pour repérer les signes de modernisation ;
\item \cle{inventorier}, \cle{classer} et \cle{interpréter} un tableau statistique de ventes, puis construire un \cle{histogramme} pour identifier le circuit de vente porteur ;
\item mobiliser les notions de \cle{qualité} et de \cle{norme} pour rédiger les mentions obligatoires d'une étiquette conforme aux exigences de l'ANOR ;
\item \cle{schématiser} la filière complète du miel sous forme d'un croquis fléché ;
\item \cle{localiser} et \cle{légender} sur un fond de carte du Cameroun la zone de production et les marchés visés ;
\item rédiger un texte argumenté de \cle{recommandations} pour la modernisation d'une filière artisanale.
\end{itemize}

Cette activité mobilise donc l'ensemble des savoir-faire du programme : observer, localiser et décrire ; représenter, schématiser et cartographier ; inventorier, classer et interpréter des documents.

\courssub{Situation}

Dans la commune de Ngaoundéré, région de l'Adamaoua, la coopérative apicole \og Miel d'Adamaoua \fg{} regroupe \num{120} apiculteurs répartis dans les villages de Belel, Mbang-Mboum et Wakwa. Les savanes arborées de l'Adamaoua, riches en karité, néré et eucalyptus, offrent un miel clair et parfumé, très apprécié. La coopérative produit environ \SI{18}{tonnes} de miel par campagne. Jusqu'ici, la récolte se fait de manière traditionnelle : enfumage des ruches sauvages ou ruches en troncs creux, extraction à la main, filtrage sommaire dans un tissu, puis mise en bouteilles recyclées vendues au marché de Ngaoundéré ou aux voyageurs le long de la Nationale n\textsuperscript{o}1.

Le gérant, M. Hamadou, veut franchir un cap : faire entrer le miel de la coopérative dans les rayons des supermarchés de Douala et de Yaoundé (Santa Lucia, Dovv, Mahima) et développer les ventes vers la diaspora camerounaise en Europe. Mais les responsables des supermarchés exigent un produit \cle{conditionné} selon les \cle{normes} en vigueur : pots en verre ou en plastique alimentaire neufs et scellés, étiquette réglementaire, traçabilité des lots. L'Agence des Normes et de la Qualité (ANOR), basée à Yaoundé, précise les exigences minimales d'étiquetage d'un miel mis en vente.

\textbf{Document 1 — Deux photographies.} Photographie A : un atelier traditionnel à Belel ; deux apiculteurs en vêtements ordinaires pressent les rayons de cire à la main au-dessus d'une bassine ; le miel est filtré dans un pagne tendu sur une marmite, puis versé dans des bouteilles d'eau minérale usagées, sans étiquette. Photographie B : une unité modernisée de conditionnement ; les apiculteurs portent combinaisons, gants et masques ; une machine à désoperculer et un extracteur centrifuge électrique libèrent le miel des rayons ; celui-ci passe dans un filtre inox, est contrôlé (densité, humidité) puis mis en pots de \SI{500}{g} scellés, étiquetés et rangés en cartons imprimés.

\textbf{Document 2 — Ventes de la coopérative (en tonnes) sur trois campagnes.}

\begin{center}
\begin{tabularx}{\linewidth}{|l|X|X|X|X|}
\hline
\rowcolor{popblueL} \textbf{Circuit de vente} & \textbf{Campagne 1} & \textbf{Campagne 2} & \textbf{Campagne 3} & \textbf{Total} \\
\hline
Marché local de Ngaoundéré & 10,0 & 9,5 & 9,0 & 28,5 \\
\hline
Boutiques et dépôts urbains & 3,0 & 4,5 & 6,0 & 13,5 \\
\hline
Export vers la diaspora & 1,0 & 2,0 & 3,0 & 6,0 \\
\hline
\textbf{Total} & 14,0 & 16,0 & 18,0 & 48,0 \\
\hline
\end{tabularx}
\end{center}

\textbf{Document 3 — Extrait des exigences de l'ANOR pour l'étiquetage du miel (extrait non exhaustif).} L'étiquette doit porter : la dénomination de vente (\og miel \fg{}) ; le nom et l'adresse du producteur ou du conditionneur ; le poids net ; la liste des ingrédients (miel 100 \% pur) ; la date de durabilité minimale (DDM) ; le numéro du lot ; les conditions de conservation ; le pays d'origine.

\textbf{Document 4 — Fond de carte du Cameroun} (contour simplifié, régions, principales villes et axes routiers).

\textbf{Problème :} comment la coopérative \og Miel d'Adamaoua \fg{} peut-elle moderniser sa filière pour conquérir les supermarchés et la diaspora ?

\courssub{Consignes}

Par groupes de quatre, et en t'appuyant sur les quatre documents :

\begin{enumerate}
\item \textbf{Décrire} la photographie A puis la photographie B (lieu, acteurs, gestes, équipements, hygiène), puis \textbf{comparer} les deux photographies en dégageant trois signes de modernisation des équipements de travail.
\item À partir du document 2, \textbf{construire} un histogramme des ventes totales par circuit (3 barres), puis \textbf{calculer} la part de chaque circuit dans le total des trois campagnes. \textbf{Identifier} le circuit dominant et le circuit le plus dynamique (le plus fort taux de croissance entre la campagne 1 et la campagne 3).
\item À partir du document 3, \textbf{rédiger} la liste complète des mentions obligatoires qui devront figurer sur l'étiquette du pot de miel de la coopérative, et \textbf{expliquer} en deux phrases pourquoi le respect de ces normes conditionne l'accès aux supermarchés.
\item \textbf{Construire} le croquis fléché de la filière complète du miel : rucher $\rightarrow$ extraction $\rightarrow$ contrôle qualité $\rightarrow$ conditionnement $\rightarrow$ distribution, en indiquant à chaque étape un équipement moderne ou une exigence de qualité.
\item Sur le fond de carte (document 4), \textbf{localiser et légender} : la zone de production (Adamaoua, Ngaoundéré), les marchés urbains visés (Yaoundé, Douala), l'axe de transport principal (Nationale n\textsuperscript{o}1 et chemin de fer Transcamerounais) et le port de sortie pour l'export (Douala).
\item \textbf{Rédiger} un texte de 10 à 15 lignes de recommandations à l'attention de M. Hamadou : équipements à acquérir, normes à respecter, circuits à développer, en t'appuyant sur les résultats des questions précédentes.
\end{enumerate}

Restitution orale de 5 minutes par groupe, puis correction collective au tableau.

\courssub{Ressources à mobiliser}

\begin{aretenir}[Boîte à outils de l'activité]
\begin{itemize}
\item \textbf{Part en pourcentage} : $\text{part} = \dfrac{\text{valeur du circuit}}{\text{total général}} \times 100$.
\item \textbf{Taux de croissance} : $\text{taux} = \dfrac{\text{valeur finale} - \text{valeur initiale}}{\text{valeur initiale}} \times 100$.
\item \textbf{Histogramme} : barres verticales de hauteur proportionnelle aux valeurs ; axes légendés avec titre et unités ; une couleur par circuit.
\item \textbf{Croquis de filière} : étapes en cases reliées par des flèches orientées ; chaque case comporte un équipement ou une exigence ; titre et légende obligatoires.
\item \textbf{Carte} : titre, légende numérotée ou fléchée, orientation (flèche du nord), échelle indicative ; les flux se représentent par des flèches d'épaisseur proportionnelle.
\item \textbf{Notions} : une \cle{norme} est un référentiel technique obligatoire ou volontaire (ANOR au Cameroun) ; la \cle{qualité} est l'aptitude d'un produit à satisfaire les attentes du consommateur et les exigences sanitaires.
\end{itemize}
\end{aretenir}

\courssub{Démarche et corrigé commenté}

\begin{exoresolu}[Corrigé commenté de l'activité]
\textbf{Question 1 — Description et comparaison des photographies.}

Photographie A : atelier traditionnel de Belel. Deux apiculteurs sans équipement de protection pressent les rayons à la main ; le miel est filtré dans un tissu et versé dans des bouteilles recyclées sans étiquette. Les conditions d'hygiène sont précaires, le rendement est faible et le produit n'est pas traçable.

Photographie B : unité modernisée. Les opérateurs portent des combinaisons, gants et masques ; la machine à désoperculer et l'extracteur centrifuge électrique remplacent le pressage manuel ; le miel est filtré sur inox, contrôlé, puis mis en pots neufs scellés et étiquetés.

Comparaison : trois signes de modernisation — (1) la \textbf{mécanisation} (extracteur électrique contre pressage manuel) ; (2) l'\textbf{hygiène et la protection} (combinaisons, filtre inox contre tissu et bassine) ; (3) le \textbf{conditionnement normalisé} (pots scellés étiquetés contre bouteilles recyclées anonymes).

\textbf{Question 2 — Histogramme et calculs.}

Parts des circuits dans le total de \SI{48}{tonnes} :
\begin{align*}
\text{Marché local} &: \frac{28{,}5}{48} \times 100 \approx 59{,}4\,\% \\
\text{Boutiques urbaines} &: \frac{13{,}5}{48} \times 100 \approx 28{,}1\,\% \\
\text{Export diaspora} &: \frac{6}{48} \times 100 = 12{,}5\,\%
\end{align*}

Taux de croissance entre la campagne 1 et la campagne 3 :
\begin{align*}
\text{Marché local} &: \frac{9 - 10}{10} \times 100 = -10\,\% \\
\text{Boutiques} &: \frac{6 - 3}{3} \times 100 = +100\,\% \\
\text{Export} &: \frac{3 - 1}{1} \times 100 = +200\,\%
\end{align*}

\begin{popfigure}
\begin{center}
\begin{tikzpicture}
\begin{axis}[
    ybar,
    bar width=1.5cm,
    width=\linewidth,
    height=8cm,
    ylabel={Ventes totales (tonnes)},
    xlabel={Circuits de vente},
    xtick={1,2,3},
    xticklabels={Marché local, Boutiques urbaines, Export diaspora},
    ymin=0, ymax=35,
    ymajorgrids=true,
    grid style=dashed,
    legend style={at={(0.5,-0.15)}, anchor=north, legend columns=-1},
    nodes near coords,
    nodes near coords align={vertical},
    every node near coord/.append style={font=\small, popdark},
    popblueL/.style={fill=popblueL, draw=popblue},
    poporangeL/.style={fill=poporangeL, draw=poporange},
    popgreenL/.style={fill=popgreenL, draw=popgreen}
]
\addplot[popblueL] coordinates {(1,28.5)};
\addplot[poporangeL] coordinates {(2,13.5)};
\addplot[popgreenL] coordinates {(3,6)};
\legend{Marché local, Boutiques urbaines, Export diaspora}
\end{axis}
\end{tikzpicture}
\end{center}
\legende{Histogramme des ventes totales par circuit sur les trois campagnes (Document 2). Le marché local domine en volume, mais les circuits modernes (boutiques, export) sont en forte progression.}
\end{popfigure}

Interprétation : le marché local reste le circuit \textbf{dominant} (près de 60 \% des volumes) mais il \textbf{régresse} ($-10\,\%$) ; l'export diaspora est le circuit le plus \textbf{dynamique} ($+200\,\%$), suivi des boutiques urbaines ($+100\,\%$). L'avenir de la coopérative se joue donc dans les circuits modernes, ce qui justifie la stratégie de M. Hamadou.

\textbf{Question 3 — Étiquette réglementaire.}

Mentions obligatoires (conformes au Document 3) :
\begin{itemize}
\item Dénomination de vente : \og Miel \fg{} (nom de marque possible : \og Miel pur d'Adamaoua \fg{}) ;
\item Nom et adresse du conditionneur : Coopérative Miel d'Adamaoua, Ngaoundéré, Cameroun ;
\item Poids net : \SI{500}{g} (ou \SI{250}{g}, \SI{1}{kg} selon les formats) ;
\item Liste des ingrédients : Miel 100 \% pur ;
\item Date de durabilité minimale (DDM) : \og À consommer de préférence avant le [jour/mois/année] \fg{} ;
\item Numéro de lot : traçabilité de la récolte (ex. : Lot 2024-ADM-001) ;
\item Conditions de conservation : \og À conserver à l'abri de la chaleur et de la lumière \fg{} ;
\item Pays d'origine : Cameroun.
\end{itemize}

Sans ces mentions, les supermarchés refusent la mise en rayon, car ils engagent leur responsabilité vis-à-vis des consommateurs et des services de contrôle (ANOR, MINCOMMERCE) ; la norme est la condition sine qua non d'entrée sur le marché moderne formel.

\textbf{Question 4 — Croquis de la filière.}

\begin{popfigure}
\begin{center}
\begin{tikzpicture}[node distance=0.6cm,
etape/.style={draw=popblue, fill=popblueL, rounded corners, minimum width=2.8cm, minimum height=1.1cm, align=center, font=\small},
fleche/.style={-{Stealth[length=3mm]}, thick, poporange}]
\node[etape, align=center] (rucher){\textbf{Rucher}\\ Ruches modernes\\ (Kenya Top Bar, Langstroth)\\ Enfumoir, lève-cadres};
\node[etape, right=of rucher, align=center] (extract){\textbf{Extraction}\\ Machine à désoperculer\\ Extracteur centrifuge\\ électrique};
\node[etape, right=of extract, align=center] (controle){\textbf{Contrôle qualité}\\ Réfractomètre (humidité $< 20\,\%$)\\ Densimètre, analyse pollen\\ Norme ANOR / Codex};
\node[etape, below=of controle, align=center] (condit){\textbf{Conditionnement}\\ Pots verre/PEHD neufs\\ Boucheuse-étiqueteuse\\ Étiquette réglementaire\\ Scellage inviolable};
\node[etape, left=of condit, align=center] (distrib){\textbf{Distribution}\\ Camions réfrigérés\\ Supermarchés (Yaoundé, Douala)\\ Plateformes export (Douala)\\ Boutiques spécialisées};
\draw[fleche] (rucher) -- node[above, font=\tiny, popdark] {Récolte} (extract);
\draw[fleche] (extract) -- node[above, font=\tiny, popdark] {Miel brut} (controle);
\draw[fleche] (controle) -- node[right, font=\tiny, popdark] {Conforme} (condit);
\draw[fleche] (condit) -- node[below, font=\tiny, popdark] {Produit fini} (distrib);
\end{tikzpicture}
\end{center}
\legende{Croquis fléché de la filière modernisée du miel d'Adamaoua : chaque étape associe un équipement moderne ou une exigence de qualité normative.}
\end{popfigure}

\textbf{Question 5 — Carte de localisation.}

\begin{popfigure}
\begin{center}
\includegraphics[width=\linewidth,height=9.5cm,keepaspectratio]{N01/cmr_map.jpg}
\end{center}
\legende{Fond de carte : carte générale du Cameroun. Le croquis attendu situe : 1. la zone de production, l'Adamaoua (Ngaoundéré) ; 2. les marchés urbains visés, Yaoundé et Douala ; 3. l'axe d'acheminement du Nord vers le Sud, par la voie ferrée Ngaoundéré--Yaoundé--Douala et les routes qui la suivent ; 4. le flux d'export par le port de Douala, sur le golfe de Guinée. L'échelle est portée par la carte (barre graduée en haut à gauche). \\ Source : Wikimedia Commons, CIA, domaine public.}
\end{popfigure}

\textbf{Question 6 — Texte de recommandations (modèle de rédaction).}

\og Monsieur le gérant, pour conquérir les supermarchés et la diaspora, la coopérative doit d'abord \textbf{moderniser ses équipements} : acquisition de ruches modernes (Kenya Top Bar, Langstroth), d'un extracteur centrifuge électrique, d'une machine à désoperculer, d'un filtre inox, d'un réfractomètre et d'une boucheuse-étiqueteuse semi-automatique. Elle doit ensuite \textbf{respecter scrupuleusement les normes de l'ANOR} : pots neufs scellés, étiquette complète (dénomination, poids net, DDM, lot, origine), contrôle systématique de l'humidité ($< 20\,\%$) et de la pureté du miel, traçabilité par lot. Nos calculs montrent que le marché local, encore dominant (\num{59,4} \%), recule ($-10\,\%$), tandis que l'export diaspora progresse de $+200\,\%$ et les boutiques de $+100\,\%$ : il faut donc \textbf{développer en priorité les circuits modernes} — contrats avec les centrales d'achat des supermarchés de Douala et Yaoundé desservis par la N1 et le Transcamerounais, et organisation d'une filière d'export via le port de Douala. Enfin, la coopérative gagnera à \textbf{structurer sa filière} : groupement des apiculteurs pour l'achat groupé d'équipements, formation aux bonnes pratiques d'hygiène (HACCP simplifié) et création d'une marque collective \og Miel d'Adamaoua \fg{} garante de la qualité. \fg{}
\tcblower
\textbf{Commentaires méthodologiques.} La comparaison des photographies repose sur des critères explicites (équipements, hygiène, conditionnement) et non sur une simple impression. Les calculs de parts et de taux sont posés avec la formule avant l'application numérique : c'est ce qui est attendu au Baccalauréat. L'histogramme visualise la rupture entre volume actuel et dynamique future. Le croquis de filière et la carte respectent les règles cartographiques : titre, légende, flèches orientées, nord. Le texte final articule les trois leviers de la modernisation vus dans la leçon : équipements, normes, circuits.
\end{exoresolu}

\begin{attention}[Pièges fréquents dans cette activité]
\begin{itemize}
\item Confondre le circuit \textbf{dominant} (le plus gros volume : le marché local) avec le circuit le plus \textbf{dynamique} (la plus forte croissance : l'export). Ce sont deux notions différentes ; le Bac sanctionne cette confusion.
\item Oublier la légende, le titre ou l'orientation (flèche Nord) sur la carte : un croquis sans ces éléments est sanctionné (0 point sur la carte).
\item Citer les mentions de l'étiquette de mémoire sans vérifier le Document 3 : l'exercice porte sur l'\textbf{exploitation des documents fournis}.
\item Décrire les photographies sans les comparer : la consigne exige un dégagement explicite des signes de modernisation (mécanisation, hygiène, conditionnement).
\item Ne pas indiquer les unités (tonnes, \%, g) dans les calculs et sur les axes de l'histogramme.
\end{itemize}
\end{attention}

\courssub{Bilan de l'activité}

Cette activité t'a permis de mettre en évidence que la \cle{modernisation de l'artisanat} camerounais ne se réduit pas à l'achat de machines : elle forme un tout cohérent qui associe la \textbf{mécanisation des équipements} (extracteur, filtre inox, réfractomètre), le \textbf{respect des normes et de la qualité} (étiquetage ANOR, contrôles physico-chimiques, traçabilité), l'\textbf{organisation des filières} (coopérative, marque collective, formation) et la \textbf{diversification des circuits de vente} (supermarchés, export diaspora). Tu as également constaté qu'un raisonnement géographique rigoureux s'appuie sur des documents variés — photographies, tableau statistique, texte normatif, carte — et que les outils quantitatifs (parts, taux de croissance, histogramme) permettent de transformer une intuition (\og le marché local est important \fg{}) en un diagnostic précis (\og il domine mais régresse, l'avenir est aux circuits modernes \fg{}). C'est exactement la démarche attendue au Baccalauréat : observer, mesurer, cartographier, puis proposer.

\newpage

\coursec{Méthodes et exercices résolus}

\coursec{L'évolution de l'artisanat}

\courssub{Savoir-faire 1 : Observer, localiser et décrire un phénomène artisanal}

\begin{methode}[Décrire un atelier, un produit ou un paysage artisanal]
La description géographique est la première étape de tout commentaire de document. Elle doit être \cle{organisée}, \cle{précise} et \cle{localisée}.
\begin{enumerate}
\item \textbf{Identifier et localiser} : nommer le lieu (quartier, ville, région), le type d'activité (menuiserie, poterie, tissage, transformation agroalimentaire) et le contexte (marché, village artisanal, zone industrielle).
\item \textbf{Décrire de manière ordonnée} : procéder du général au particulier (d'abord l'ensemble de l'atelier, puis les équipements, puis les gestes et les produits). Pour une photographie, décrire du premier plan vers l'arrière-plan.
\item \textbf{Employer le vocabulaire du cours} : machine à écraser, scie électrique, ponceuse, perceuse, filière, circuit de vente, norme, conditionnement, produit fini.
\item \textbf{Repérer les indices de modernisation} : présence d'électricité, machines, emballages étiquetés, étalages organisés, panneaux publicitaires.
\item \textbf{Conclure par une idée} : dégager ce que le document montre de l'évolution de l'artisanat (gain de productivité, diversification, amélioration de la qualité).
\end{enumerate}
\textbf{Réflexes} : ne jamais interpréter avant d'avoir décrit ; toujours appuyer chaque affirmation sur un élément visible du document ; chiffrer quand c'est possible (nombre d'artisans, de machines, d'étals).
\end{methode}

\begin{exoresolu}[Description d'une photographie : un atelier de menuiserie à Foumban]
\textbf{Énoncé.} Sur une photographie prise dans le quartier artisanal de Foumban (région de l'Ouest), on observe : au premier plan, un menuisier utilisant une scie électrique fixée sur un établi ; à côté, une ponceuse et une perceuse à colonne branchées sur une prise murale ; au deuxième plan, des statues en bois verni alignées sur des étagères, certaines emballées dans du papier bulle et des cartons étiquetés ; à l'arrière-plan, une enseigne « Art et Tradition SARL — exportation » et deux apprentis qui poncent des masques à la main.
Décrivez cette photographie et dégagez ce qu'elle révèle de l'évolution de l'artisanat camerounais. (6 points)
\tcblower
\textbf{Solution détaillée.}

\emph{Étape 1 — Identification et localisation.} Le document est une photographie prise à Foumban, chef-lieu du département du Noun (région de l'Ouest), ville réputée pour son artisanat d'art (sculpture sur bois, bronze, perles). Il s'agit d'un atelier de menuiserie-sculpture.

\emph{Étape 2 — Description ordonnée, du premier plan vers l'arrière-plan.}
\begin{itemize}
\item Au premier plan, un menuisier travaille avec une \cle{scie électrique} fixée sur un établi : l'outil manuel traditionnel (scie à main, herminette) est remplacé par une machine.
\item À côté, une \cle{ponceuse} et une \cle{perceuse} à colonne sont branchées sur le réseau électrique : l'atelier dispose d'équipements modernes et d'une source d'énergie.
\item Au deuxième plan, les produits finis (statues en bois verni) sont rangés sur des étagères ; certains sont \cle{conditionnés} dans du papier bulle et des cartons étiquetés : le conditionnement et la présentation sont soignés.
\item À l'arrière-plan, l'enseigne « Art et Tradition SARL — exportation » indique que l'atelier est organisé en entreprise (forme sociétaire) et vise un marché extérieur ; deux apprentis poncent des masques à la main, signe que le savoir-faire traditionnel et la formation par apprentissage persistent.
\end{itemize}

\emph{Étape 3 — Interprétation.} La photographie révèle trois évolutions majeures de l'artisanat camerounais :
\begin{enumerate}
\item la \cle{modernisation des équipements} (scie électrique, ponceuse, perceuse), qui augmente la productivité et la précision du travail ;
\item l'\cle{amélioration de la qualité et du conditionnement} (vernis, emballage, étiquetage), qui permet de respecter les exigences des clients lointains ;
\item l'\cle{ouverture sur de nouveaux marchés} (exportation), rendue possible par l'organisation de l'atelier en entreprise structurée.
\end{enumerate}

\emph{Conclusion.} Ce document illustre la transition de l'artisanat camerounais : d'un savoir-faire purement manuel destiné au village, il devient une activité mécanisée, normalisée et tournée vers des marchés nationaux et internationaux, tout en conservant ses techniques traditionnelles (ponçage manuel, apprentissage).
\end{exoresolu}

\courssub{Savoir-faire 2 : Légender une carte de l'artisanat camerounais}

\begin{methode}[Construire et légender une carte ou un croquis]
\begin{enumerate}
\item \textbf{Lire la consigne et choisir le fond de carte} : contour simplifié du Cameroun, limites régionales si nécessaire, flèche du nord et échelle indicative.
\item \textbf{Choisir les informations à représenter} : localiser les grands foyers artisanaux (Foumban pour la sculpture et le bronze, Bamessingui près de Ndop pour la poterie, Bafoussam et Bandjoun pour le tissage et le bambou, Maroua pour le cuir et la vannerie, Douala et Yaoundé pour la menuiserie urbaine et la transformation alimentaire).
\item \textbf{Associer à chaque information un symbole simple} : points ou pictogrammes pour les villes-artisanats, zones colorées pour les grandes régions de production, flèches pour les flux de vente (vers les villes, les ports, l'étranger).
\item \textbf{Rédiger la légende} : chaque symbole est expliqué, numéroté et classé (d'abord les surfaces, puis les points, puis les flux). La légende doit être \emph{hiérarchisée} et \emph{exhaustive} : tout symbole du croquis figure dans la légende, et rien d'autre.
\item \textbf{Soigner la présentation} : titre précis (quoi, où, quand), tracés nets, noms de lieux lisibles et correctement orthographiés.
\end{enumerate}
\textbf{Formule à retenir} : une bonne carte = un titre + un fond + des symboles + une légende hiérarchisée + une orientation.
\end{methode}

\begin{exoresolu}[Croquis : les grands foyers de l'artisanat camerounais et leurs débouchés]
\textbf{Énoncé.} À l'aide d'un fond de carte du Cameroun, construisez un croquis intitulé « Les principaux foyers artisanaux du Cameroun et leurs circuits de vente ». Vous localiserez au moins cinq foyers artisanaux, vous distinguerez les grandes spécialités et vous représenterez les flux de commercialisation vers les métropoles et le port de Douala. (8 points)
\tcblower
\textbf{Solution détaillée.}

\emph{Étape 1 — Choix des informations.} On retient cinq foyers représentatifs :
\begin{itemize}
\item Foumban (Ouest) : sculpture sur bois, bronze, objets d'art ;
\item Bamessingui près de Ndop (Nord-Ouest) : poterie ;
\item Bafoussam et Bandjoun (Ouest) : tissage, bambou, vannerie ;
\item Maroua (Extrême-Nord) : cuir, peausserie, vannerie ;
\item Douala et Yaoundé : menuiserie urbaine, transformation agroalimentaire artisanale (huile de palme, jus, farines).
\end{itemize}
Les flux de vente convergent vers Douala (port d'exportation) et Yaoundé (marché national et touristique).

\emph{Étape 2 — Réalisation du croquis.}

\begin{popfigure}
\begin{center}
\includegraphics[width=\linewidth,height=9.5cm,keepaspectratio]{N01/cmr_map.jpg}
\end{center}
\legende{Fond de carte : carte générale du Cameroun (régions, villes, routes, chemin de fer). Le croquis attendu situe : 1. les foyers artisanaux de production : Foumban et Bafoussam (Ouest), Bamenda (Nord-Ouest), Maroua (Extrême-Nord) ; 2. les métropoles de commercialisation : Douala (port d'exportation) et Yaoundé (marché national) ; 3. des flèches de flux de vente de ces foyers vers les deux métropoles. L'échelle est lisible sur la barre graduée de la carte. \\ Source : Wikimedia Commons, CIA, domaine public.}
\end{popfigure}

\emph{Étape 3 — Vérifications.} Le titre précise le sujet, le lieu et le contenu ; la flèche du nord oriente le croquis ; chaque symbole (points de production, points de commercialisation, flèches de flux) est expliqué dans une légende hiérarchisée (points de production, puis points de vente, puis flux) ; les noms de villes sont lisibles et bien orthographiés.

\emph{Conclusion.} Le croquis montre que les foyers artisanaux sont répartis sur tout le territoire (Ouest, Nord-Ouest, Extrême-Nord, métropoles du Sud) et que leurs produits convergent vers Douala et Yaoundé, ce qui illustre l'organisation des circuits de vente et le rôle du port de Douala dans l'exportation de l'artisanat camerounais.
\end{exoresolu}

\courssub{Savoir-faire 3 : Construire un diagramme (schéma de filière)}

\begin{methode}[Schématiser une filière artisanale]
Une \cle{filière} est l'ensemble des étapes qui vont de la matière première au produit vendu au consommateur. Pour la schématiser :
\begin{enumerate}
\item \textbf{Identifier les maillons} : approvisionnement en matière première $\rightarrow$ transformation (atelier) $\rightarrow$ conditionnement et contrôle qualité $\rightarrow$ distribution (circuits de vente) $\rightarrow$ consommateur.
\item \textbf{Disposer les maillons dans des cases} reliées par des flèches orientées dans le sens du flux du produit.
\item \textbf{Annoter les flèches} avec les équipements, les acteurs ou les lieux (machine à écraser, coopérative, marché, supermarché, exportation).
\item \textbf{Ajouter éventuellement les flux annexes} : formation, financement (microcrédit), contrôle des normes, en les plaçant au-dessus ou en dessous de la chaîne principale.
\item \textbf{Titrer le schéma} et vérifier qu'il se lit de gauche à droite sans ambiguïté.
\end{enumerate}
\textbf{Réflexe} : un schéma de filière n'est pas une liste : les flèches montrent les \emph{transformations} et les \emph{déplacements} du produit.
\end{methode}

\begin{exoresolu}[Schéma : la filière artisanale de l'huile de palme dans le Littoral]
\textbf{Énoncé.} Dans les environs de Dibombari (région du Littoral), des artisanes transforment les noix de palme en huile rouge vendue à Douala. Construisez un schéma de cette filière en faisant apparaître les maillons, les équipements modernes et les circuits de vente. (6 points)
\tcblower
\textbf{Solution détaillée.}

\emph{Étape 1 — Identification des maillons.} Régimes de palme (matière première) $\rightarrow$ atelier de transformation (cuisson, machine à écraser, décantation) $\rightarrow$ conditionnement (bidons étiquetés, contrôle de qualité) $\rightarrow$ vente (marchés de Douala, boutiques, grossistes) $\rightarrow$ consommateurs.

\emph{Étape 2 — Construction du schéma.}

\begin{popfigure}
\begin{tikzpicture}[node distance=0.35cm and 0.55cm,
  box/.style={draw=popdark,fill=popblueL,rounded corners,minimum height=0.9cm,align=center,font=\small,text width=2.5cm},
  flux/.style={-{Stealth},very thick,poporange}]
\node[box] (a) {Récolte des régimes de palme (villages de Dibombari)};
\node[box,right=of a] (b) {Atelier : cuisson, machine à écraser, décantation};
\node[box,right=of b] (c) {Conditionnement : bidons étiquetés, contrôle qualité};
\node[box,right=of c] (d) {Circuits de vente : marchés et boutiques de Douala};
\node[box,right=of d] (e) {Consommateurs};
\draw[flux] (a) -- (b);
\draw[flux] (b) -- (c);
\draw[flux] (c) -- (d);
\draw[flux] (d) -- (e);
\node[above=0.25cm of b,font=\small\itshape,popdark] {modernisation des équipements};
\node[above=0.25cm of c,font=\small\itshape,popdark] {respect des normes};
\end{tikzpicture}
\legende{Schéma de la filière artisanale de l'huile de palme (environs de Dibombari, Littoral). Les cases représentent les maillons, les flèches orange le flux du produit, les annotations les apports de la modernisation.}
\end{popfigure}

\emph{Étape 3 — Analyse.} Le schéma fait apparaître :
\begin{itemize}
\item la \cle{modernisation} au maillon transformation : la machine à écraser remplace le pilage au mortier, ce qui multiplie le rendement (une artisane traite en une journée plusieurs dizaines de kilogrammes de noix contre quelques kilogrammes autrefois) ;
\item le \cle{respect des normes} au maillon conditionnement : bidons propres, étiquettes indiquant le produit, le poids et le producteur ;
\item l'\cle{organisation des circuits de vente} : le produit ne se vend plus seulement au bord de la route mais rejoint les marchés structurés de Douala.
\end{itemize}

\emph{Conclusion.} La filière artisanale de l'huile de palme illustre l'évolution de l'artisanat camerounais : mécanisation de la transformation, amélioration du conditionnement et intégration à des circuits commerciaux urbains organisés.
\end{exoresolu}

\courssub{Savoir-faire 4 : Inventorier, classer et interpréter un tableau statistique}

\begin{methode}[Exploiter un tableau de données sur l'artisanat]
\begin{enumerate}
\item \textbf{Lire le titre, les unités et la source} avant tout calcul (effectifs, tonnes, millions de FCFA, pourcentages).
\item \textbf{Classer les données} : ranger les modalités par ordre décroissant, regrouper éventuellement en catégories (artisanat d'art, artisanat alimentaire, artisanat du bois, artisanat du textile).
\item \textbf{Calculer les indicateurs utiles} :
\begin{itemize}
\item part en pourcentage : $\text{part} = \dfrac{\text{effectif de la catégorie}}{\text{effectif total}} \times 100$ ;
\item taux de variation entre deux dates : $t = \dfrac{V_{2} - V_{1}}{V_{1}} \times 100$ ;
\item valeur moyenne par unité (ex. production par artisan) : $\text{moyenne} = \dfrac{\text{production totale}}{\text{nombre d'artisans}}$.
\end{itemize}
\item \textbf{Interpréter} : dégager deux ou trois faits saillants (catégorie dominante, progression, disparité) et les relier aux savoirs du cours (modernisation, produits nouveaux, normes, filières).
\item \textbf{Rédiger} une conclusion courte qui répond exactement à la question posée.
\end{enumerate}
\textbf{Réflexes} : arrondir les pourcentages au dixième ; vérifier que la somme des parts vaut environ 100 \% ; toujours citer les chiffres dans l'interprétation.
\end{methode}

\begin{exoresolu}[Étude statistique : la production artisanale d'une coopérative]
\textbf{Énoncé.} Une coopérative artisanale de Bafoussam regroupant 50 artisans a vendu en 2023 les produits suivants :

\begin{center}
\begin{tabularx}{\linewidth}{|l|X|X|}\hline
\rowcolor{popblueL} \textbf{Catégorie de produits} & \textbf{Chiffre d'affaires (millions de FCFA)} & \textbf{Chiffre d'affaires 2019 (millions de FCFA)} \\ \hline
Transformation agroalimentaire (jus, farines, huiles) & 60 & 25 \\ \hline
Habillement et décoration (pagne, tissage, perles) & 45 & 30 \\ \hline
Meubles et objets en bois & 30 & 20 \\ \hline
Vannerie et poterie & 15 & 15 \\ \hline
\end{tabularx}
\end{center}

1. Calculez la part de chaque catégorie dans le chiffre d'affaires de 2023. 2. Calculez le taux de variation du chiffre d'affaires total entre 2019 et 2023. 3. Interprétez ces résultats en les reliant à l'évolution de l'artisanat camerounais. (8 points)
\tcblower
\textbf{Solution détaillée.}

\emph{Question 1 — Parts en 2023.} Le chiffre d'affaires total de 2023 est :
\[ 60 + 45 + 30 + 15 = 150 \text{ millions de FCFA.} \]
On applique la formule $\text{part} = \dfrac{\text{catégorie}}{\text{total}} \times 100$ :
\begin{align*}
\text{Agroalimentaire} &: \frac{60}{150} \times 100 = 40{,}0\ \% \\
\text{Habillement-décoration} &: \frac{45}{150} \times 100 = 30{,}0\ \% \\
\text{Meubles et bois} &: \frac{30}{150} \times 100 = 20{,}0\ \% \\
\text{Vannerie et poterie} &: \frac{15}{150} \times 100 = 10{,}0\ \%
\end{align*}
Vérification : $40 + 30 + 20 + 10 = 100\ \%$. Le calcul est cohérent.

\emph{Question 2 — Taux de variation 2019-2023.} Le total de 2019 est :
\[ 25 + 30 + 20 + 15 = 90 \text{ millions de FCFA.} \]
\[ t = \frac{150 - 90}{90} \times 100 = \frac{60}{90} \times 100 \approx 66{,}7\ \%. \]
Le chiffre d'affaires total a augmenté d'environ 66,7 \% en quatre ans. Par artisan, il est passé de $\dfrac{90}{50} = 1{,}8$ à $\dfrac{150}{50} = 3$ millions de FCFA, soit une progression de $\dfrac{3 - 1{,}8}{1{,}8} \times 100 \approx 66{,}7\ \%$ également (l'effectif étant constant).

\emph{Question 3 — Interprétation.}
\begin{itemize}
\item La \cle{transformation agroalimentaire} est devenue la première catégorie (40 \% du chiffre d'affaires, contre $\frac{25}{90} \times 100 \approx 27{,}8\ \%$ en 2019) : c'est l'illustration des \cle{produits nouveaux} (jus conditionnés, farines, huiles) apparus grâce aux machines à écraser et aux nouvelles méthodes de conditionnement.
\item L'habillement et la décoration progressent en valeur (de 30 à 45 millions, soit $+50\ \%$) : la demande de produits artisanaux modernisés (pagnes tissés, articles de décoration) reste forte.
\item La vannerie et la poterie stagnent (15 millions, part passée de 16,7 \% à 10 \%) : les productions les plus traditionnelles, moins mécanisées, bénéficient moins de la dynamique.
\item La forte hausse globale ($+66{,}7\ \%$) traduit les effets conjugués de la \cle{modernisation des équipements}, de l'\cle{amélioration de la qualité} et de l'\cle{organisation des circuits de vente} (coopérative, marchés urbains).
\end{itemize}

\emph{Conclusion.} Entre 2019 et 2023, la coopérative de Bafoussam a vu son chiffre d'affaires croître des deux tiers, tiré par les produits nouveaux issus de la transformation agroalimentaire : l'évolution de l'artisanat camerounais passe par la mécanisation, la diversification des produits et la structuration des filières.
\end{exoresolu}

\courssub{Savoir-faire 5 : Rédiger un commentaire de document (synthèse)}

\begin{methode}[Construire un commentaire composé sur l'évolution de l'artisanat]
\begin{enumerate}
\item \textbf{Analyser le sujet et les documents} : souligner les mots-clés de la consigne, dégager le problème posé (comment l'artisanat camerounais se modernise-t-il ?).
\item \textbf{Dégager une problématique} et un plan en deux ou trois parties équilibrées, par exemple : I. un héritage traditionnel encore vivant ; II. une modernisation visible (équipements, produits nouveaux, qualité et normes) ; III. une intégration aux filières et aux marchés.
\item \textbf{Rédiger l'introduction} : accroche (situation concrète), définition des termes clés (artisanat, \cle{qualité}, \cle{norme}), problématique, annonce du plan.
\item \textbf{Développer} : une idée par paragraphe, appuyée sur un exemple précis (Foumban, Maroua, Dibombari) et un chiffre du document ; enchaîner avec des connecteurs logiques (d'abord, ensuite, enfin ; cependant, ainsi).
\item \textbf{Conclure} : réponse à la problématique et ouverture (défis restants : accès au financement, à l'électricité, aux normes d'exportation).
\end{enumerate}
\textbf{Réflexe} : chaque partie doit mobiliser au moins un savoir du programme et un élément du document.
\end{methode}

\begin{exoresolu}[Commentaire : « L'artisanat camerounais entre tradition et modernité »]
\textbf{Énoncé.} À partir de vos connaissances et des documents étudiés (photographie de l'atelier de Foumban, tableau de la coopérative de Bafoussam, croquis des foyers artisanaux), rédigez un commentaire organisé répondant à la question : comment l'artisanat camerounais évolue-t-il tout en préservant son héritage traditionnel ? (10 points)
\tcblower
\textbf{Solution détaillée (corrigé rédigé).}

\emph{Introduction.} Sur les marchés de Foumban, les statues en bois verni côtoient désormais les cartons étiquetés destinés à l'exportation. L'\cle{artisanat}, ensemble des activités de production manuelle ou semi-mécanisée exercées par de petits producteurs, est un pilier de l'économie camerounaise et de l'intégration nationale. Longtemps fondé sur des techniques héritées, il connaît de profondes transformations. Comment l'artisanat camerounais évolue-t-il tout en préservant son héritage traditionnel ? Nous verrons d'abord que l'artisanat repose sur un héritage traditionnel encore vivant, puis qu'il se modernise dans ses équipements, ses produits et sa qualité, enfin qu'il s'organise en filières intégrées aux marchés.

\emph{I. Un héritage traditionnel encore vivant.} L'artisanat camerounais puise ses racines dans des savoir-faire séculaires répartis sur tout le territoire : sculpture sur bois et bronze à Foumban, poterie à Bamessingui, cuir et vannerie à Maroua, tissage à Bafoussam et Bandjoun. La transmission se fait par apprentissage : dans l'atelier photographié à Foumban, deux apprentis poncent encore les masques à la main. Ces productions traditionnelles (vannerie, poterie) demeurent, même si leur poids relatif recule : dans la coopérative de Bafoussam, leur part est tombée de 16,7 \% à 10 \% du chiffre d'affaires entre 2019 et 2023.

\emph{II. Une modernisation visible.} Cependant, l'artisanat se transforme. D'abord, les \cle{équipements se modernisent} : scies électriques, ponceuses, perceuses et machines à écraser augmentent la productivité et la précision, comme le montre l'atelier de Foumban. Ensuite, des \cle{produits nouveaux} apparaissent dans la décoration, l'habillement et surtout la transformation alimentaire : celle-ci est devenue la première source de revenus de la coopérative de Bafoussam (40 \% du chiffre d'affaires en 2023, contre 27,8 \% en 2019). Enfin, la \cle{qualité} s'améliore et le respect des \cle{normes} progresse : vernis, emballages en papier bulle, cartons étiquetés, bidons propres et contrôles de fabrication permettent d'atteindre des clients exigeants.

\emph{III. Une intégration aux filières et aux marchés.} Ainsi, l'artisanat s'organise en \cle{filières} complètes, de la matière première au consommateur, comme dans la filière huile de palme de Dibombari (récolte, machine à écraser, conditionnement, vente à Douala). Les \cle{circuits de vente} se structurent : coopératives, marchés urbains, boutiques, et exportation via le port de Douala, comme l'indique l'enseigne « exportation » de l'atelier de Foumban. Le résultat est spectaculaire : le chiffre d'affaires de la coopérative de Bafoussam a progressé d'environ 66,7 \% entre 2019 et 2023.

\emph{Conclusion.} L'artisanat camerounais évolue donc sans renier son héritage : il conjugue gestes traditionnels et machines modernes, produits ancestraux et produits nouveaux, marchés villageois et circuits d'exportation. Restent à relever des défis : l'accès à l'électricité et au financement, la formation aux normes internationales, qui conditionneront la poursuite de cette mutation.
\end{exoresolu}

\begin{attention}[Les 5 erreurs qui coûtent des points au Bac]
\begin{enumerate}
\item \textbf{Confondre description et interprétation.} Décrire un document, c'est dire ce qu'on \emph{voit} ; interpréter, c'est expliquer ce que cela \emph{signifie}. Mélanger les deux dès la première phrase fait perdre des points : décris d'abord, ordonné et localisé, puis dégage l'idée.
\item \textbf{Oublier un élément obligatoire du croquis.} Un croquis sans titre, sans flèche du nord ou avec une légende incomplète (symbole non expliqué, légende non hiérarchisée) est lourdement pénalisé. Vérifie toujours : titre + orientation + légende exhaustive où chaque symbole du croquis est expliqué.
\item \textbf{Se tromper dans les formules statistiques.} Erreurs classiques : calculer une part sur le mauvais total (celui de 2019 au lieu de 2023), oublier le $\times 100$, ou écrire le taux de variation comme $\dfrac{V_2}{V_1} \times 100$ au lieu de $\dfrac{V_2 - V_1}{V_1} \times 100$. Vérifie enfin que la somme des parts vaut 100 \%.
\item \textbf{Rédiger un schéma de filière comme une simple liste.} Sans flèches orientées entre les maillons, le schéma ne montre ni les transformations ni les flux : il ne vaut presque rien. Chaque case doit être reliée par une flèche, idéalement annotée (équipement, acteur, lieu).
\item \textbf{Employer un vocabulaire vague au lieu des notions du programme.} Écrire « les artisans travaillent mieux » au lieu de « modernisation des équipements », « amélioration de la qualité et respect des normes », « organisation des filières et des circuits de vente » prive la copie des mots-clés attendus. Définis aussi précisément \cle{qualité} (aptitude d'un produit à satisfaire l'usage attendu) et \cle{norme} (règle technique de fabrication à respecter) dès que la question les mobilise.
\end{enumerate}
\end{attention}

\newpage

\coursec{Exercices}

\courssub{Je vérifie mes connaissances}

\begin{exercice}[QCM : les notions clés]
Pour chaque question, choisis la (ou les) bonne(s) réponse(s).

\textbf{1.} Une \cle{norme} est :
\begin{itemize}
\item[a)] une règle imposée par les clients étrangers uniquement ;
\item[b)] un document de référence fixant les caractéristiques qu'un produit doit respecter pour être mis sur le marché ;
\item[c)] une taxe payée par l'artisan ;
\item[d)] un label réservé aux produits industriels.
\end{itemize}

\textbf{2.} La \cle{qualité} d'un produit artisanal désigne :
\begin{itemize}
\item[a)] son prix de vente élevé ;
\item[b)] l'ensemble de ses caractéristiques lui permettant de satisfaire les besoins et les attentes du consommateur ;
\item[c)] sa fabrication à la main ;
\item[d)] sa rareté.
\end{itemize}

\textbf{3.} Parmi ces équipements, lesquels traduisent la modernisation de l'artisanat ?
\begin{itemize}
\item[a)] la scie électrique ;
\item[b)] la ponceuse ;
\item[c)] la herse tirée par des bœufs ;
\item[d)] la machine à écraser (manioc, huile de palme).
\end{itemize}

\textbf{4.} Un \cle{circuit court} de vente est un circuit :
\begin{itemize}
\item[a)] qui passe par plusieurs intermédiaires ;
\item[b)] où le producteur vend directement au consommateur ou avec un seul intermédiaire ;
\item[c)] réservé à l'exportation ;
\item[d)] qui utilise obligatoirement Internet.
\end{itemize}
\end{exercice}

\begin{exercice}[Vrai ou faux ? Justifie]
Réponds par vrai ou faux, puis justifie ta réponse en une ou deux phrases.
\begin{enumerate}
\item La modernisation des équipements de travail supprime totalement le savoir-faire traditionnel de l'artisan.
\item Le contrôle de la qualité et le respect des normes de fabrication facilitent l'accès des produits artisanaux camerounais aux marchés extérieurs.
\item Une filière regroupe uniquement les artisans qui fabriquent le produit.
\item Les nouvelles méthodes de conditionnement (emballage, étiquetage) participent de l'amélioration de la qualité du produit.
\end{enumerate}
\end{exercice}

\begin{exercice}[Définitions]
Définis en deux ou trois phrases chacun des termes suivants, en donnant à chaque fois un exemple camerounais :
\begin{enumerate}
\item artisanat ;
\item filière ;
\item norme ;
\item circuit de vente.
\end{enumerate}
\end{exercice}

\begin{exercice}[Calculs directs : lire une évolution]
Une coopérative d'artisans potiers de la région de l'Ouest a vendu \num{1200} pièces en 2018. Après l'achat de tours électriques et de fours améliorés, elle en a vendu \num{2100} en 2023.
\begin{enumerate}
\item Calcule le taux de variation des ventes entre 2018 et 2023 (arrondis à 0,1 \%).
\item Sur les \num{2100} pièces vendues en 2023, 630 ont été vendues dans les supermarchés urbains. Calcule la part de ce circuit de vente en pourcentage.
\item Le prix moyen d'une pièce est passé de \num{1500} FCFA à \num{2300} FCFA grâce au meilleur conditionnement. Calcule le chiffre d'affaires de la coopérative en 2023.
\end{enumerate}
\end{exercice}

\courssub{Je m'entraîne}

\begin{exercice}[Classer les équipements]
Voici une liste d'équipements et d'outils : machette, scie électrique, ponceuse, couteau à évider, perceuse, machine à écraser le manioc, mortier et pilon, tour de potier électrique, fer à repasser à charbon, machine à coudre électrique.
\begin{enumerate}
\item Classe ces équipements en deux colonnes : « équipements traditionnels » et « équipements modernisés ».
\item Pour trois équipements modernisés de ton choix, indique le métier artisanal concerné et l'effet de l'équipement sur le travail de l'artisan (rapidité, précision, pénibilité, volume de production).
\end{enumerate}
\end{exercice}

\begin{exercice}[Identifier les produits nouveaux]
Le marché artisanal de Foumban (région de l'Ouest) propose aujourd'hui, à côté des statues et masques traditionnels, des meubles de décoration d'intérieur, des accessoires de mode en tissu Ndop, des sacs en perles et des objets utilitaires (lampes, cadres, porte-clés).
\begin{enumerate}
\item Distingue, dans cette liste, les produits traditionnels des produits nouveaux.
\item Explique en quelques lignes à quels marchés nouveaux (clientèles) ces produits nouveaux s'adressent.
\item Montre le lien entre l'apparition de ces produits et la modernisation des équipements.
\end{enumerate}
\end{exercice}

\begin{exercice}[Lire une étiquette]
Sur le pot de miel conditionné par un groupement d'apiculteurs de l'Adamaoua, on lit : « Miel pur de l'Adamaoua — 500 g — GIC Miel du Vina, Banyo — Date de production : 03/2024 — À consommer de préférence avant : 03/2026 — Lot n° 45 ».
\begin{enumerate}
\item Releve sur cette étiquette quatre mentions qui rassurent le consommateur.
\item Explique pourquoi ces mentions traduisent le respect des normes de conditionnement et de présentation.
\item Quel avantage commercial ce conditionnement donne-t-il à ce miel par rapport au miel vendu en vrac dans des bouteilles non étiquetées ?
\end{enumerate}
\end{exercice}

\begin{exercice}[Schématiser une filière]
La filière du poivre de Penja (région du Littoral), première Indication Géographique Protégée (IGP) d'Afrique subsaharienne (2013), comprend : des producteurs, un groupement qui contrôle la qualité, des conditionneurs, des exportateurs et des distributeurs au Cameroun et en Europe.
\begin{enumerate}
\item Construis un schéma de la filière sous forme de cases reliées par des flèches, de la production à la vente.
\item Indique sur ton schéma, à chaque étape, l'opération réalisée (culture et récolte, tri et séchage, contrôle qualité, emballage, exportation, distribution).
\item Explique en cinq lignes comment l'organisation de cette filière et le respect des normes ont ouvert le marché international au poivre de Penja.
\end{enumerate}
\end{exercice}

\courssub{Je me prépare au Bac}

\begin{exercice}[Étude documentaire : deux ateliers, deux époques]
\textbf{Document 1.} Photographie d'un atelier de sculpture à Foumban dans les années 1970 : l'artisan travaille assis par terre, à la main, avec couteau, herminette et racloir ; les statues finies sont exposées à même le sol devant l'atelier.

\textbf{Document 2.} Photographie d'un atelier de menuiserie-ébénisterie à Yaoundé en 2022 : scie circulaire électrique, ponceuse, perceuse à colonne ; les meubles finis sont vernis, emballés dans du film plastique et étiquetés avant livraison.

\textbf{Document 3.} Production d'une coopérative d'ébénistes (ordres de grandeur) :

\begin{center}
\begin{tabularx}{\linewidth}{|l|X|X|X|}
\hline
\rowcolor{popblueL} Année & Pièces produites & Pièces refusées au contrôle qualité & Part vendue hors de la région \\
\hline
2010 & 400 & 120 & 10 \% \\
\hline
2015 & 900 & 90 & 25 \% \\
\hline
2022 & 1800 & 72 & 45 \% \\
\hline
\end{tabularx}
\end{center}

\begin{enumerate}
\item Décris précisément les différences d'équipement entre les deux ateliers (documents 1 et 2).
\item À partir du document 3, calcule le taux de refus au contrôle qualité en 2010 et en 2022, puis le taux de variation de la production entre 2010 et 2022.
\item Dégage les transformations de l'artisanat camerounais qui apparaissent dans les trois documents (équipements, qualité, débouchés).
\item Montre comment ces transformations élargissent les marchés des artisans camerounais.
\end{enumerate}
\end{exercice}

\begin{exercice}[Construction graphique : les circuits de vente]
Une filière artisanale de transformation du manioc en gari conditionné (région du Centre) a réalisé en 2023 les ventes suivantes :

\begin{center}
\begin{tabularx}{\linewidth}{|l|X|}
\hline
\rowcolor{popblueL} Circuit de vente & Ventes (tonnes) \\
\hline
Vente directe au marché local & 90 \\
\hline
Détaillants et boutiques de quartier & 60 \\
\hline
Supermarchés de Yaoundé et Douala & 75 \\
\hline
Exportation vers la diaspora (Europe) & 15 \\
\hline
Total & 240 \\
\hline
\end{tabularx}
\end{center}

\begin{enumerate}
\item Calcule la part de chaque circuit de vente en pourcentage du total, puis l'angle correspondant pour un diagramme circulaire (arrondis au degré).
\item Construis le diagramme circulaire des ventes de cette filière (titre et légende obligatoires).
\item Construis un croquis de la filière, du producteur de manioc au consommateur final, en faisant apparaître les quatre circuits de vente.
\item Interprète les résultats : quel circuit domine ? Que révèle la place des supermarchés et de l'exportation sur la modernisation de cette filière ?
\end{enumerate}
\end{exercice}

\begin{exercice}[Dissertation guidée : artisanat et intégration nationale]
\textbf{Sujet :} « Montrez, à partir d'exemples camerounais précis, que la modernisation de l'artisanat contribue à l'intégration économique nationale. »

Tu rédigeras un devoir structuré en suivant le plan imposé ci-dessous.

\textbf{Introduction :} amène le sujet à partir d'une situation concrète (par exemple l'achat à Douala d'un meuble fabriqué à Foumban ou d'un tissu Ndop du Nord-Ouest), définis l'artisanat et sa modernisation, puis annonce le plan.

\textbf{I. Un artisanat hérité, ancré dans les régions du Cameroun}
\begin{enumerate}
\item Présente les grands pôles artisanaux et leurs spécialités (sculpture et bronzes de Foumban, poterie, tissage du Ndop, vannerie, transformation alimentaire...).
\item Montre les limites de l'artisanat traditionnel (faible productivité, qualité inégale, marchés locaux étroits).
\end{enumerate}

\textbf{II. Une modernisation qui transforme la production}
\begin{enumerate}
\item Explique le rôle des équipements modernes (machines à écraser, scies électriques, ponceuses, perceuses) et des produits nouveaux (décoration, habillement, transformation alimentaire).
\item Montre comment le contrôle de la qualité, le respect des normes et les nouvelles méthodes de conditionnement ouvrent des marchés nationaux et internationaux (exemple : l'IGP du poivre de Penja).
\end{enumerate}

\textbf{III. Une contribution à l'intégration économique nationale}
\begin{enumerate}
\item Montre comment l'organisation des filières et des circuits de vente relie les régions entre elles (flux de produits artisanaux vers Douala, Yaoundé et les autres régions).
\item Dégage les effets : revenus des artisans, emplois, mise en valeur des ressources locales, rayonnement du Cameroun à l'extérieur.
\end{enumerate}

\textbf{Conclusion :} bilan de la démonstration et ouverture (par exemple : les défis restants — financement, formation, accès à l'électricité).
\end{exercice}

\newpage

\coursec{Corrigés des exercices}

\begin{corrige}[Exercice 1 — QCM : les notions clés]
\textbf{1. Réponse b).} Une norme est un document de référence, élaboré par un organisme de normalisation (au Cameroun : l'ANOR, Agence des Normes et de la Qualité), qui fixe les caractéristiques qu'un produit doit respecter (dimensions, composition, sécurité, étiquetage) pour être mis sur le marché. Les réponses a), c) et d) sont fausses : une norme n'est ni une exigence réservée aux étrangers, ni une taxe, ni un label industriel.

\textbf{2. Réponse b).} La qualité désigne l'ensemble des caractéristiques d'un produit (aspect, solidité, goût, sécurité, régularité) qui lui permettent de satisfaire les besoins et les attentes du consommateur. Le prix (a), la fabrication manuelle (c) et la rareté (d) ne définissent pas la qualité : un produit fait main peut être de mauvaise qualité, un produit courant de très bonne qualité.

\textbf{3. Réponses a), b) et d).} La scie électrique, la ponceuse et la machine à écraser (manioc, huile de palme) traduisent la modernisation des équipements de travail : elles utilisent une énergie mécanique ou électrique qui augmente la rapidité et le volume de production. La herse tirée par des bœufs (c) est un outil agricole traditionnel à traction animale, sans rapport avec la modernisation artisanale.

\textbf{4. Réponse b).} Un circuit court est un circuit de vente où le producteur vend directement au consommateur (vente à l'atelier, au marché) ou avec un seul intermédiaire. Les réponses a), c) et d) sont fausses : la multiplication des intermédiaires définit au contraire un circuit long, et le circuit court n'est ni réservé à l'exportation ni dépendant d'Internet.
\end{corrige}

\begin{corrige}[Exercice 2 — Vrai ou faux ? Justifie]
\begin{enumerate}
\item \textbf{Faux.} La modernisation des équipements (scie électrique, ponceuse, machine à écraser) soulage l'artisan et accélère le travail, mais le geste, le dessin et le choix des formes restent issus du savoir-faire traditionnel : à Foumban, le sculpteur utilise la ponceuse pour finir, mais c'est toujours lui qui conçoit et sculpte la statue.

\item \textbf{Vrai.} Le contrôle de la qualité et le respect des normes de fabrication garantissent au client un produit fiable et conforme à ses attentes ; c'est une condition d'accès aux supermarchés et aux marchés extérieurs, comme le montre le poivre de Penja, exporté en Europe grâce à son label IGP et à des contrôles stricts.

\item \textbf{Faux.} Une filière regroupe l'ensemble des acteurs d'un produit, de la matière première à la vente : producteurs, artisans transformateurs, conditionneurs, transporteurs, commerçants, exportateurs. Les artisans fabricants n'en sont qu'un maillon.

\item \textbf{Vrai.} Le conditionnement (emballage propre, étiquetage, fermeture hermétique) protège le produit, le conserve plus longtemps, informe le consommateur et valorise sa présentation : il participe donc pleinement de l'amélioration de la qualité, comme pour le miel de l'Adamaoua vendu en pots étiquetés.
\end{enumerate}
\end{corrige}

\begin{corrige}[Exercice 3 — Définitions]
\begin{enumerate}
\item \textbf{Artisanat :} activité de production et de transformation exercée à petite échelle, dans un atelier, à l'aide d'outils manuels ou de petites machines, et reposant sur un savoir-faire professionnel souvent hérité. \emph{Exemple :} la sculpture sur bois et la fonte des bronzes à Foumban (région de l'Ouest).

\item \textbf{Filière :} ensemble des activités et des acteurs qui se succèdent pour un même produit, de l'approvisionnement en matière première jusqu'à la vente au consommateur (production, transformation, conditionnement, transport, commercialisation). \emph{Exemple :} la filière du poivre de Penja, des producteurs du Littoral aux exportateurs vers l'Europe.

\item \textbf{Norme :} document de référence élaboré par un organisme de normalisation (au Cameroun, l'ANOR) qui fixe les caractéristiques qu'un produit doit respecter (composition, dimensions, hygiène, étiquetage) pour être mis sur le marché. \emph{Exemple :} les normes d'hygiène et d'étiquetage exigées pour vendre du gari conditionné en supermarché.

\item \textbf{Circuit de vente :} cheminement suivi par un produit, avec ses intermédiaires éventuels, entre le producteur et le consommateur final ; on distingue les circuits courts (vente directe ou un seul intermédiaire) et les circuits longs (plusieurs intermédiaires). \emph{Exemple :} la vente directe de poteries au marché de Foumbot (circuit court) ou la vente de meubles de Foumban aux supermarchés de Douala (circuit long).
\end{enumerate}
\end{corrige}

\begin{corrige}[Exercice 4 — Calculs directs : lire une évolution]
\begin{enumerate}
\item Le taux de variation se calcule par :
\[
t = \frac{V_{2023} - V_{2018}}{V_{2018}} \times 100 = \frac{\num{2100} - \num{1200}}{\num{1200}} \times 100 = \frac{900}{\num{1200}} \times 100 = 75,0\,\%.
\]
Les ventes ont augmenté de \textbf{75,0 \%} entre 2018 et 2023.

\item La part des supermarchés urbains est :
\[
\frac{630}{\num{2100}} \times 100 = 30\,\%.
\]
Les supermarchés urbains représentent \textbf{30 \%} des ventes de 2023.

\item Le chiffre d'affaires est le produit des quantités vendues par le prix moyen :
\[
CA = \num{2100} \times \num{2300} = \num{4830000}\ \text{FCFA}.
\]
Le chiffre d'affaires de la coopérative en 2023 s'élève à \textbf{\num{4830000} FCFA} (4,83 millions de FCFA).
\end{enumerate}
\end{corrige}

\begin{corrige}[Exercice 5 — Classer les équipements]
\begin{enumerate}
\item Classement des équipements :
\begin{center}
\begin{tabularx}{\linewidth}{|X|X|}
\hline
\rowcolor{popblueL} Équipements traditionnels & Équipements modernisés \\
\hline
Machette & Scie électrique \\
\hline
Couteau à évider & Ponceuse \\
\hline
Mortier et pilon & Perceuse \\
\hline
Fer à repasser à charbon & Machine à écraser le manioc \\
\hline
 & Tour de potier électrique \\
\hline
 & Machine à coudre électrique \\
\hline
\end{tabularx}
\end{center}

\item Trois exemples d'effets (toute réponse cohérente est acceptée) :
\begin{itemize}
\item \textbf{Scie électrique} (menuiserie-ébénisterie) : découpe plus rapide et plus précise, effort physique réduit, volume de production fortement augmenté.
\item \textbf{Machine à écraser le manioc} (transformation alimentaire) : remplace le pilage au mortier, diminue fortement la pénibilité et permet de produire de grandes quantités de gari ou de bâton de manioc.
\item \textbf{Machine à coudre électrique} (couture, habillement) : couture régulière et rapide, finition de meilleure qualité, possibilité de répondre à des commandes importantes (uniformes, pagnes Ndop).
\end{itemize}
\end{enumerate}
\end{corrige}

\begin{corrige}[Exercice 6 — Identifier les produits nouveaux]
\begin{enumerate}
\item \textbf{Produits traditionnels :} les statues et les masques. \textbf{Produits nouveaux :} les meubles de décoration d'intérieur, les accessoires de mode en tissu Ndop, les sacs en perles et les objets utilitaires (lampes, cadres, porte-clés).

\item Ces produits nouveaux s'adressent à des clientèles nouvelles : les citadins camerounais qui décorent leurs intérieurs (Yaoundé, Douala), les touristes en quête de souvenirs modernes, la diaspora attachée aux produits du pays, et une clientèle de mode qui porte le Ndop en vêtements et accessoires contemporains.

\item L'apparition de ces produits est rendue possible par la modernisation des équipements : la scie électrique et la ponceuse permettent de fabriquer des meubles aux finitions soignées ; la machine à coudre électrique permet de confectionner sacs et vêtements en série ; la perceuse facilite le montage des lampes et cadres. La modernisation des outils élargit donc la gamme des produits offerts.
\end{enumerate}
\end{corrige}

\begin{corrige}[Exercice 7 — Lire une étiquette]
\begin{enumerate}
\item Quatre mentions qui rassurent le consommateur : la mention « miel pur » (nature du produit), la quantité « 500 g » (poids net), l'identité du producteur « GIC Miel du Vina, Banyo » (traçabilité), la date de production et la date limite de consommation (03/2024 et 03/2026), ainsi que le numéro de lot (lot n° 45).

\item Ces mentions traduisent le respect des normes de conditionnement et de présentation : les normes d'étiquetage imposent d'informer le consommateur sur la nature du produit, la quantité, le fabricant et les dates ; le numéro de lot permet de retracer le produit en cas de problème (traçabilité), ce qui est une exigence du contrôle de la qualité.

\item Par rapport au miel vendu en vrac dans des bouteilles non étiquetées, ce conditionnement inspire confiance, garantit l'hygiène et la conservation, et permet la vente dans les boutiques et supermarchés urbains, voire à l'exportation : le miel étiqueté accède à des marchés plus larges et se vend à meilleur prix.
\end{enumerate}
\end{corrige}

\begin{corrige}[Exercice 8 — Schématiser une filière]
\begin{enumerate}
\item[1. et 2.] Schéma de la filière du poivre de Penja (cases et opérations) :

\begin{popfigure}
\begin{center}
\begin{tikzpicture}[node distance=0.55cm,
case/.style={draw=popblue, fill=popblueL, rounded corners=2pt, text width=2.6cm, align=center, minimum height=0.9cm, font=\small},
op/.style={font=\scriptsize\itshape, text=popdark},
fleche/.style={-{Stealth[length=2.5mm]}, thick, poporange}]
\node[case, align=center] (prod){\textbf{Producteurs}\\ (Penja, Littoral)};
\node[case, right=of prod, align=center] (recolte){\textbf{Tri et séchage}\\ du poivre};
\node[case, right=of recolte, align=center] (controle){\textbf{Groupement}\\ contrôle qualité IGP};
\node[case, below=0.8cm of controle, align=center] (cond){\textbf{Conditionneurs}\\ emballage, étiquetage};
\node[case, left=of cond, align=center] (export){\textbf{Exportateurs}\\ (port de Douala)};
\node[case, left=of export, align=center] (distrib){\textbf{Distributeurs}\\ Cameroun et Europe};
\draw[fleche] (prod) -- (recolte);
\draw[fleche] (recolte) -- (controle);
\draw[fleche] (controle) -- (cond);
\draw[fleche] (cond) -- (export);
\draw[fleche] (export) -- (distrib);
\node[op, above=0.05cm of prod] {culture et récolte};
\node[op, below=0.05cm of distrib] {vente au consommateur};
\end{tikzpicture}
\end{center}
\legende{Schéma de la filière du poivre de Penja : de la production à la vente, chaque case correspond à une étape et à une opération (culture et récolte, tri et séchage, contrôle qualité, emballage, exportation, distribution).}
\end{popfigure}

\item[3.] L'organisation de la filière associe tous les acteurs autour d'un cahier des charges commun : le groupement contrôle la qualité à chaque étape (tri, séchage, conditionnement), ce qui garantit un produit régulier et conforme aux normes. L'obtention de l'IGP en 2013, première d'Afrique subsaharienne, protège le nom « poivre de Penja » et certifie son origine et sa qualité. Cette reconnaissance officielle rassure les acheteurs étrangers : le poivre de Penja est ainsi exporté et vendu à prix élevé sur les marchés européens, ce qui valorise le travail des producteurs.
\end{enumerate}
\end{corrige}

\begin{corrige}[Exercice 9 — Étude documentaire : deux ateliers, deux époques]
\begin{enumerate}
\item Dans l'atelier de Foumban des années 1970 (document 1), l'artisan travaille à la main, assis par terre, avec des outils traditionnels (couteau, herminette, racloir), et les produits sont exposés à même le sol. Dans l'atelier de Yaoundé en 2022 (document 2), le travail repose sur des machines électriques (scie circulaire, ponceuse, perceuse à colonne) et les meubles finis sont vernis, emballés sous film plastique et étiquetés : l'équipement est motorisé et le conditionnement soigné.

\item Taux de refus au contrôle qualité :
\[
2010 : \frac{120}{400} \times 100 = 30\,\% \qquad ; \qquad 2022 : \frac{72}{\num{1800}} \times 100 = 4\,\%.
\]
Taux de variation de la production entre 2010 et 2022 :
\[
\frac{\num{1800} - 400}{400} \times 100 = \frac{\num{1400}}{400} \times 100 = 350\,\%.
\]
Le taux de refus est tombé de 30 \% à 4 \%, tandis que la production a augmenté de 350 \%.

\item Trois transformations apparaissent : la \textbf{modernisation des équipements} (passage de l'outil manuel à la machine électrique, documents 1 et 2) ; l'\textbf{amélioration de la qualité} (taux de refus divisé par plus de 7, vernis, emballage, étiquetage, documents 2 et 3) ; l'\textbf{élargissement des débouchés} (la part vendue hors de la région passe de 10 \% à 45 \%, document 3).

\item Ces transformations élargissent les marchés : la mécanisation augmente les volumes produits, ce qui permet de répondre à de grosses commandes ; la qualité régulière et le conditionnement soigné (emballage, étiquette) donnent accès aux supermarchés et aux clients des autres régions ; la confiance ainsi créée ouvre même la voie à l'exportation. L'artisan ne vend plus seulement devant son atelier, mais sur l'ensemble du territoire national et au-delà.
\end{enumerate}

\textbf{Barème indicatif (sur 20) :} description comparée des ateliers : 4 pts ; calculs exacts et rédigés (taux de refus, taux de variation) : 5 pts ; dégagement des trois transformations : 6 pts ; explication de l'élargissement des marchés : 4 pts ; présentation : 1 pt.

\textbf{Points de vigilance :} ne pas confondre taux de refus (refusés/production) et part vendue hors région ; citer les documents dans la réponse (« d'après le document 2... ») ; un taux de variation se calcule toujours par rapport à la valeur de départ.
\end{corrige}

\begin{corrige}[Exercice 10 — Construction graphique : les circuits de vente]
\begin{enumerate}
\item Parts et angles (angle $= \dfrac{\text{ventes}}{240} \times 360$) :
\begin{center}
\begin{tabularx}{\linewidth}{|X|c|c|}
\hline
\rowcolor{popblueL} Circuit de vente & Part (\%) & Angle (arrondi) \\
\hline
Vente directe au marché local & $90/240 = 37,5\,\%$ & $135^{\circ}$ \\
\hline
Détaillants et boutiques de quartier & $60/240 = 25\,\%$ & $90^{\circ}$ \\
\hline
Supermarchés de Yaoundé et Douala & $75/240 = 31,25\,\%$ & $112^{\circ}$ (arrondi) \\
\hline
Exportation vers la diaspora & $15/240 = 6,25\,\%$ & $23^{\circ}$ (arrondi) \\
\hline
\textbf{Total} & \textbf{100 \%} & \textbf{360$^{\circ}$} \\
\hline
\end{tabularx}
\end{center}
(Vérification : $135 + 90 + 112 + 23 = 360^{\circ}$ ; $37,5 + 25 + 31,25 + 6,25 = 100\,\%$.)

\item Diagramme circulaire :

\begin{popfigure}
\begin{center}
\begin{tikzpicture}[scale=1.1]
% Marché local : 135° (de 90° à -45°), boutiques : 90°, supermarchés : 112.5°, export : 22.5°
\fill[popblueL] (0,0) -- (90:2) arc (90:-45:2) -- cycle;
\fill[popgreenL] (0,0) -- (-45:2) arc (-45:-135:2) -- cycle;
\fill[poporangeL] (0,0) -- (-135:2) arc (-135:-247.5:2) -- cycle;
\fill[poppinkL] (0,0) -- (-247.5:2) arc (-247.5:-270:2) -- cycle;
\draw[popdark, thick] (0,0) circle (2);
\draw[popdark] (0,0) -- (90:2); \draw[popdark] (0,0) -- (-45:2);
\draw[popdark] (0,0) -- (-135:2); \draw[popdark] (0,0) -- (-247.5:2);
\node[font=\scriptsize] at (22.5:1.3) {37,5 \%};
\node[font=\scriptsize] at (-90:1.3) {25 \%};
\node[font=\scriptsize] at (-191:1.35) {31,25 \%};
\node[font=\scriptsize] at (-258:1.35) {6,25 \%};
% Légende
\node[anchor=west, font=\scriptsize] at (2.5,1.2) {\textcolor{popblue}{\rule{0.35cm}{0.2cm}} Vente directe au marché local};
\node[anchor=west, font=\scriptsize] at (2.5,0.7) {\textcolor{popgreen}{\rule{0.35cm}{0.2cm}} Détaillants et boutiques};
\node[anchor=west, font=\scriptsize] at (2.5,0.2) {\textcolor{poporange}{\rule{0.35cm}{0.2cm}} Supermarchés (Yaoundé, Douala)};
\node[anchor=west, font=\scriptsize] at (2.5,-0.3) {\textcolor{poppink}{\rule{0.35cm}{0.2cm}} Exportation (diaspora, Europe)};
\node[font=\small\bfseries, text=popdark] at (0,-2.6) {Répartition des ventes de gari conditionné par circuit (2023)};
\end{tikzpicture}
\end{center}
\legende{Diagramme circulaire des ventes de la filière gari en 2023 (total : 240 tonnes). Source : données de l'énoncé ; calculs de l'élève.}
\end{popfigure}

\item Croquis de la filière :

\begin{popfigure}
\begin{center}
\begin{tikzpicture}[node distance=0.5cm,
case/.style={draw=popblue, fill=popblueL, rounded corners=2pt, text width=2.7cm, align=center, minimum height=0.85cm, font=\small},
fleche/.style={-{Stealth[length=2.5mm]}, thick, poporange}]
\node[case] (prod) {\textbf{Producteurs de manioc} (villages du Centre)};
\node[case, right=1cm of prod] (transfo) {\textbf{Artisans transformateurs} (écrasement, gari conditionné)};
\node[case, right=1cm of transfo, text width=3.1cm] (circuits) {\textbf{Circuits de vente}\\ marché local ; boutiques ; supermarchés ; exportation};
\node[case, right=1cm of circuits] (conso) {\textbf{Consommateurs} (Cameroun et diaspora)};
\draw[fleche] (prod) -- (transfo);
\draw[fleche] (transfo) -- (circuits);
\draw[fleche] (circuits) -- (conso);
\end{tikzpicture}
\end{center}
\legende{Croquis de la filière gari conditionné : 1. production du manioc ; 2. transformation artisanale (machine à écraser, conditionnement) ; 3. quatre circuits de vente ; 4. consommateurs finaux.}
\end{popfigure}

\item Le circuit dominant est la \textbf{vente directe au marché local} (37,5 \%) : le circuit court reste le principal débouché. Mais la place des supermarchés (31,25 \%, près d'un tiers des ventes) et de l'exportation (6,25 \%) révèle la modernisation de la filière : seul un produit conditionné, étiqueté et de qualité contrôlée peut entrer en supermarché et traverser les frontières. La filière combine donc encore circuits traditionnels et circuits modernes.
\end{enumerate}

\textbf{Barème indicatif (sur 20) :} calculs des parts et des angles (tableau vérifié à 360$^{\circ}$) : 5 pts ; diagramme circulaire soigné avec titre et légende : 5 pts ; croquis de la filière complet et légendé : 5 pts ; interprétation argumentée : 4 pts ; présentation : 1 pt.

\textbf{Points de vigilance :} l'angle se calcule sur le total (240 t), pas sur chaque ligne ; vérifier que la somme des angles fait 360$^{\circ}$ ; titre et légende sont obligatoires sur toute construction graphique au Bac.
\end{corrige}

\begin{corrige}[Exercice 11 — Dissertation guidée : artisanat et intégration nationale]
\textbf{Introduction (éléments attendus).} Amorce : à Douala, on achète aujourd'hui dans un supermarché un meuble fabriqué à Foumban ou un pagne en tissu Ndop tissé dans le Nord-Ouest. Définitions : l'artisanat est une activité de production à petite échelle reposant sur un savoir-faire professionnel ; sa modernisation désigne le renouvellement des équipements, des produits, de la qualité et de l'organisation des ventes. Problématique : en quoi cette modernisation tisse-t-elle des liens économiques entre les régions du Cameroun ? Annonce du plan en trois parties.

\textbf{I. Un artisanat hérité, ancré dans les régions du Cameroun.}
\begin{enumerate}
\item Le Cameroun possède de grands pôles artisanaux spécialisés : sculpture sur bois et bronzes de Foumban (Ouest), poterie de l'Ouest et du Nord, tissage du Ndop (Nord-Ouest), vannerie, travail du cuir et transformation alimentaire (gari, huile de palme, miel de l'Adamaoua). Chaque région met ainsi en valeur ses ressources et ses traditions.
\item Mais l'artisanat traditionnel souffre de limites : outillage manuel et faible productivité, qualité inégale d'une pièce à l'autre, absence de conditionnement, marchés presque uniquement locaux (vente devant l'atelier ou au marché). Ces limites confinent l'artisan à un rayon de vente étroit et à de faibles revenus.
\end{enumerate}

\textbf{II. Une modernisation qui transforme la production.}
\begin{enumerate}
\item Les équipements modernes — machines à écraser le manioc ou le palmier, scies électriques, ponceuses, perceuses, machines à coudre — augmentent la rapidité, la précision et les volumes, tout en réduisant la pénibilité. Apparaissent des produits nouveaux : meubles de décoration d'intérieur, accessoires de mode en Ndop, produits alimentaires transformés et conditionnés.
\item Le contrôle de la qualité, le respect des normes (ANOR) et les nouvelles méthodes de conditionnement (emballage, étiquetage, traçabilité) ouvrent les marchés nationaux (supermarchés de Yaoundé et Douala) et internationaux : le poivre de Penja, première IGP d'Afrique subsaharienne (2013), est ainsi exporté vers l'Europe.
\end{enumerate}

\textbf{III. Une contribution à l'intégration économique nationale.}
\begin{enumerate}
\item L'organisation des filières (producteurs, artisans, conditionneurs, transporteurs, commerçants) et des circuits de vente relie les régions entre elles : les meubles de Foumban partent vers Douala et Yaoundé, le miel de l'Adamaoua descend vers le Sud, le gari du Centre alimente les villes ; ces flux empruntent les axes routiers et le port de Douala.
\item Les effets sont sensibles : hausse des revenus des artisans, création d'emplois dans les ateliers et les services, mise en valeur des ressources locales de chaque région, devises rapportées par les exportations et rayonnement de l'image du Cameroun à l'extérieur. L'artisanat modernisé devient un maillon de l'intégration économique nationale.
\end{enumerate}

\textbf{Conclusion (éléments attendus).} Bilan : en modernisant ses équipements, sa qualité et son organisation, l'artisanat camerounais dépasse le cadre local et relie les régions au marché national et au monde. Ouverture : des défis demeurent — financement des équipements, formation des artisans, accès à l'électricité en zone rurale.

\textbf{Barème indicatif (sur 20) :} introduction (amorce, définitions, problématique, plan) : 3 pts ; partie I (pôles + limites) : 4 pts ; partie II (équipements, produits nouveaux, normes, exemple Penja) : 5 pts ; partie III (flux interrégionaux + effets) : 5 pts ; conclusion : 2 pts ; langue et présentation : 1 pt.

\textbf{Points de vigilance :} respecter le plan imposé et annoncer les parties ; citer au moins trois exemples camerounais précis (lieu + produit) ; employer correctement les mots-clés du programme (\cle{qualité}, \cle{norme}, \cle{filière}, circuit de vente) ; éviter les listes sans phrase d'explication : chaque exemple doit être relié à l'intégration nationale.
\end{corrige}

\newpage

\coursec{Fiche bilan}

\begin{popfigure}
\begin{center}\tcbox[colback=popblue,colframe=popblue,coltext=white,arc=9pt,boxsep=4pt,left=6pt,right=6pt]{\bfseries\small L'évolution de l'artisanat camerounais}\end{center}
\vspace{-2pt}
\begin{tcbraster}[raster columns=3,raster equal height=rows,raster column skip=6pt,raster row skip=6pt,enhanced,blankest]
\begin{tcolorbox}[enhanced,colback=white,colframe=poporange,boxrule=0.9pt,arc=3pt,title={Héritage traditionnel},fonttitle=\bfseries\scriptsize,coltitle=white,colbacktitle=poporange,left=4pt,right=4pt,top=2pt,bottom=2pt,boxsep=1pt,toptitle=1.5pt,bottomtitle=1.5pt,halign=left]
\begin{itemize}[nosep,leftmargin=*,itemsep=1pt,font=\scriptsize]
\item Sculpture, poterie, vannerie, tissage
\item Savoir-faire locaux (Foumban, Bafoussam\ldots)
\end{itemize}
\end{tcolorbox}
\begin{tcolorbox}[enhanced,colback=white,colframe=popgreen,boxrule=0.9pt,arc=3pt,title={Modernisation des équipements},fonttitle=\bfseries\scriptsize,coltitle=white,colbacktitle=popgreen,left=4pt,right=4pt,top=2pt,bottom=2pt,boxsep=1pt,toptitle=1.5pt,bottomtitle=1.5pt,halign=left]
\begin{itemize}[nosep,leftmargin=*,itemsep=1pt,font=\scriptsize]
\item Scies électriques, ponceuses, perceuses
\item Machines à écraser (agroalimentaire)
\end{itemize}
\end{tcolorbox}
\begin{tcolorbox}[enhanced,colback=white,colframe=poppurple,boxrule=0.9pt,arc=3pt,title={Produits nouveaux},fonttitle=\bfseries\scriptsize,coltitle=white,colbacktitle=poppurple,left=4pt,right=4pt,top=2pt,bottom=2pt,boxsep=1pt,toptitle=1.5pt,bottomtitle=1.5pt,halign=left]
\begin{itemize}[nosep,leftmargin=*,itemsep=1pt,font=\scriptsize]
\item Décoration, habillement (pagne, design)
\item Transformation alimentaire
\end{itemize}
\end{tcolorbox}
\begin{tcolorbox}[enhanced,colback=white,colframe=poppink,boxrule=0.9pt,arc=3pt,title={Qualité et normes},fonttitle=\bfseries\scriptsize,coltitle=white,colbacktitle=poppink,left=4pt,right=4pt,top=2pt,bottom=2pt,boxsep=1pt,toptitle=1.5pt,bottomtitle=1.5pt,halign=left]
\begin{itemize}[nosep,leftmargin=*,itemsep=1pt,font=\scriptsize]
\item Conditionnement, présentation
\item Contrôle qualité, normes (ANOR)
\end{itemize}
\end{tcolorbox}
\begin{tcolorbox}[enhanced,colback=white,colframe=popgold,boxrule=0.9pt,arc=3pt,title={Filières et circuits de vente},fonttitle=\bfseries\scriptsize,coltitle=white,colbacktitle=popgold,left=4pt,right=4pt,top=2pt,bottom=2pt,boxsep=1pt,toptitle=1.5pt,bottomtitle=1.5pt,halign=left]
\begin{itemize}[nosep,leftmargin=*,itemsep=1pt,font=\scriptsize]
\item GIC, coopératives, centres artisanaux
\item Marchés, foires (SIAO), e-commerce
\end{itemize}
\end{tcolorbox}
\begin{tcolorbox}[enhanced,colback=white,colframe=popmuted,boxrule=0.9pt,arc=3pt,title={Enjeux et limites},fonttitle=\bfseries\scriptsize,coltitle=white,colbacktitle=popmuted,left=4pt,right=4pt,top=2pt,bottom=2pt,boxsep=1pt,toptitle=1.5pt,bottomtitle=1.5pt,halign=left]
\begin{itemize}[nosep,leftmargin=*,itemsep=1pt,font=\scriptsize]
\item Financement, formation, énergie
\item Concurrence des importations
\end{itemize}
\end{tcolorbox}
\end{tcbraster}

\legende{Carte mentale de la leçon : le nœud central (bleu) résume le thème ; les six branches colorées reprennent les grandes parties du cours (héritage, équipements, produits, qualité-normes, filières, enjeux).}
\end{popfigure}

\begin{bilan}
\textbf{Définitions clés à connaître par c\oe ur}
\begin{itemize}
  \item \cle{Artisanat} : activité de production et de transformation réalisée à petite échelle, souvent manuellement ou avec des outils simples, par un artisan seul ou avec peu d'ouvriers (menuiserie, forge, poterie, couture, transformation alimentaire).
  \item \cle{Qualité} : ensemble des caractéristiques d'un produit (aspect, solidité, hygiène, durabilité) qui lui permettent de répondre aux attentes du consommateur et de se vendre durablement.
  \item \cle{Norme} : règle technique de référence (dimensions, composition, hygiène, sécurité, étiquetage) définie par un organisme (au Cameroun : l'\textbf{ANOR}, Agence des Normes et de la Qualité) que le produit doit respecter pour être commercialisé, notamment à l'exportation.
  \item \cle{Filière} : enchaînement organisé des étapes allant de l'approvisionnement en matière première jusqu'à la vente du produit fini (ex. filière bois : forêt $\rightarrow$ scierie $\rightarrow$ menuisier $\rightarrow$ vendeur).
  \item \cle{Circuit de vente} : chemin parcouru par le produit de l'artisan au consommateur (vente directe à l'atelier, marché, coopérative, foire-exposition, boutique, internet).
\end{itemize}

\textbf{Les quatre mutations de l'artisanat camerounais}
\begin{center}
\begin{tabularx}{\linewidth}{|l|X|X|}\hline
\rowcolor{popblueL} \textbf{Mutation} & \textbf{Contenu} & \textbf{Exemples camerounais} \\ \hline
Équipements modernes & Machines à écraser, scies électriques, ponceuses, perceuses, machines à coudre industrielles & Menuiseries de Yaoundé et Douala ; moulins à manioc et presses à huile de palme en zone rurale \\ \hline
Produits nouveaux & Décoration intérieure, habillement moderne (pagne revisité, accessoires), produits alimentaires transformés et emballés & Ndop et tissus teints de l'Ouest ; jus, épices et farines conditionnés \\ \hline
Qualité et normes & Nouveaux conditionnements, présentation soignée, contrôle qualité, respect des normes de fabrication & Étiquetage, emballages hygiéniques, certification ANOR pour l'export \\ \hline
Filières et circuits & Regroupement en GIC et coopératives, centres artisanaux, foires, vente en ligne & Village artisanal de Foumban ; foires artisanales ; vente via réseaux sociaux \\ \hline
\end{tabularx}
\end{center}

\textbf{Méthodes express}
\begin{itemize}
  \item \emph{Commenter un document} (photo, tableau, carte) : 1) identifier et situer ; 2) décrire les faits visibles ; 3) expliquer en mobilisant les notions du cours (modernisation, qualité, filière) ; 4) conclure sur l'enjeu (intégration nationale, emploi, revenus).
  \item \emph{Construire un croquis} : titre, légende numérotée, flèche du nord, échelle indicative ; figurer les flux par des flèches d'épaisseur proportionnelle.
  \item \emph{Analyser une filière} : toujours suivre la chaîne matière première $\rightarrow$ transformation $\rightarrow$ conditionnement $\rightarrow$ commercialisation, en nommant les acteurs à chaque étape.
\end{itemize}

\textbf{Repères chiffrés (ordres de grandeur prudents)}
\begin{itemize}
  \item L'artisanat et le secteur informel occupent une part très importante de l'emploi non agricole au Cameroun (l'informel représente l'essentiel des emplois urbains, ordre de grandeur : plusieurs millions d'actifs).
  \item L'artisanat est encadré par le ministère en charge de l'Artisanat et des Petites et Moyennes Entreprises ; l'ANOR veille aux normes et à la qualité.
\end{itemize}
\end{bilan}

\begin{aretenir}[Checklist avant le Bac]
\begin{itemize}
  \item $\square$ Je sais définir \cle{artisanat}, \cle{qualité}, \cle{norme}, \cle{filière} et \cle{circuit de vente} en une phrase chacun.
  \item $\square$ Je cite au moins quatre équipements modernes (machine à écraser, scie électrique, ponceuse, perceuse) et leur effet sur la production.
  \item $\square$ Je donne des exemples de produits nouveaux dans la décoration, l'habillement et l'agroalimentaire.
  \item $\square$ J'explique comment le conditionnement et la présentation améliorent la valeur marchande d'un produit.
  \item $\square$ Je connais le rôle de l'ANOR dans le contrôle de la qualité et le respect des normes.
  \item $\square$ Je décris les étapes d'une filière artisanale complète (ex. bois, cacao, textile).
  \item $\square$ Je cite les circuits de vente : atelier, marché, coopérative/GIC, foire, boutique, internet.
  \item $\square$ Je localise sur une carte schématique les grands pôles artisanaux (Foumban, Bafoussam, Douala, Yaoundé, Maroua).
  \item $\square$ Je sais construire et légender un croquis ou un diagramme de filière avec titre, légende et flèches de flux.
  \item $\square$ Je sais relier l'évolution de l'artisanat à la famille de situations : l'\cle{intégration nationale} (emplois, revenus, identité culturelle, désenclavement).
  \item $\square$ Je mentionne les limites : accès au financement, à l'électricité, à la formation, concurrence des produits importés.
  \item $\square$ Je rédige mes réponses en français soigné, avec des exemples camerounais précis et chiffrés avec prudence.
\end{itemize}
\end{aretenir}

\findecours

\end{document}
